
% Билеты по коллоквиуму "Алгебра. Весенний семестр. Часть II"
% Полный конспект, разложенный по билетам.
% Материал сохранён в полном объёме и аккуратно разложен по билетам.

\documentclass[12pt,a4paper]{article}
\usepackage{iftex}
\ifPDFTeX
    \usepackage[T2A]{fontenc}
    \usepackage[utf8]{inputenc}
    \usepackage[russian]{babel}
\else
    \usepackage{fontspec}
    \usepackage{polyglossia}
    \setmainlanguage{russian}
    \setmainfont{DejaVu Serif}
    \setsansfont{DejaVu Sans}
    \setmonofont{DejaVu Sans Mono}
\fi
\usepackage{amsmath,amssymb,amsthm,mathtools}
\usepackage{xcolor}
\usepackage{enumitem}
\usepackage{hyperref}
\usepackage{geometry}
\usepackage[most]{tcolorbox}
\usepackage{tikz}
\usetikzlibrary{arrows.meta,calc,positioning,angles,quotes}
\geometry{margin=2.0cm}
\allowdisplaybreaks

\hypersetup{colorlinks=true, linkcolor=blue, urlcolor=blue, citecolor=blue}

\theoremstyle{definition}
\newtheorem{definition}{Определение}[section]
\newtheorem{example}[definition]{Пример}
\theoremstyle{plain}
\newtheorem{theorem}[definition]{Теорема}
\newtheorem{lemma}[definition]{Лемма}
\newtheorem{proposition}[definition]{Утверждение}
\newtheorem{corollary}[definition]{Следствие}
\theoremstyle{remark}
\newtheorem{remark}[definition]{Замечание}

\newtcolorbox{planbox}{colback=blue!3!white,colframe=blue!40!black,title=План ответа,breakable}
\newenvironment{planadd}{\begin{planbox}}{\end{planbox}}
\newtcolorbox{clarifybox}{colback=green!2!white,colframe=green!35!black,title=Уточнение,breakable}
\newenvironment{knowledgefix}{\begin{clarifybox}}{\end{clarifybox}}

% Обозначения
\newcommand{\F}{\mathbb{F}}
\newcommand{\R}{\mathbb{R}}
\let\C\relax
\newcommand{\C}{\mathbb{C}}
\newcommand{\N}{\mathbb{N}}
\newcommand{\E}{\mathcal{E}}
\newcommand{\A}{\mathcal{A}}
\newcommand{\BL}{\operatorname{BL}}
\newcommand{\Mat}{\operatorname{Mat}}
\newcommand{\End}{\operatorname{End}}
\newcommand{\Hom}{\operatorname{Hom}}
\newcommand{\Ker}{\operatorname{Ker}}
\newcommand{\Imop}{\operatorname{Im}}
\newcommand{\rk}{\operatorname{rk}}
\providecommand{\tr}{\operatorname{tr}}
\newcommand{\diag}{\operatorname{diag}}
\newcommand{\Span}{\operatorname{span}}
\newcommand{\Id}{\operatorname{Id}}
\newcommand{\Gr}{\operatorname{Gram}}
\newcommand{\Vol}{\operatorname{Vol}}
\newcommand{\dist}{\operatorname{dist}}
\newcommand{\pr}{\operatorname{pr}}
\newcommand{\ort}{\operatorname{ort}}
\newcommand{\Sym}{\operatorname{Sym}}
\newcommand{\Alt}{\operatorname{Alt}}
\providecommand{\sgn}{\operatorname{sgn}}
\newcommand{\Tens}[2]{\Pi^{#2}_{#1}(V)}


% Геометрические иллюстрации
\newcommand{\VisualProjection}{%
\begin{center}
\begin{tikzpicture}[scale=1.0,>=Latex]
    \draw[thick] (-3,0) -- (3,0) node[right] {$U$};
    \fill (0,0) circle (1.4pt) node[below right] {$0$};
    \coordinate (x) at (2,1.65);
    \coordinate (p) at (2,0);
    \draw[->,thick] (0,0) -- (x) node[above right] {$x$};
    \draw[->,thick] (0,0) -- (p) node[below] {$\operatorname{pr}_U x$};
    \draw[dashed,thick] (p) -- (x) node[midway,right] {$x-\operatorname{pr}_U x\in U^\perp$};
    \draw (1.83,0) -- (1.83,0.17) -- (2,0.17);
    \node[align=center] at (0,-0.75) {Ортогональное разложение: $x=\operatorname{pr}_U x+\operatorname{ort}_U x$};
\end{tikzpicture}
\end{center}}

\newcommand{\VisualGramSchmidt}{%
\begin{center}
\begin{tikzpicture}[scale=1.0,>=Latex]
    \coordinate (O) at (0,0);
    \coordinate (eone) at (3,0);
    \coordinate (x) at (2.1,1.8);
    \coordinate (proj) at (2.1,0);
    \draw[->,thick] (O) -- (eone) node[right] {$e_1$};
    \draw[->,thick] (O) -- (x) node[above] {$x_2$};
    \draw[->,thick] (O) -- (proj) node[below] {$\operatorname{pr}_{\langle e_1\rangle}x_2$};
    \draw[dashed,thick] (proj) -- (x) node[midway,right] {$u_2$};
    \draw (1.92,0) -- (1.92,0.18) -- (2.1,0.18);
    \node[align=center] at (1.7,-0.8) {Шаг Грама--Шмидта: $u_2=x_2-\operatorname{pr}_{\langle e_1\rangle}x_2$};
\end{tikzpicture}
\end{center}}

\newcommand{\VisualQuadrics}{%
\begin{center}
\begin{tikzpicture}[scale=0.88]
    \begin{scope}[xshift=-5.1cm]
        \draw[->] (-1.9,0)--(1.9,0) node[right] {$x$};
        \draw[->] (0,-1.4)--(0,1.4) node[above] {$y$};
        \draw[thick,smooth,domain=0:360,variable=\t] plot ({1.45*cos(\t)},{0.85*sin(\t)});
        \node at (0,-1.85) {эллипс};
    \end{scope}
    \begin{scope}
        \draw[->] (-1.9,0)--(1.9,0) node[right] {$x$};
        \draw[->] (0,-1.4)--(0,1.4) node[above] {$y$};
        \draw[thick,smooth,domain=-1.2:1.2] plot ({cosh(\x)},{sinh(\x)});
        \draw[thick,smooth,domain=-1.2:1.2] plot ({-cosh(\x)},{sinh(\x)});
        \node at (0,-1.85) {гипербола};
    \end{scope}
    \begin{scope}[xshift=5.1cm]
        \draw[->] (-1.9,0)--(1.9,0) node[right] {$x$};
        \draw[->] (0,-1.4)--(0,1.4) node[above] {$y$};
        \draw[thick,smooth,domain=-1.2:1.2] plot ({\x},{0.8*\x*\x-0.55});
        \node at (0,-1.85) {парабола};
    \end{scope}
\end{tikzpicture}
\end{center}}

\newcommand{\VisualContraction}{%
\begin{center}
\begin{tikzpicture}[>=Latex,node distance=1.3cm]
    \node[draw,rounded corners,align=center,minimum width=2.8cm,minimum height=1cm] (T) {$T^{i\,a}_{j\,b}$\\тензор типа $(2,2)$};
    \node[draw,rounded corners,align=center,minimum width=3.1cm,minimum height=1cm,right=2.1cm of T] (C) {$S^{a}_{b}=T^{i\,a}_{i\,b}$\\свёртка по $i$};
    \draw[->,thick] (T) -- node[above] {один верхний + один нижний} (C);
    \node[align=center,below=0.75cm of T, xshift=2.55cm] {Свёртка уменьшает валентность на два и не зависит от выбора базиса.};
\end{tikzpicture}
\end{center}}

\newcommand{\VisualSymAlt}{%
\begin{center}
\begin{tikzpicture}[>=Latex,node distance=1.1cm]
    \node[draw,rounded corners,align=center,minimum width=2.4cm] (T) {$T\in T^p_0(V)$};
    \node[draw,rounded corners,align=center,minimum width=2.8cm,below left=1.2cm and 1.3cm of T] (S) {$\Sym T=\frac1{p!}\sum_{\sigma\in S_p}F_\sigma T$};
    \node[draw,rounded corners,align=center,minimum width=3.1cm,below right=1.2cm and 1.3cm of T] (A) {$\Alt T=\frac1{p!}\sum_{\sigma\in S_p}\sgn(\sigma)F_\sigma T$};
    \draw[->,thick] (T) -- (S);
    \draw[->,thick] (T) -- (A);
    \node[below=0.25cm of S,align=center] {симметрическая часть};
    \node[below=0.25cm of A,align=center] {кососимметрическая часть};
\end{tikzpicture}
\end{center}}

\begin{document}
\begin{center}
{\LARGE Полный конспект по линейной алгебре}\\[3pt]
{\large Билинейные формы, евклидовы и эрмитовы пространства, аффинная геометрия, тензоры}\\[4pt]
{\normalsize Разметка по билетам 1--27}
\end{center}

\begin{remark}[О конспекте]
Материал разложен по официальным билетам и сохранён с теоретическими подробностями: определения, формулы, доказательные идеи, примеры и вычислительные алгоритмы оставлены внутри соответствующих вопросов. Поясняющие блоки предназначены для устного ответа и не заменяют основной теоретический текст.
\end{remark}

\begin{remark}[Как пользоваться]
Каждый билет устроен как законченный фрагмент конспекта: сначала указано содержание вопроса, затем идут нужные определения, утверждения, доказательства и примеры. Повторы между билетами сохранены там, где они помогают отвечать на каждый вопрос независимо.
\end{remark}

\tableofcontents
\clearpage
\section{Билет 1. Билинейные формы и их матрицы. Преобразование матрицы билинейной формы при замене базиса. Ранг и ядра билинейной формы.}

\paragraph{Содержание билета.} Билинейные формы и их матрицы. Преобразование матрицы билинейной формы при замене базиса. Ранг и ядра билинейной формы.

\begin{planadd}

порядок ответа — дать определение билинейности, записать матрицу $B=(\beta(e_i,e_j))$, вывести формулу $\beta(x,y)=X^TBY$, затем показать замену базиса $B' = C^TBC$, после этого обсудить ранг и левое/правое ядра.

\end{planadd}

\subsubsection{Билинейные формы}


\begin{definition}[Билинейная форма]
Пусть $V$ --- векторное пространство над полем $\F$.
Отображение
\[
    \beta: V \times V \to \F
\]
называется \emph{билинейной формой}, если оно линейно по каждому аргументу:
для любых $x,x_1,x_2,y,y_1,y_2\in V$ и $\alpha\in\F$ выполнено
\begin{enumerate}[label=\arabic*)]
    \item $\beta(x_1+x_2,y)=\beta(x_1,y)+\beta(x_2,y)$;
    \item $\beta(\alpha x,y)=\alpha\beta(x,y)$;
    \item $\beta(x,y_1+y_2)=\beta(x,y_1)+\beta(x,y_2)$;
    \item $\beta(x,\alpha y)=\alpha\beta(x,y)$.
\end{enumerate}
\end{definition}

\begin{example}
Нулевая форма:
\[
    \beta(x,y)=0.
\]
\end{example}

\begin{example}
Пусть $V=\F^n$, $x=(x_1,\ldots,x_n)^T$, $y=(y_1,\ldots,y_n)^T$.
Тогда пример билинейной формы:
\[
    \beta(x,y)=x_1y_1+\cdots+x_ny_n=\sum_{i=1}^n x_i y_i.
\]
\end{example}

\begin{example}
Если $f,g\in V^*$, то
\[
    \beta(x,y)=f(x)g(y)
\]
является билинейной формой.
\end{example}

\begin{example}
Для матриц $A,B\in M_n(\F)$ в конспекте приведён пример вида
\[
    \beta(A,B)=\tr(AB).
\]
\begin{knowledgefix}
На пространстве матриц стандартный проверенный пример билинейной формы записывается как
\(\langle A,B\rangle=\tr(A^{T}B)\) над \(\R\) и как \(\tr(A^{*}B)\) над \(\C\). Запись \(\tr(AB)\) тоже билинейна, но не является стандартным положительным скалярным произведением на всех матрицах.
\end{knowledgefix}
\end{example}

\begin{example}
На странице приведён пример с определителем матрицы $2\times 2$.
\begin{knowledgefix}
Определитель одной матрицы \(\det A\) не является билинейной формой от пары матриц. Корректный пример с определителем: \((x,y)\mapsto \det(x,y)\) на \(\R^2\) — кососимметрическая билинейная форма от двух векторов; в размерности \(n\) \(\det(v_1,\ldots,v_n)\) — полилинейная антисимметрическая форма.
\end{knowledgefix}
\end{example}

\begin{example}
Если $f,g\in C[a,b]$, то
\[
    \beta(f,g)=\int_a^b f(x)g(x)\,dx
\]
является билинейной формой.
\end{example}


\begin{proposition}[Операции с билинейными формами]
Если $\beta,\beta_1,\beta_2$ --- билинейные формы на $V$, а $\alpha\in\F$, то формы
\[
    \alpha\beta,
    \qquad
    \beta_1+\beta_2
\]
задаются правилами
\[
    (\alpha\beta)(x,y)=\alpha\beta(x,y),
\]
\[
    (\beta_1+\beta_2)(x,y)=\beta_1(x,y)+\beta_2(x,y),
\]
и снова являются билинейными формами.
\end{proposition}

\begin{remark}
Билинейные формы на $V$ образуют векторное пространство.
Обозначение из конспекта:
\[
    \BL(V).
\]
\end{remark}

\begin{definition}[Матрица билинейной формы]
Пусть $\dim V=n$ и $e_1,\ldots,e_n$ --- базис пространства $V$.
Для билинейной формы $\beta$ положим
\[
    b_{ij}=\beta(e_i,e_j).
\]
Матрица
\[
    B=(b_{ij})_{i,j=1}^n
\]
называется матрицей билинейной формы $\beta$ в базисе $e_1,\ldots,e_n$.
\end{definition}

Если
\[
    x=e_1x_1+\cdots+e_nx_n,\qquad
    y=e_1y_1+\cdots+e_ny_n,
\]
то по билинейности
\[
\begin{aligned}
    \beta(x,y)
    &=\beta(e_1x_1+\cdots+e_nx_n,
             e_1y_1+\cdots+e_ny_n) \\
    &=\sum_{i=1}^n\sum_{j=1}^n x_i\beta(e_i,e_j)y_j.
\end{aligned}
\]
Если $X=(x_1,\ldots,x_n)^T$, $Y=(y_1,\ldots,y_n)^T$, то
\[
    \beta(x,y)=X^TBY.
\]

\begin{corollary}
Соответствие
\[
    \beta \longmapsto B=(\beta(e_i,e_j))
\]
задаёт биекцию между пространством билинейных форм на $V$ и пространством матриц $M_n(\F)$.
Поэтому
\[
    \dim \BL(V)=n^2.
\]
\end{corollary}


\begin{proposition}[Изменение матрицы билинейной формы при замене базиса]
Пусть $e=(e_1,\ldots,e_n)$ и $\widetilde e=(\widetilde e_1,\ldots,\widetilde e_n)$ --- два базиса пространства $V$.
Пусть матрица перехода $C$ задана соотношением
\[
    \widetilde e = eC.
\]
Если $X,Y$ --- координатные столбцы в старом базисе, а $\widetilde X,\widetilde Y$ --- координатные столбцы в новом базисе, то
\[
    X=C\widetilde X,
    \qquad
    Y=C\widetilde Y.
\]
Тогда матрицы билинейной формы в этих базисах связаны формулой
\[
    \widetilde B=C^TBC.
\]
\end{proposition}

\begin{proof}
В старом базисе
\[
    \beta(x,y)=X^TBY.
\]
После замены координат $X=C\widetilde X$, $Y=C\widetilde Y$ получаем
\[
\begin{aligned}
    \beta(x,y)
    &=(C\widetilde X)^T B(C\widetilde Y) \\
    &=\widetilde X^T C^TBC\widetilde Y.
\end{aligned}
\]
Значит, матрица формы в новом базисе равна
\[
    \widetilde B=C^TBC.
\]
\end{proof}


\begin{definition}[Левое и правое ядро билинейной формы]
Пусть $\beta$ --- билинейная форма на $V$.
Определим
\[
    \Ker_L\beta=\{x\in V\mid \beta(x,y)=0\quad \forall y\in V\},
\]
\[
    \Ker_R\beta=\{y\in V\mid \beta(x,y)=0\quad \forall x\in V\}.
\]
\begin{knowledgefix}
Фиксируем обозначения: \(\Ker_L\beta=\{x\in V\mid \beta(x,y)=0\ \forall y\in V\}\), \(\Ker_R\beta=\{y\in V\mid \beta(x,y)=0\ \forall x\in V\}\). В матрице \(B\) это соответственно \(X^TB=0\) и \(BY=0\); поэтому при конечной размерности обе размерности равны \(n-\rk B\).
\end{knowledgefix}
\end{definition}

\begin{proposition}
$\Ker_L\beta$ и $\Ker_R\beta$ являются подпространствами пространства $V$.
\end{proposition}

\begin{proof}
Для левого ядра:
если $x_1,x_2\in\Ker_L\beta$, то для любого $y\in V$
\[
    \beta(x_1+x_2,y)=\beta(x_1,y)+\beta(x_2,y)=0.
\]
Если $\alpha\in\F$ и $x\in\Ker_L\beta$, то
\[
    \beta(\alpha x,y)=\alpha\beta(x,y)=0.
\]
Значит, $\Ker_L\beta$ --- подпространство. Для правого ядра доказательство аналогично.
\end{proof}

Пусть в некотором базисе матрица формы равна $B$, а
\[
    \beta(x,y)=X^TBY.
\]
Тогда условия принадлежности ядрам выражаются матрично:
\[
    x\in\Ker_L\beta
    \quad\Longleftrightarrow\quad
    X^TB=0
    \quad\Longleftrightarrow\quad
    B^TX=0,
\]
\[
    y\in\Ker_R\beta
    \quad\Longleftrightarrow\quad
    BY=0.
\]
Отсюда
\[
    \dim \Ker_L\beta=n-\rk B,
    \qquad
    \dim \Ker_R\beta=n-\rk B.
\]
В частности,
\[
    \dim \Ker_L\beta=\dim \Ker_R\beta.
\]

\begin{definition}[Невырожденная билинейная форма]
Билинейная форма $\beta$ называется \emph{невырожденной}, если её ядро тривиально.
В конечномерном случае это равносильно условиям
\[
    \Ker_L\beta=0,
    \qquad
    \Ker_R\beta=0,
\]
а также условию
\[
    \det B\ne 0.
\]
\begin{knowledgefix}
Используем стандартную формулировку через два ядра: левое \(\Ker_L\beta\) зануляет форму по первому аргументу, правое \(\Ker_R\beta\) — по второму. Билинейная форма невырождена тогда и только тогда, когда эти ядра нулевые, что эквивалентно \(\det B\ne0\) в некотором/любом базисе.
\end{knowledgefix}
\end{definition}

\begin{corollary}
Если одно из ядер невырожденной формы равно нулю, то и другое ядро равно нулю.
\end{corollary}

\subsubsection*{Пример}

На странице приведён вычислительный пример билинейной формы с матрицей
\[
    B=\begin{pmatrix}
        1 & 1 & 2\\
        -1 & 1 & -1\\
        1 & 3 & 1
    \end{pmatrix}
\]
\begin{knowledgefix}
Для матрицы билинейной формы точная общая схема такая: если \(B=(\beta(e_i,e_j))\), координаты новых базисных векторов образуют столбцы матрицы \(C\), то \(B'=C^TBC\). Ранг формы равен рангу этой матрицы и не зависит от базиса, так как умножение на невырожденные \(C^T\) и \(C\) не меняет ранг.
\end{knowledgefix}

Идея примера: найти ядро формы через систему линейных уравнений, соответствующую матрице $B$ или $B^T$.

\subsection*{Дополнительные уточнения}

\subsubsection*{1. Билинейные формы, сопряжённое пространство и невырожденность}
Пусть $V$ --- конечномерное пространство над полем $\F$, а
\[
    \beta:V\times V\to \F
\]
--- билинейная форма. При фиксации второго аргумента $v\in V$ получаем линейный функционал
\[
    x\mapsto \beta(x,v),
\]
поэтому билинейной форме соответствует линейное отображение
\[
    \Phi_\beta:V\to V^*,\qquad
    \Phi_\beta(v)(x)=\beta(x,v).
\]
Это даёт изоморфизм
\[
    \operatorname{Bil}(V)\simeq \Hom(V,V^*).
\]

В этой конвенции правое ядро формы записывается как
\[
    \Ker_R\beta=\{v\in V\mid \beta(x,v)=0\ \forall x\in V\}=\Ker \Phi_\beta.
\]
Аналогично, если фиксировать первый аргумент, возникает отображение
\[
    \Psi_\beta:V\to V^*,\qquad \Psi_\beta(u)(y)=\beta(u,y),
\]
и левое ядро
\[
    \Ker_L\beta=\{u\in V\mid \beta(u,y)=0\ \forall y\in V\}=\Ker\Psi_\beta.
\]
Для конечномерного пространства следующие условия эквивалентны:
\[
    \beta\text{ невырождена}
    \quad\Longleftrightarrow\quad
    \Phi_\beta\text{ --- изоморфизм}
    \quad\Longleftrightarrow\quad
    \det B\ne 0
    \quad\Longleftrightarrow\quad
    \Ker_L\beta=\Ker_R\beta=0.
\]

Если на $V$ задано скалярное произведение $(\cdot,\cdot)$, то оно задаёт изоморфизм
\[
    V\xrightarrow{\sim} V^*,\qquad v\mapsto (\cdot,v).
\]
В ортонормированном базисе $e_1,\ldots,e_n$ вектор $e_i$ переходит в сопряжённый ковектор $e^i$, поскольку
\[
    e^i(e_j)=\delta^i_j,
    \qquad
    (e_j,e_i)=\delta_{ji}.
\]

В евклидовом пространстве билинейной форме $\beta$ соответствует единственный оператор $A\in\End(V)$ такой, что
\[
    \beta(x,y)=(x,Ay).
\]
В ортонормированном базисе матрица оператора $A$ совпадает с матрицей формы $\beta$.
Сопряжённый оператор задаётся условием
\[
    (Ax,y)=(x,A^*y).
\]
Тогда
\[
    \beta(x,y)=\beta(y,x)\Longleftrightarrow A^*=A,
    \qquad
    \beta(x,y)=-\beta(y,x)\Longleftrightarrow A^*=-A.
\]

\clearpage

\section{Билет 2. Симметрические и кососимметрические билинейные формы, их матрицы. Ортогональное дополнение к подпространству относительно билинейной формы, его свойства.}

\paragraph{Содержание билета.} Симметрические и кососимметрические билинейные формы, их матрицы. Ортогональное дополнение к подпространству относительно билинейной формы, его свойства.

\begin{planadd}

в ответе удобно сначала отделить симметричность и кососимметричность на языке форм, затем на языке матриц, а потом перейти к $U^{\perp_\beta}$ и его размерностным свойствам.

\end{planadd}

\subsubsection{Билинейные формы}


\begin{definition}[Билинейная форма]
Пусть $V$ --- векторное пространство над полем $\F$.
Отображение
\[
    \beta: V \times V \to \F
\]
называется \emph{билинейной формой}, если оно линейно по каждому аргументу:
для любых $x,x_1,x_2,y,y_1,y_2\in V$ и $\alpha\in\F$ выполнено
\begin{enumerate}[label=\arabic*)]
    \item $\beta(x_1+x_2,y)=\beta(x_1,y)+\beta(x_2,y)$;
    \item $\beta(\alpha x,y)=\alpha\beta(x,y)$;
    \item $\beta(x,y_1+y_2)=\beta(x,y_1)+\beta(x,y_2)$;
    \item $\beta(x,\alpha y)=\alpha\beta(x,y)$.
\end{enumerate}
\end{definition}

\begin{example}
Нулевая форма:
\[
    \beta(x,y)=0.
\]
\end{example}

\begin{example}
Пусть $V=\F^n$, $x=(x_1,\ldots,x_n)^T$, $y=(y_1,\ldots,y_n)^T$.
Тогда пример билинейной формы:
\[
    \beta(x,y)=x_1y_1+\cdots+x_ny_n=\sum_{i=1}^n x_i y_i.
\]
\end{example}

\begin{example}
Если $f,g\in V^*$, то
\[
    \beta(x,y)=f(x)g(y)
\]
является билинейной формой.
\end{example}

\begin{example}
Для матриц $A,B\in M_n(\F)$ в конспекте приведён пример вида
\[
    \beta(A,B)=\tr(AB).
\]
\begin{knowledgefix}
На пространстве матриц стандартный проверенный пример билинейной формы записывается как
\(\langle A,B\rangle=\tr(A^{T}B)\) над \(\R\) и как \(\tr(A^{*}B)\) над \(\C\). Запись \(\tr(AB)\) тоже билинейна, но не является стандартным положительным скалярным произведением на всех матрицах.
\end{knowledgefix}
\end{example}

\begin{example}
На странице приведён пример с определителем матрицы $2\times 2$.
\begin{knowledgefix}
Определитель одной матрицы \(\det A\) не является билинейной формой от пары матриц. Корректный пример с определителем: \((x,y)\mapsto \det(x,y)\) на \(\R^2\) — кососимметрическая билинейная форма от двух векторов; в размерности \(n\) \(\det(v_1,\ldots,v_n)\) — полилинейная антисимметрическая форма.
\end{knowledgefix}
\end{example}

\begin{example}
Если $f,g\in C[a,b]$, то
\[
    \beta(f,g)=\int_a^b f(x)g(x)\,dx
\]
является билинейной формой.
\end{example}


\begin{proposition}[Операции с билинейными формами]
Если $\beta,\beta_1,\beta_2$ --- билинейные формы на $V$, а $\alpha\in\F$, то формы
\[
    \alpha\beta,
    \qquad
    \beta_1+\beta_2
\]
задаются правилами
\[
    (\alpha\beta)(x,y)=\alpha\beta(x,y),
\]
\[
    (\beta_1+\beta_2)(x,y)=\beta_1(x,y)+\beta_2(x,y),
\]
и снова являются билинейными формами.
\end{proposition}

\begin{remark}
Билинейные формы на $V$ образуют векторное пространство.
Обозначение из конспекта:
\[
    \BL(V).
\]
\end{remark}

\begin{definition}[Матрица билинейной формы]
Пусть $\dim V=n$ и $e_1,\ldots,e_n$ --- базис пространства $V$.
Для билинейной формы $\beta$ положим
\[
    b_{ij}=\beta(e_i,e_j).
\]
Матрица
\[
    B=(b_{ij})_{i,j=1}^n
\]
называется матрицей билинейной формы $\beta$ в базисе $e_1,\ldots,e_n$.
\end{definition}

Если
\[
    x=e_1x_1+\cdots+e_nx_n,\qquad
    y=e_1y_1+\cdots+e_ny_n,
\]
то по билинейности
\[
\begin{aligned}
    \beta(x,y)
    &=\beta(e_1x_1+\cdots+e_nx_n,
             e_1y_1+\cdots+e_ny_n) \\
    &=\sum_{i=1}^n\sum_{j=1}^n x_i\beta(e_i,e_j)y_j.
\end{aligned}
\]
Если $X=(x_1,\ldots,x_n)^T$, $Y=(y_1,\ldots,y_n)^T$, то
\[
    \beta(x,y)=X^TBY.
\]

\begin{corollary}
Соответствие
\[
    \beta \longmapsto B=(\beta(e_i,e_j))
\]
задаёт биекцию между пространством билинейных форм на $V$ и пространством матриц $M_n(\F)$.
Поэтому
\[
    \dim \BL(V)=n^2.
\]
\end{corollary}


\begin{proposition}[Изменение матрицы билинейной формы при замене базиса]
Пусть $e=(e_1,\ldots,e_n)$ и $\widetilde e=(\widetilde e_1,\ldots,\widetilde e_n)$ --- два базиса пространства $V$.
Пусть матрица перехода $C$ задана соотношением
\[
    \widetilde e = eC.
\]
Если $X,Y$ --- координатные столбцы в старом базисе, а $\widetilde X,\widetilde Y$ --- координатные столбцы в новом базисе, то
\[
    X=C\widetilde X,
    \qquad
    Y=C\widetilde Y.
\]
Тогда матрицы билинейной формы в этих базисах связаны формулой
\[
    \widetilde B=C^TBC.
\]
\end{proposition}

\begin{proof}
В старом базисе
\[
    \beta(x,y)=X^TBY.
\]
После замены координат $X=C\widetilde X$, $Y=C\widetilde Y$ получаем
\[
\begin{aligned}
    \beta(x,y)
    &=(C\widetilde X)^T B(C\widetilde Y) \\
    &=\widetilde X^T C^TBC\widetilde Y.
\end{aligned}
\]
Значит, матрица формы в новом базисе равна
\[
    \widetilde B=C^TBC.
\]
\end{proof}


\begin{definition}[Левое и правое ядро билинейной формы]
Пусть $\beta$ --- билинейная форма на $V$.
Определим
\[
    \Ker_L\beta=\{x\in V\mid \beta(x,y)=0\quad \forall y\in V\},
\]
\[
    \Ker_R\beta=\{y\in V\mid \beta(x,y)=0\quad \forall x\in V\}.
\]
\begin{knowledgefix}
Фиксируем обозначения: \(\Ker_L\beta=\{x\in V\mid \beta(x,y)=0\ \forall y\in V\}\), \(\Ker_R\beta=\{y\in V\mid \beta(x,y)=0\ \forall x\in V\}\). В матрице \(B\) это соответственно \(X^TB=0\) и \(BY=0\); поэтому при конечной размерности обе размерности равны \(n-\rk B\).
\end{knowledgefix}
\end{definition}

\begin{proposition}
$\Ker_L\beta$ и $\Ker_R\beta$ являются подпространствами пространства $V$.
\end{proposition}

\begin{proof}
Для левого ядра:
если $x_1,x_2\in\Ker_L\beta$, то для любого $y\in V$
\[
    \beta(x_1+x_2,y)=\beta(x_1,y)+\beta(x_2,y)=0.
\]
Если $\alpha\in\F$ и $x\in\Ker_L\beta$, то
\[
    \beta(\alpha x,y)=\alpha\beta(x,y)=0.
\]
Значит, $\Ker_L\beta$ --- подпространство. Для правого ядра доказательство аналогично.
\end{proof}

Пусть в некотором базисе матрица формы равна $B$, а
\[
    \beta(x,y)=X^TBY.
\]
Тогда условия принадлежности ядрам выражаются матрично:
\[
    x\in\Ker_L\beta
    \quad\Longleftrightarrow\quad
    X^TB=0
    \quad\Longleftrightarrow\quad
    B^TX=0,
\]
\[
    y\in\Ker_R\beta
    \quad\Longleftrightarrow\quad
    BY=0.
\]
Отсюда
\[
    \dim \Ker_L\beta=n-\rk B,
    \qquad
    \dim \Ker_R\beta=n-\rk B.
\]
В частности,
\[
    \dim \Ker_L\beta=\dim \Ker_R\beta.
\]

\begin{definition}[Невырожденная билинейная форма]
Билинейная форма $\beta$ называется \emph{невырожденной}, если её ядро тривиально.
В конечномерном случае это равносильно условиям
\[
    \Ker_L\beta=0,
    \qquad
    \Ker_R\beta=0,
\]
а также условию
\[
    \det B\ne 0.
\]
\begin{knowledgefix}
Используем стандартную формулировку через два ядра: левое \(\Ker_L\beta\) зануляет форму по первому аргументу, правое \(\Ker_R\beta\) — по второму. Билинейная форма невырождена тогда и только тогда, когда эти ядра нулевые, что эквивалентно \(\det B\ne0\) в некотором/любом базисе.
\end{knowledgefix}
\end{definition}

\begin{corollary}
Если одно из ядер невырожденной формы равно нулю, то и другое ядро равно нулю.
\end{corollary}

\subsubsection*{Пример}

На странице приведён вычислительный пример билинейной формы с матрицей
\[
    B=\begin{pmatrix}
        1 & 1 & 2\\
        -1 & 1 & -1\\
        1 & 3 & 1
    \end{pmatrix}
\]
\begin{knowledgefix}
Для матрицы билинейной формы точная общая схема такая: если \(B=(\beta(e_i,e_j))\), координаты новых базисных векторов образуют столбцы матрицы \(C\), то \(B'=C^TBC\). Ранг формы равен рангу этой матрицы и не зависит от базиса, так как умножение на невырожденные \(C^T\) и \(C\) не меняет ранг.
\end{knowledgefix}

Идея примера: найти ядро формы через систему линейных уравнений, соответствующую матрице $B$ или $B^T$.

\subsubsection{Симметрические, кососимметрические и квадратичные формы}


\begin{definition}[Симметрическая билинейная форма]
Билинейная форма $\beta$ называется \emph{симметрической}, если
\[
    \beta(x,y)=\beta(y,x)
    \qquad \forall x,y\in V.
\]
В матричной форме это условие соответствует
\[
    B^T=B.
\]
\end{definition}

\begin{definition}[Кососимметрическая билинейная форма]
Билинейная форма $\beta$ называется \emph{кососимметрической}, если
\[
    \beta(x,y)=-\beta(y,x)
    \qquad \forall x,y\in V.
\]
В матричной форме это условие соответствует
\[
    B^T=-B.
\]
\end{definition}

\begin{remark}
Если $\beta$ кососимметрична, то
\[
    \beta(x,x)=-\beta(x,x),
\]
поэтому
\[
    2\beta(x,x)=0.
\]
Если $\operatorname{char}\F\ne 2$, то отсюда
\[
    \beta(x,x)=0.
\]
\end{remark}

\begin{proposition}[Разложение формы на симметрическую и кососимметрическую части]
Если $\operatorname{char}\F\ne2$, то всякая билинейная форма $\beta$ раскладывается в сумму симметрической и кососимметрической форм:
\[
    \beta=\beta^+ + \beta^-,
\]
где
\[
    \beta^+(x,y)=\frac{\beta(x,y)+\beta(y,x)}{2},
    \qquad
    \beta^-(x,y)=\frac{\beta(x,y)-\beta(y,x)}{2}.
\]
\end{proposition}

\begin{remark}
В конспекте используются обозначения вида
\[
    \BL^+(V),\qquad \BL^-(V),
\]
для пространств симметрических и кососимметрических билинейных форм.
\begin{knowledgefix}
Для невырожденной симметрической формы выполняется \(\dim U^\perp=n-\dim U\). Если ограничение \(\beta|_U\) невырождено, то \(U\cap U^\perp=0\), а значит \(V=U\oplus U^\perp\). Если форма вырождена или несимметрична, нужно различать левое и правое дополнения.
\end{knowledgefix}
\end{remark}


\begin{definition}[Квадратичная форма]
Квадратичной формой, связанной с билинейной формой $\beta$, называется отображение
\[
    q(x)=\beta(x,x).
\]
\end{definition}

\begin{example}
На странице приведён пример квадратичной формы от нескольких переменных.
\begin{knowledgefix}
Для иллюстрации метода рассмотрим форму \(q(x,y)=x^2+2xy+3y^2\) матрица равна \(\begin{pmatrix}1&1\\1&3\end{pmatrix}\), а приведение Лагранжем даёт \(q=(x+y)^2+2y^2\). Этот пример показывает сам алгоритм приведения.
\end{knowledgefix}
\end{example}

Если $\operatorname{char}\F\ne2$ и $q(x)=\beta(x,x)$ для симметрической билинейной формы $\beta$, то $\beta$ восстанавливается по $q$ формулой поляризации:
\[
    \beta(x,y)=\frac{q(x+y)-q(x)-q(y)}{2}.
\]

\begin{remark}
В конспекте отмечено, что соответствие
\[
    \BL^+(V)\longrightarrow \{\text{квадратичные формы на }V\},
    \qquad
    \beta\longmapsto q,\quad q(x)=\beta(x,x),
\]
является биекцией при $\operatorname{char}\F\ne2$.
\begin{knowledgefix}
Связь с отображением \(\Phi_\beta:V\to V^*\), \(\Phi_\beta(x)=\beta(x,\cdot)\): правое/левое ядро зависит от выбранного порядка аргументов. При соглашении \(\Phi_\beta(x)(y)=\beta(x,y)\) ядро \(\Phi_\beta\) есть левое ядро формы.
\end{knowledgefix}
\end{remark}

\subsubsection*{Ортогональное дополнение относительно билинейной формы}

\begin{definition}[Ортогональное дополнение]
Пусть $U\le V$. Для билинейной формы $\beta$ определим
\[
    U^\perp=\{y\in V\mid \beta(x,y)=0\quad \forall x\in U\}.
\]
\begin{knowledgefix}
Для общей билинейной формы вводим правое дополнение \(U^\perp_R=\{v\mid \beta(u,v)=0\ \forall u\in U\}\) и левое \({}^\perp U=\{v\mid \beta(v,u)=0\ \forall u\in U\}\). Для симметрической формы они совпадают, поэтому можно писать просто \(U^\perp\).
\end{knowledgefix}
\end{definition}

\begin{lemma}
Пусть $\beta$ --- невырожденная билинейная форма на конечномерном пространстве $V$.
Тогда для любого подпространства $U\le V$
\[
    \dim U^\perp=\dim V-\dim U.
\]
\end{lemma}

\begin{proof}
Пусть $e_1,\ldots,e_k$ --- базис $U$, дополненный до базиса
$e_1,\ldots,e_k,e_{k+1},\ldots,e_n$ пространства $V$.
Условие $y\in U^\perp$ эквивалентно системе
\[
    \beta(e_i,y)=0,
    \qquad i=1,\ldots,k.
\]
В координатах это система $k$ линейных уравнений.
Так как форма невырождена, соответствующая система имеет ранг $k$.
Следовательно,
\[
    \dim U^\perp=n-k=\dim V-\dim U.
\]
\begin{knowledgefix}
Блочная запись свойства ортогонального дополнения: если в базисе, согласованном с \(U\), матрица ограничения на \(U\) невырождена, то после выделения координат \(u\) и \(w\) условие \(\beta(u,w)=0\) задаёт систему, из которой следует \(V=U\oplus U^\perp\).
\end{knowledgefix}
\end{proof}

\begin{corollary}
Для невырожденной формы
\[
    (U^\perp)^\perp=U.
\]
\end{corollary}

\begin{definition}[Невырожденное подпространство]
Подпространство $U\le V$ называется невырожденным относительно формы $\beta$, если ограничение формы на $U$ невырождено, то есть
\[
    \Ker(\beta|_U)=0.
\]
\begin{knowledgefix}
Обозначение \(\beta|_U\) означает ограничение формы на \(U\times U\). Подпространство \(U\) невырождено относительно \(\beta\), если \(\Ker(\beta|_U)=0\), то есть \(U\cap U^\perp=0\) в симметрическом случае.
\end{knowledgefix}
\end{definition}

\begin{proposition}
Для невырожденной билинейной формы $\beta$ и подпространства $U\le V$ следующие условия эквивалентны:
\begin{enumerate}[label=\arabic*)]
    \item $U$ невырождено;
    \item $U\cap U^\perp=0$;
    \item $V=U\oplus U^\perp$.
\end{enumerate}
\end{proposition}

\begin{proof}
Из определения ядра ограничения формы:
\[
    \Ker(\beta|_U)=U\cap U^\perp.
\]
Поэтому $U$ невырождено тогда и только тогда, когда $U\cap U^\perp=0$.
С учётом размерностной формулы
\[
    \dim U^\perp=\dim V-\dim U
\]
получаем
\[
    \dim U+\dim U^\perp=\dim V.
\]
Значит, условие $U\cap U^\perp=0$ равносильно прямой сумме
\[
    V=U\oplus U^\perp.
\]
\end{proof}

\begin{remark}
Если $V=U\oplus U^\perp$, то в базисе, составленном из базиса $U$ и базиса $U^\perp$, матрица формы имеет блочно-диагональный вид.
\end{remark}

\subsubsection*{Ортогональная система}

\begin{definition}[Ортогональная система]
Система векторов называется ортогональной относительно симметрической билинейной формы $\beta$, если
\[
    \beta(e_i,e_j)=0
    \qquad \text{при } i\ne j.
\]
\end{definition}

Если $\beta\in\BL^+(V)$ и $e_1,
\ldots,e_n$ --- ортогональная система, то матрица формы в этом базисе диагональна:
\[
    B=\diag\bigl(\beta(e_1,e_1),\ldots,\beta(e_n,e_n)\bigr).
\]
В этом случае
\[
    q(x)=\beta(x,x)=\sum_{i=1}^n \beta(e_i,e_i)x_i^2.
\]

\clearpage

\section{Билет 3. Квадратичные формы, поляризация. Канонический и нормальный виды симметрической билинейной и квадратичной формы.}

\paragraph{Содержание билета.} Квадратичные формы, поляризация. Канонический и нормальный виды симметрической билинейной и квадратичной формы.

\begin{planadd}

держи связку: квадратичная форма $q(x)=\beta(x,x)$, восстановление через поляризацию, затем канонический/нормальный вид.

\end{planadd}

\subsubsection{Симметрические, кососимметрические и квадратичные формы}


\begin{definition}[Симметрическая билинейная форма]
Билинейная форма $\beta$ называется \emph{симметрической}, если
\[
    \beta(x,y)=\beta(y,x)
    \qquad \forall x,y\in V.
\]
В матричной форме это условие соответствует
\[
    B^T=B.
\]
\end{definition}

\begin{definition}[Кососимметрическая билинейная форма]
Билинейная форма $\beta$ называется \emph{кососимметрической}, если
\[
    \beta(x,y)=-\beta(y,x)
    \qquad \forall x,y\in V.
\]
В матричной форме это условие соответствует
\[
    B^T=-B.
\]
\end{definition}

\begin{remark}
Если $\beta$ кососимметрична, то
\[
    \beta(x,x)=-\beta(x,x),
\]
поэтому
\[
    2\beta(x,x)=0.
\]
Если $\operatorname{char}\F\ne 2$, то отсюда
\[
    \beta(x,x)=0.
\]
\end{remark}

\begin{proposition}[Разложение формы на симметрическую и кососимметрическую части]
Если $\operatorname{char}\F\ne2$, то всякая билинейная форма $\beta$ раскладывается в сумму симметрической и кососимметрической форм:
\[
    \beta=\beta^+ + \beta^-,
\]
где
\[
    \beta^+(x,y)=\frac{\beta(x,y)+\beta(y,x)}{2},
    \qquad
    \beta^-(x,y)=\frac{\beta(x,y)-\beta(y,x)}{2}.
\]
\end{proposition}

\begin{remark}
В конспекте используются обозначения вида
\[
    \BL^+(V),\qquad \BL^-(V),
\]
для пространств симметрических и кососимметрических билинейных форм.
\begin{knowledgefix}
Для невырожденной симметрической формы выполняется \(\dim U^\perp=n-\dim U\). Если ограничение \(\beta|_U\) невырождено, то \(U\cap U^\perp=0\), а значит \(V=U\oplus U^\perp\). Если форма вырождена или несимметрична, нужно различать левое и правое дополнения.
\end{knowledgefix}
\end{remark}


\begin{definition}[Квадратичная форма]
Квадратичной формой, связанной с билинейной формой $\beta$, называется отображение
\[
    q(x)=\beta(x,x).
\]
\end{definition}

\begin{example}
На странице приведён пример квадратичной формы от нескольких переменных.
\begin{knowledgefix}
Для иллюстрации метода рассмотрим форму \(q(x,y)=x^2+2xy+3y^2\) матрица равна \(\begin{pmatrix}1&1\\1&3\end{pmatrix}\), а приведение Лагранжем даёт \(q=(x+y)^2+2y^2\). Этот пример показывает сам алгоритм приведения.
\end{knowledgefix}
\end{example}

Если $\operatorname{char}\F\ne2$ и $q(x)=\beta(x,x)$ для симметрической билинейной формы $\beta$, то $\beta$ восстанавливается по $q$ формулой поляризации:
\[
    \beta(x,y)=\frac{q(x+y)-q(x)-q(y)}{2}.
\]

\begin{remark}
В конспекте отмечено, что соответствие
\[
    \BL^+(V)\longrightarrow \{\text{квадратичные формы на }V\},
    \qquad
    \beta\longmapsto q,\quad q(x)=\beta(x,x),
\]
является биекцией при $\operatorname{char}\F\ne2$.
\begin{knowledgefix}
Связь с отображением \(\Phi_\beta:V\to V^*\), \(\Phi_\beta(x)=\beta(x,\cdot)\): правое/левое ядро зависит от выбранного порядка аргументов. При соглашении \(\Phi_\beta(x)(y)=\beta(x,y)\) ядро \(\Phi_\beta\) есть левое ядро формы.
\end{knowledgefix}
\end{remark}

\subsubsection*{Ортогональное дополнение относительно билинейной формы}

\begin{definition}[Ортогональное дополнение]
Пусть $U\le V$. Для билинейной формы $\beta$ определим
\[
    U^\perp=\{y\in V\mid \beta(x,y)=0\quad \forall x\in U\}.
\]
\begin{knowledgefix}
Для общей билинейной формы вводим правое дополнение \(U^\perp_R=\{v\mid \beta(u,v)=0\ \forall u\in U\}\) и левое \({}^\perp U=\{v\mid \beta(v,u)=0\ \forall u\in U\}\). Для симметрической формы они совпадают, поэтому можно писать просто \(U^\perp\).
\end{knowledgefix}
\end{definition}

\begin{lemma}
Пусть $\beta$ --- невырожденная билинейная форма на конечномерном пространстве $V$.
Тогда для любого подпространства $U\le V$
\[
    \dim U^\perp=\dim V-\dim U.
\]
\end{lemma}

\begin{proof}
Пусть $e_1,\ldots,e_k$ --- базис $U$, дополненный до базиса
$e_1,\ldots,e_k,e_{k+1},\ldots,e_n$ пространства $V$.
Условие $y\in U^\perp$ эквивалентно системе
\[
    \beta(e_i,y)=0,
    \qquad i=1,\ldots,k.
\]
В координатах это система $k$ линейных уравнений.
Так как форма невырождена, соответствующая система имеет ранг $k$.
Следовательно,
\[
    \dim U^\perp=n-k=\dim V-\dim U.
\]
\begin{knowledgefix}
Блочная запись свойства ортогонального дополнения: если в базисе, согласованном с \(U\), матрица ограничения на \(U\) невырождена, то после выделения координат \(u\) и \(w\) условие \(\beta(u,w)=0\) задаёт систему, из которой следует \(V=U\oplus U^\perp\).
\end{knowledgefix}
\end{proof}

\begin{corollary}
Для невырожденной формы
\[
    (U^\perp)^\perp=U.
\]
\end{corollary}

\begin{definition}[Невырожденное подпространство]
Подпространство $U\le V$ называется невырожденным относительно формы $\beta$, если ограничение формы на $U$ невырождено, то есть
\[
    \Ker(\beta|_U)=0.
\]
\begin{knowledgefix}
Обозначение \(\beta|_U\) означает ограничение формы на \(U\times U\). Подпространство \(U\) невырождено относительно \(\beta\), если \(\Ker(\beta|_U)=0\), то есть \(U\cap U^\perp=0\) в симметрическом случае.
\end{knowledgefix}
\end{definition}

\begin{proposition}
Для невырожденной билинейной формы $\beta$ и подпространства $U\le V$ следующие условия эквивалентны:
\begin{enumerate}[label=\arabic*)]
    \item $U$ невырождено;
    \item $U\cap U^\perp=0$;
    \item $V=U\oplus U^\perp$.
\end{enumerate}
\end{proposition}

\begin{proof}
Из определения ядра ограничения формы:
\[
    \Ker(\beta|_U)=U\cap U^\perp.
\]
Поэтому $U$ невырождено тогда и только тогда, когда $U\cap U^\perp=0$.
С учётом размерностной формулы
\[
    \dim U^\perp=\dim V-\dim U
\]
получаем
\[
    \dim U+\dim U^\perp=\dim V.
\]
Значит, условие $U\cap U^\perp=0$ равносильно прямой сумме
\[
    V=U\oplus U^\perp.
\]
\end{proof}

\begin{remark}
Если $V=U\oplus U^\perp$, то в базисе, составленном из базиса $U$ и базиса $U^\perp$, матрица формы имеет блочно-диагональный вид.
\end{remark}

\subsubsection*{Ортогональная система}

\begin{definition}[Ортогональная система]
Система векторов называется ортогональной относительно симметрической билинейной формы $\beta$, если
\[
    \beta(e_i,e_j)=0
    \qquad \text{при } i\ne j.
\]
\end{definition}

Если $\beta\in\BL^+(V)$ и $e_1,
\ldots,e_n$ --- ортогональная система, то матрица формы в этом базисе диагональна:
\[
    B=\diag\bigl(\beta(e_1,e_1),\ldots,\beta(e_n,e_n)\bigr).
\]
В этом случае
\[
    q(x)=\beta(x,x)=\sum_{i=1}^n \beta(e_i,e_i)x_i^2.
\]

\subsubsection{Канонический вид квадратичной формы}


\begin{theorem}[Существование ортогонального базиса]
Пусть $\operatorname{char}\F\ne2$ и $\beta$ --- симметрическая билинейная форма на конечномерном пространстве $V$.
Тогда существует ортогональный базис пространства $V$ относительно формы $\beta$.
\end{theorem}

\begin{proof}
Доказательство ведётся индукцией по $n=\dim V$.

Если $n=1$, утверждение очевидно.

Пусть $n>1$ и утверждение уже доказано для пространств меньшей размерности.
Рассмотрим квадратичную форму
\[
    q(x)=\beta(x,x).
\]
Если $q(x)=0$ для всех $x\in V$, то по формуле поляризации
\[
    \beta(x,y)=\frac{q(x+y)-q(x)-q(y)}{2}=0,
\]
то есть форма нулевая, и любой базис является ортогональным.

Если же существует $e_1\in V$ такой, что
\[
    q(e_1)=\beta(e_1,e_1)\ne0,
\]
то подпространство $\langle e_1\rangle$ невырождено. Тогда
\[
    V=\langle e_1\rangle\oplus \langle e_1\rangle^\perp.
\]
К подпространству $\langle e_1\rangle^\perp$ применяем индукционное предположение и получаем в нём ортогональный базис
\[
    e_2,\ldots,e_n.
\]
Тогда
\[
    e_1,e_2,\ldots,e_n
\]
--- ортогональный базис всего $V$.
\end{proof}

\begin{corollary}
Любую квадратичную форму можно привести к каноническому виду.
\end{corollary}

\subsubsection*{Метод Лагранжа}

Рассмотрим квадратичную форму
\[
    q(x_1,\ldots,x_n)
    =\sum_{i=1}^n b_{ii}x_i^2
    +2\sum_{1\le i<j\le n} b_{ij}x_ix_j.
\]
Если $b_{11}\ne0$, то выделяем полный квадрат:
\[
\begin{aligned}
    q(x_1,\ldots,x_n)
    &=b_{11}\left(x_1+\frac{b_{12}}{b_{11}}x_2+\cdots+\frac{b_{1n}}{b_{11}}x_n\right)^2 \\
    &\quad + q_1(x_2,\ldots,x_n),
\end{aligned}
\]
где $q_1$ --- квадратичная форма от $n-1$ переменных.
Далее применяем тот же процесс к $q_1$.

Если $b_{11}=0$, но есть ненулевой коэффициент при смешанном члене, то выполняют линейную замену переменных, чтобы появился ненулевой квадрат.
\begin{knowledgefix}
Стандартный ход метода Лагранжа: если при \(x_1^2\) коэффициент \(a_{11}\ne0\), выделяют полный квадрат
\[q=a_{11}\left(x_1+\frac{a_{12}}{a_{11}}x_2+\cdots+\frac{a_{1n}}{a_{11}}x_n\right)^2+q_1(x_2,
\ldots,x_n),\]
после чего повторяют процедуру для \(q_1\). Если подходящего квадратного коэффициента нет, сначала делают линейную замену, создающую ненулевой квадрат.
\end{knowledgefix}

\begin{remark}
В результате метода Лагранжа получается новый базис, в котором матрица квадратичной формы диагональна.
\end{remark}

\subsubsection*{Метод Якоби}

Пусть $B$ --- матрица симметрической билинейной формы в базисе $e_1,\ldots,e_n$.
Обозначим угловые миноры:
\[
    \Delta_k=\det B_k,
\]
где $B_k$ --- левый верхний $k\times k$ блок матрицы $B$, и положим
\[
    \Delta_0=1.
\]

\begin{theorem}[Формула Якоби]
Если
\[
    \Delta_k\ne0,
    \qquad k=1,\ldots,n,
\]
то существует ортогональный базис $f_1,\ldots,f_n$, в котором квадратичная форма имеет вид
\[
    q(x_1,\ldots,x_n)
    =\frac{\Delta_1}{\Delta_0}x_1^2
    +\frac{\Delta_2}{\Delta_1}x_2^2
    +\cdots+
    \frac{\Delta_n}{\Delta_{n-1}}x_n^2.
\]
\begin{knowledgefix}
Точная стандартная формулировка — метод Якоби: если ведущие главные миноры \(\Delta_k\ne0\), то квадратичная форма приводится треугольной заменой к виду
\[q=\Delta_1 y_1^2+\frac{\Delta_2}{\Delta_1}y_2^2+\cdots+\frac{\Delta_n}{\Delta_{n-1}}y_n^2.\]
Эта же формула используется для быстрого определения инерции при невырожденных ведущих минорах.
\end{knowledgefix}
\end{theorem}

\begin{proof}[Идея доказательства]
Пусть
\[
    V_k=\langle e_1,\ldots,e_k\rangle.
\]
Условие $\Delta_k\ne0$ означает, что ограничение формы на $V_k$ невырождено.
Тогда можно строить ортогональный базис последовательно:
\[
    V_1\subset V_2\subset\cdots\subset V_n=V,
\]
выбирая
\[
    f_k\in V_k\cap V_{k-1}^{\perp}
\]
так, чтобы
\[
    V_k=V_{k-1}\oplus\langle f_k\rangle.
\]
В этом базисе матрица формы диагональна.
Диагональные элементы определяются из соотношения для определителей:
\[
    \Delta_k=\prod_{i=1}^k \beta(f_i,f_i).
\]
Поэтому
\[
    \beta(f_k,f_k)=\frac{\Delta_k}{\Delta_{k-1}}.
\]
\begin{knowledgefix}
Смысл геометрического доказательства: симметрическая билинейная форма задаётся самосопряжённым оператором, поэтому в ортонормированном базисе из собственных векторов её матрица диагональна. Геометрически это переход к главным осям.
\end{knowledgefix}
\end{proof}

\subsubsection*{Положительная определённость и критерий Сильвестра}

\begin{definition}[Положительно определённая квадратичная форма]
Квадратичная форма $q$ называется \emph{положительно определённой}, если
\[
    q(x)>0
    \qquad \forall x\ne0.
\]
\end{definition}

\begin{definition}[Отрицательно определённая квадратичная форма]
Квадратичная форма $q$ называется \emph{отрицательно определённой}, если
\[
    q(x)<0
    \qquad \forall x\ne0.
\]
\end{definition}

\begin{definition}[Полуопределённые формы]
Форма называется положительно полуопределённой, если
\[
    q(x)\ge0
    \qquad \forall x\in V,
\]
и отрицательно полуопределённой, если
\[
    q(x)\le0
    \qquad \forall x\in V.
\]
\begin{knowledgefix}
Классификация по знаку: положительно определённая \(q(x)>0\) при \(x\ne0\); отрицательно определённая \(q(x)<0\); положительно полуопределённая \(q(x)\ge0\); отрицательно полуопределённая \(q(x)\le0\); неопределённая — принимает и положительные, и отрицательные значения.
\end{knowledgefix}
\end{definition}

\begin{theorem}[Критерий через канонический вид]
Пусть квадратичная форма приведена к каноническому виду
\[
    q(x_1,
    \ldots,x_n)=a_1x_1^2+
    \cdots+a_nx_n^2.
\]
Тогда:
\begin{enumerate}[label=\arabic*)]
    \item $q$ положительно определена тогда и только тогда, когда $a_i>0$ для всех $i$;
    \item $q$ отрицательно определена тогда и только тогда, когда $a_i<0$ для всех $i$;
    \item если среди $a_i$ есть положительные и отрицательные, то форма неопределённая.
\end{enumerate}
\end{theorem}

\begin{theorem}[Критерий Сильвестра]
Пусть $q$ --- вещественная квадратичная форма с симметрической матрицей $B$, а
\[
    \Delta_k
\]
--- её угловые миноры. Тогда $q$ положительно определена тогда и только тогда, когда
\[
    \Delta_1>0,
    \quad
    \Delta_2>0,
    \quad
    \ldots,
    \quad
    \Delta_n>0.
\]
\end{theorem}

\begin{proof}
Если все угловые миноры положительны, то по формуле Якоби
\[
    q(x_1,
    \ldots,x_n)=
    \frac{\Delta_1}{\Delta_0}x_1^2+
    \frac{\Delta_2}{\Delta_1}x_2^2+
    \cdots+
    \frac{\Delta_n}{\Delta_{n-1}}x_n^2.
\]
Так как все отношения
\[
    \frac{\Delta_k}{\Delta_{k-1}}
\]
положительны, форма положительно определена.
Обратное направление следует из того, что ограничение положительно определённой формы на любое подпространство также положительно определено, поэтому все соответствующие угловые миноры положительны.
\begin{knowledgefix}
Доказательство закона инерции через переходную матрицу: если две диагональные записи одной формы связаны невырожденной заменой \(x=Cy\), то максимальная размерность подпространства, на котором форма положительно определена, не меняется. Эта максимальная размерность равна числу положительных квадратов; аналогично для отрицательных. Поэтому числа положительных, отрицательных и нулевых коэффициентов не зависят от приведения.
\end{knowledgefix}
\end{proof}

\begin{corollary}[Критерий отрицательной определённости]
Форма $q$ отрицательно определена тогда и только тогда, когда
\[
    (-1)^k\Delta_k>0,
    \qquad k=1,
    \ldots,n.
\]
Иначе говоря,
\[
    \Delta_1<0,
    \quad
    \Delta_2>0,
    \quad
    \Delta_3<0,
    \quad
    \ldots
\]
\begin{knowledgefix}
Для отрицательной определённости критерий записывается как \(\Delta_1<0,\Delta_2>0,\Delta_3<0,\ldots\), то есть \((-1)^k\Delta_k>0\). Для положительной определённости: \(\Delta_k>0\) для всех \(k\).
\end{knowledgefix}
\end{corollary}

\begin{example}
На странице 19 приведён пример квадратичной формы от двух переменных и вычисление её матрицы/миноров.
\begin{knowledgefix}
Численный пример заменён проверенным стандартным: матрица \(\begin{pmatrix}2&1\\1&2\end{pmatrix}\) положительно определена, так как \(\Delta_1=2>0\), \(\Delta_2=3>0\). Матрица \(\begin{pmatrix}0&1\\1&0\end{pmatrix}\) неопределённа, поскольку \(q(x,y)=2xy\) принимает оба знака.
\end{knowledgefix}
\end{example}

\clearpage

\section{Билет 4. Методы Лагранжа и Якоби приведения квадратичной формы к каноническому виду. Закон инерции. Приведение кососимметрической билинейной формы к каноническому виду.}

\paragraph{Содержание билета.} Методы Лагранжа и Якоби приведения квадратичной формы к каноническому виду. Закон инерции. Приведение кососимметрической билинейной формы к каноническому виду.

\begin{planadd}

этот билет вычислительный; важно не просто назвать методы, а объяснить алгоритм: выделение квадратов Лагранжа, главные миноры Якоби, закон инерции и отдельный канонический вид кососимметрической формы.

\end{planadd}

\subsubsection{Канонический вид квадратичной формы}


\begin{theorem}[Существование ортогонального базиса]
Пусть $\operatorname{char}\F\ne2$ и $\beta$ --- симметрическая билинейная форма на конечномерном пространстве $V$.
Тогда существует ортогональный базис пространства $V$ относительно формы $\beta$.
\end{theorem}

\begin{proof}
Доказательство ведётся индукцией по $n=\dim V$.

Если $n=1$, утверждение очевидно.

Пусть $n>1$ и утверждение уже доказано для пространств меньшей размерности.
Рассмотрим квадратичную форму
\[
    q(x)=\beta(x,x).
\]
Если $q(x)=0$ для всех $x\in V$, то по формуле поляризации
\[
    \beta(x,y)=\frac{q(x+y)-q(x)-q(y)}{2}=0,
\]
то есть форма нулевая, и любой базис является ортогональным.

Если же существует $e_1\in V$ такой, что
\[
    q(e_1)=\beta(e_1,e_1)\ne0,
\]
то подпространство $\langle e_1\rangle$ невырождено. Тогда
\[
    V=\langle e_1\rangle\oplus \langle e_1\rangle^\perp.
\]
К подпространству $\langle e_1\rangle^\perp$ применяем индукционное предположение и получаем в нём ортогональный базис
\[
    e_2,\ldots,e_n.
\]
Тогда
\[
    e_1,e_2,\ldots,e_n
\]
--- ортогональный базис всего $V$.
\end{proof}

\begin{corollary}
Любую квадратичную форму можно привести к каноническому виду.
\end{corollary}

\subsubsection*{Метод Лагранжа}

Рассмотрим квадратичную форму
\[
    q(x_1,\ldots,x_n)
    =\sum_{i=1}^n b_{ii}x_i^2
    +2\sum_{1\le i<j\le n} b_{ij}x_ix_j.
\]
Если $b_{11}\ne0$, то выделяем полный квадрат:
\[
\begin{aligned}
    q(x_1,\ldots,x_n)
    &=b_{11}\left(x_1+\frac{b_{12}}{b_{11}}x_2+\cdots+\frac{b_{1n}}{b_{11}}x_n\right)^2 \\
    &\quad + q_1(x_2,\ldots,x_n),
\end{aligned}
\]
где $q_1$ --- квадратичная форма от $n-1$ переменных.
Далее применяем тот же процесс к $q_1$.

Если $b_{11}=0$, но есть ненулевой коэффициент при смешанном члене, то выполняют линейную замену переменных, чтобы появился ненулевой квадрат.
\begin{knowledgefix}
Стандартный ход метода Лагранжа: если при \(x_1^2\) коэффициент \(a_{11}\ne0\), выделяют полный квадрат
\[q=a_{11}\left(x_1+\frac{a_{12}}{a_{11}}x_2+\cdots+\frac{a_{1n}}{a_{11}}x_n\right)^2+q_1(x_2,
\ldots,x_n),\]
после чего повторяют процедуру для \(q_1\). Если подходящего квадратного коэффициента нет, сначала делают линейную замену, создающую ненулевой квадрат.
\end{knowledgefix}

\begin{remark}
В результате метода Лагранжа получается новый базис, в котором матрица квадратичной формы диагональна.
\end{remark}

\subsubsection*{Метод Якоби}

Пусть $B$ --- матрица симметрической билинейной формы в базисе $e_1,\ldots,e_n$.
Обозначим угловые миноры:
\[
    \Delta_k=\det B_k,
\]
где $B_k$ --- левый верхний $k\times k$ блок матрицы $B$, и положим
\[
    \Delta_0=1.
\]

\begin{theorem}[Формула Якоби]
Если
\[
    \Delta_k\ne0,
    \qquad k=1,\ldots,n,
\]
то существует ортогональный базис $f_1,\ldots,f_n$, в котором квадратичная форма имеет вид
\[
    q(x_1,\ldots,x_n)
    =\frac{\Delta_1}{\Delta_0}x_1^2
    +\frac{\Delta_2}{\Delta_1}x_2^2
    +\cdots+
    \frac{\Delta_n}{\Delta_{n-1}}x_n^2.
\]
\begin{knowledgefix}
Точная стандартная формулировка — метод Якоби: если ведущие главные миноры \(\Delta_k\ne0\), то квадратичная форма приводится треугольной заменой к виду
\[q=\Delta_1 y_1^2+\frac{\Delta_2}{\Delta_1}y_2^2+\cdots+\frac{\Delta_n}{\Delta_{n-1}}y_n^2.\]
Эта же формула используется для быстрого определения инерции при невырожденных ведущих минорах.
\end{knowledgefix}
\end{theorem}

\begin{proof}[Идея доказательства]
Пусть
\[
    V_k=\langle e_1,\ldots,e_k\rangle.
\]
Условие $\Delta_k\ne0$ означает, что ограничение формы на $V_k$ невырождено.
Тогда можно строить ортогональный базис последовательно:
\[
    V_1\subset V_2\subset\cdots\subset V_n=V,
\]
выбирая
\[
    f_k\in V_k\cap V_{k-1}^{\perp}
\]
так, чтобы
\[
    V_k=V_{k-1}\oplus\langle f_k\rangle.
\]
В этом базисе матрица формы диагональна.
Диагональные элементы определяются из соотношения для определителей:
\[
    \Delta_k=\prod_{i=1}^k \beta(f_i,f_i).
\]
Поэтому
\[
    \beta(f_k,f_k)=\frac{\Delta_k}{\Delta_{k-1}}.
\]
\begin{knowledgefix}
Смысл геометрического доказательства: симметрическая билинейная форма задаётся самосопряжённым оператором, поэтому в ортонормированном базисе из собственных векторов её матрица диагональна. Геометрически это переход к главным осям.
\end{knowledgefix}
\end{proof}

\subsubsection*{Положительная определённость и критерий Сильвестра}

\begin{definition}[Положительно определённая квадратичная форма]
Квадратичная форма $q$ называется \emph{положительно определённой}, если
\[
    q(x)>0
    \qquad \forall x\ne0.
\]
\end{definition}

\begin{definition}[Отрицательно определённая квадратичная форма]
Квадратичная форма $q$ называется \emph{отрицательно определённой}, если
\[
    q(x)<0
    \qquad \forall x\ne0.
\]
\end{definition}

\begin{definition}[Полуопределённые формы]
Форма называется положительно полуопределённой, если
\[
    q(x)\ge0
    \qquad \forall x\in V,
\]
и отрицательно полуопределённой, если
\[
    q(x)\le0
    \qquad \forall x\in V.
\]
\begin{knowledgefix}
Классификация по знаку: положительно определённая \(q(x)>0\) при \(x\ne0\); отрицательно определённая \(q(x)<0\); положительно полуопределённая \(q(x)\ge0\); отрицательно полуопределённая \(q(x)\le0\); неопределённая — принимает и положительные, и отрицательные значения.
\end{knowledgefix}
\end{definition}

\begin{theorem}[Критерий через канонический вид]
Пусть квадратичная форма приведена к каноническому виду
\[
    q(x_1,
    \ldots,x_n)=a_1x_1^2+
    \cdots+a_nx_n^2.
\]
Тогда:
\begin{enumerate}[label=\arabic*)]
    \item $q$ положительно определена тогда и только тогда, когда $a_i>0$ для всех $i$;
    \item $q$ отрицательно определена тогда и только тогда, когда $a_i<0$ для всех $i$;
    \item если среди $a_i$ есть положительные и отрицательные, то форма неопределённая.
\end{enumerate}
\end{theorem}

\begin{theorem}[Критерий Сильвестра]
Пусть $q$ --- вещественная квадратичная форма с симметрической матрицей $B$, а
\[
    \Delta_k
\]
--- её угловые миноры. Тогда $q$ положительно определена тогда и только тогда, когда
\[
    \Delta_1>0,
    \quad
    \Delta_2>0,
    \quad
    \ldots,
    \quad
    \Delta_n>0.
\]
\end{theorem}

\begin{proof}
Если все угловые миноры положительны, то по формуле Якоби
\[
    q(x_1,
    \ldots,x_n)=
    \frac{\Delta_1}{\Delta_0}x_1^2+
    \frac{\Delta_2}{\Delta_1}x_2^2+
    \cdots+
    \frac{\Delta_n}{\Delta_{n-1}}x_n^2.
\]
Так как все отношения
\[
    \frac{\Delta_k}{\Delta_{k-1}}
\]
положительны, форма положительно определена.
Обратное направление следует из того, что ограничение положительно определённой формы на любое подпространство также положительно определено, поэтому все соответствующие угловые миноры положительны.
\begin{knowledgefix}
Доказательство закона инерции через переходную матрицу: если две диагональные записи одной формы связаны невырожденной заменой \(x=Cy\), то максимальная размерность подпространства, на котором форма положительно определена, не меняется. Эта максимальная размерность равна числу положительных квадратов; аналогично для отрицательных. Поэтому числа положительных, отрицательных и нулевых коэффициентов не зависят от приведения.
\end{knowledgefix}
\end{proof}

\begin{corollary}[Критерий отрицательной определённости]
Форма $q$ отрицательно определена тогда и только тогда, когда
\[
    (-1)^k\Delta_k>0,
    \qquad k=1,
    \ldots,n.
\]
Иначе говоря,
\[
    \Delta_1<0,
    \quad
    \Delta_2>0,
    \quad
    \Delta_3<0,
    \quad
    \ldots
\]
\begin{knowledgefix}
Для отрицательной определённости критерий записывается как \(\Delta_1<0,\Delta_2>0,\Delta_3<0,\ldots\), то есть \((-1)^k\Delta_k>0\). Для положительной определённости: \(\Delta_k>0\) для всех \(k\).
\end{knowledgefix}
\end{corollary}

\begin{example}
На странице 19 приведён пример квадратичной формы от двух переменных и вычисление её матрицы/миноров.
\begin{knowledgefix}
Численный пример заменён проверенным стандартным: матрица \(\begin{pmatrix}2&1\\1&2\end{pmatrix}\) положительно определена, так как \(\Delta_1=2>0\), \(\Delta_2=3>0\). Матрица \(\begin{pmatrix}0&1\\1&0\end{pmatrix}\) неопределённа, поскольку \(q(x,y)=2xy\) принимает оба знака.
\end{knowledgefix}
\end{example}

\subsubsection{Кососимметрические билинейные формы}

\subsubsection*{Канонический вид кососимметрической формы}

Пусть $\beta\in\BL^-(V)$, то есть
\[
    \beta(x,y)=-\beta(y,x)
    \qquad \forall x,y\in V.
\]
Если $\operatorname{char}\F\ne2$, то
\[
    \beta(x,x)=0
    \qquad \forall x\in V.
\]

\begin{theorem}[Канонический вид кососимметрической билинейной формы]
Пусть $\operatorname{char}\F\ne2$ и $\beta$ --- кососимметрическая билинейная форма на конечномерном пространстве $V$.
Тогда существует базис, в котором матрица формы имеет блочно-диагональный вид
\[
    \begin{pmatrix}
        0 & 1 &&&&& 0\\
        -1& 0 &&&&& 0\\
        && \ddots &&&& \\
        &&&0&1&&\\
        &&&-1&0&&\\
        &&&&&0&\\
        0&&&&&&\ddots
    \end{pmatrix},
\]
то есть является прямой суммой блоков
\[
    \begin{pmatrix}0&1\\-1&0\end{pmatrix}
\]
и нулевого блока на ядре формы.
\begin{knowledgefix}
Для кососимметрической формы канонический базис обычно называется симплектическим: \((e_1,f_1,\ldots,e_r,f_r,k_1,
\ldots,k_s)\), где \(\beta(e_i,f_i)=1\), \(\beta(f_i,e_i)=-1\), остальные пары между блоками равны нулю, а \(k_j\) порождают ядро.
\end{knowledgefix}
\end{theorem}

\begin{proof}[Идея доказательства из конспекта]
Доказательство идёт индукцией по $n=\dim V$.

Если форма нулевая, то утверждение очевидно: вся матрица формы нулевая.

Если форма не нулевая, то существуют $e_1,e_2\in V$ такие, что
\[
    \beta(e_1,e_2)\ne0.
\]
Нормируя один из векторов, можно добиться
\[
    \beta(e_1,e_2)=1.
\]
Тогда на подпространстве
\[
    U=\langle e_1,e_2\rangle
\]
матрица ограничения формы имеет вид
\[
    \begin{pmatrix}
        0&1\\
        -1&0
    \end{pmatrix},
\]
поэтому $U$ невырождено относительно $\beta$.

Далее берётся ортогональное дополнение $U^\perp$ относительно формы $\beta$:
\[
    V=U\oplus U^\perp.
\]
К $U^\perp$ применяется индукционное предположение.
\begin{knowledgefix}
Для кососимметрической формы \(\beta(y,x)=-\beta(x,y)\), поэтому левое и правое ортогональные дополнения совпадают как множества: условие \(\beta(u,v)=0\) эквивалентно \(\beta(v,u)=0\). Знаки важны в матрице блока, но не меняют само дополнение.
\end{knowledgefix}
\end{proof}

\begin{corollary}
Ранг кососимметрической матрицы над полем характеристики $\ne2$ чётен:
\[
    \rk B=2r.
\]
\begin{knowledgefix}
Следствие: ранг кососимметрической билинейной формы чётен, а её матрица в подходящем базисе имеет блочно-диагональный вид с блоками \(\begin{pmatrix}0&1\\-1&0\end{pmatrix}\) и, возможно, нулевым блоком ядра.
\end{knowledgefix}
\end{corollary}

\subsubsection{Метод Лагранжа: вычислительный пример}


На страницах 23--24 приведён длинный вычислительный пример приведения квадратичной формы к каноническому виду методом Лагранжа.
Смысл примера: последовательно выделять полные квадраты и вводить новые переменные.

В общем виде, если
\[
    q(x_1,\ldots,x_n)=b_{11}x_1^2+2\sum_{j=2}^n b_{1j}x_1x_j+q_1(x_2,\ldots,x_n)
\]
и $b_{11}\ne0$, то первый шаг метода Лагранжа имеет вид
\[
\begin{aligned}
    q(x_1,\ldots,x_n)
    &=b_{11}\left(x_1+\frac{b_{12}}{b_{11}}x_2+\cdots+\frac{b_{1n}}{b_{11}}x_n\right)^2
      +\widetilde q(x_2,\ldots,x_n).
\end{aligned}
\]

Для конкретного примера в конспекте вводятся новые переменные $y_i$ и выполняются преобразования вида
\[
    y_1=x_1+\cdots,
    \qquad
    y_2=x_2+\cdots,
    \qquad
    y_3=x_3+\cdots.
\]
После замены получается диагональный вид квадратичной формы.

\begin{example}[Пример со страниц 23--24]
\begin{knowledgefix}
Общий алгоритм Лагранжа: последовательно выделять полные квадраты и выполнять невырожденные линейные замены. Итоговая диагональная матрица зависит от коэффициентов формы, а знаки диагональных коэффициентов дают инерцию формы.
\end{knowledgefix}
\end{example}

\clearpage

\section{Билет 5. Положительно/отрицательно определённые билинейные и квадратичные формы. Критерий Сильвестра.}

\paragraph{Содержание билета.} Положительно/отрицательно определённые билинейные и квадратичные формы. Критерий Сильвестра.

\begin{planadd}

главный скелет ответа — определение положительной/отрицательной определённости, связь с нормальным видом, затем критерий Сильвестра через главные угловые миноры.

\end{planadd}

\subsubsection{Канонический вид квадратичной формы}


\begin{theorem}[Существование ортогонального базиса]
Пусть $\operatorname{char}\F\ne2$ и $\beta$ --- симметрическая билинейная форма на конечномерном пространстве $V$.
Тогда существует ортогональный базис пространства $V$ относительно формы $\beta$.
\end{theorem}

\begin{proof}
Доказательство ведётся индукцией по $n=\dim V$.

Если $n=1$, утверждение очевидно.

Пусть $n>1$ и утверждение уже доказано для пространств меньшей размерности.
Рассмотрим квадратичную форму
\[
    q(x)=\beta(x,x).
\]
Если $q(x)=0$ для всех $x\in V$, то по формуле поляризации
\[
    \beta(x,y)=\frac{q(x+y)-q(x)-q(y)}{2}=0,
\]
то есть форма нулевая, и любой базис является ортогональным.

Если же существует $e_1\in V$ такой, что
\[
    q(e_1)=\beta(e_1,e_1)\ne0,
\]
то подпространство $\langle e_1\rangle$ невырождено. Тогда
\[
    V=\langle e_1\rangle\oplus \langle e_1\rangle^\perp.
\]
К подпространству $\langle e_1\rangle^\perp$ применяем индукционное предположение и получаем в нём ортогональный базис
\[
    e_2,\ldots,e_n.
\]
Тогда
\[
    e_1,e_2,\ldots,e_n
\]
--- ортогональный базис всего $V$.
\end{proof}

\begin{corollary}
Любую квадратичную форму можно привести к каноническому виду.
\end{corollary}

\subsubsection*{Метод Лагранжа}

Рассмотрим квадратичную форму
\[
    q(x_1,\ldots,x_n)
    =\sum_{i=1}^n b_{ii}x_i^2
    +2\sum_{1\le i<j\le n} b_{ij}x_ix_j.
\]
Если $b_{11}\ne0$, то выделяем полный квадрат:
\[
\begin{aligned}
    q(x_1,\ldots,x_n)
    &=b_{11}\left(x_1+\frac{b_{12}}{b_{11}}x_2+\cdots+\frac{b_{1n}}{b_{11}}x_n\right)^2 \\
    &\quad + q_1(x_2,\ldots,x_n),
\end{aligned}
\]
где $q_1$ --- квадратичная форма от $n-1$ переменных.
Далее применяем тот же процесс к $q_1$.

Если $b_{11}=0$, но есть ненулевой коэффициент при смешанном члене, то выполняют линейную замену переменных, чтобы появился ненулевой квадрат.
\begin{knowledgefix}
Стандартный ход метода Лагранжа: если при \(x_1^2\) коэффициент \(a_{11}\ne0\), выделяют полный квадрат
\[q=a_{11}\left(x_1+\frac{a_{12}}{a_{11}}x_2+\cdots+\frac{a_{1n}}{a_{11}}x_n\right)^2+q_1(x_2,
\ldots,x_n),\]
после чего повторяют процедуру для \(q_1\). Если подходящего квадратного коэффициента нет, сначала делают линейную замену, создающую ненулевой квадрат.
\end{knowledgefix}

\begin{remark}
В результате метода Лагранжа получается новый базис, в котором матрица квадратичной формы диагональна.
\end{remark}

\subsubsection*{Метод Якоби}

Пусть $B$ --- матрица симметрической билинейной формы в базисе $e_1,\ldots,e_n$.
Обозначим угловые миноры:
\[
    \Delta_k=\det B_k,
\]
где $B_k$ --- левый верхний $k\times k$ блок матрицы $B$, и положим
\[
    \Delta_0=1.
\]

\begin{theorem}[Формула Якоби]
Если
\[
    \Delta_k\ne0,
    \qquad k=1,\ldots,n,
\]
то существует ортогональный базис $f_1,\ldots,f_n$, в котором квадратичная форма имеет вид
\[
    q(x_1,\ldots,x_n)
    =\frac{\Delta_1}{\Delta_0}x_1^2
    +\frac{\Delta_2}{\Delta_1}x_2^2
    +\cdots+
    \frac{\Delta_n}{\Delta_{n-1}}x_n^2.
\]
\begin{knowledgefix}
Точная стандартная формулировка — метод Якоби: если ведущие главные миноры \(\Delta_k\ne0\), то квадратичная форма приводится треугольной заменой к виду
\[q=\Delta_1 y_1^2+\frac{\Delta_2}{\Delta_1}y_2^2+\cdots+\frac{\Delta_n}{\Delta_{n-1}}y_n^2.\]
Эта же формула используется для быстрого определения инерции при невырожденных ведущих минорах.
\end{knowledgefix}
\end{theorem}

\begin{proof}[Идея доказательства]
Пусть
\[
    V_k=\langle e_1,\ldots,e_k\rangle.
\]
Условие $\Delta_k\ne0$ означает, что ограничение формы на $V_k$ невырождено.
Тогда можно строить ортогональный базис последовательно:
\[
    V_1\subset V_2\subset\cdots\subset V_n=V,
\]
выбирая
\[
    f_k\in V_k\cap V_{k-1}^{\perp}
\]
так, чтобы
\[
    V_k=V_{k-1}\oplus\langle f_k\rangle.
\]
В этом базисе матрица формы диагональна.
Диагональные элементы определяются из соотношения для определителей:
\[
    \Delta_k=\prod_{i=1}^k \beta(f_i,f_i).
\]
Поэтому
\[
    \beta(f_k,f_k)=\frac{\Delta_k}{\Delta_{k-1}}.
\]
\begin{knowledgefix}
Смысл геометрического доказательства: симметрическая билинейная форма задаётся самосопряжённым оператором, поэтому в ортонормированном базисе из собственных векторов её матрица диагональна. Геометрически это переход к главным осям.
\end{knowledgefix}
\end{proof}

\subsubsection*{Положительная определённость и критерий Сильвестра}

\begin{definition}[Положительно определённая квадратичная форма]
Квадратичная форма $q$ называется \emph{положительно определённой}, если
\[
    q(x)>0
    \qquad \forall x\ne0.
\]
\end{definition}

\begin{definition}[Отрицательно определённая квадратичная форма]
Квадратичная форма $q$ называется \emph{отрицательно определённой}, если
\[
    q(x)<0
    \qquad \forall x\ne0.
\]
\end{definition}

\begin{definition}[Полуопределённые формы]
Форма называется положительно полуопределённой, если
\[
    q(x)\ge0
    \qquad \forall x\in V,
\]
и отрицательно полуопределённой, если
\[
    q(x)\le0
    \qquad \forall x\in V.
\]
\begin{knowledgefix}
Классификация по знаку: положительно определённая \(q(x)>0\) при \(x\ne0\); отрицательно определённая \(q(x)<0\); положительно полуопределённая \(q(x)\ge0\); отрицательно полуопределённая \(q(x)\le0\); неопределённая — принимает и положительные, и отрицательные значения.
\end{knowledgefix}
\end{definition}

\begin{theorem}[Критерий через канонический вид]
Пусть квадратичная форма приведена к каноническому виду
\[
    q(x_1,
    \ldots,x_n)=a_1x_1^2+
    \cdots+a_nx_n^2.
\]
Тогда:
\begin{enumerate}[label=\arabic*)]
    \item $q$ положительно определена тогда и только тогда, когда $a_i>0$ для всех $i$;
    \item $q$ отрицательно определена тогда и только тогда, когда $a_i<0$ для всех $i$;
    \item если среди $a_i$ есть положительные и отрицательные, то форма неопределённая.
\end{enumerate}
\end{theorem}

\begin{theorem}[Критерий Сильвестра]
Пусть $q$ --- вещественная квадратичная форма с симметрической матрицей $B$, а
\[
    \Delta_k
\]
--- её угловые миноры. Тогда $q$ положительно определена тогда и только тогда, когда
\[
    \Delta_1>0,
    \quad
    \Delta_2>0,
    \quad
    \ldots,
    \quad
    \Delta_n>0.
\]
\end{theorem}

\begin{proof}
Если все угловые миноры положительны, то по формуле Якоби
\[
    q(x_1,
    \ldots,x_n)=
    \frac{\Delta_1}{\Delta_0}x_1^2+
    \frac{\Delta_2}{\Delta_1}x_2^2+
    \cdots+
    \frac{\Delta_n}{\Delta_{n-1}}x_n^2.
\]
Так как все отношения
\[
    \frac{\Delta_k}{\Delta_{k-1}}
\]
положительны, форма положительно определена.
Обратное направление следует из того, что ограничение положительно определённой формы на любое подпространство также положительно определено, поэтому все соответствующие угловые миноры положительны.
\begin{knowledgefix}
Доказательство закона инерции через переходную матрицу: если две диагональные записи одной формы связаны невырожденной заменой \(x=Cy\), то максимальная размерность подпространства, на котором форма положительно определена, не меняется. Эта максимальная размерность равна числу положительных квадратов; аналогично для отрицательных. Поэтому числа положительных, отрицательных и нулевых коэффициентов не зависят от приведения.
\end{knowledgefix}
\end{proof}

\begin{corollary}[Критерий отрицательной определённости]
Форма $q$ отрицательно определена тогда и только тогда, когда
\[
    (-1)^k\Delta_k>0,
    \qquad k=1,
    \ldots,n.
\]
Иначе говоря,
\[
    \Delta_1<0,
    \quad
    \Delta_2>0,
    \quad
    \Delta_3<0,
    \quad
    \ldots
\]
\begin{knowledgefix}
Для отрицательной определённости критерий записывается как \(\Delta_1<0,\Delta_2>0,\Delta_3<0,\ldots\), то есть \((-1)^k\Delta_k>0\). Для положительной определённости: \(\Delta_k>0\) для всех \(k\).
\end{knowledgefix}
\end{corollary}

\begin{example}
На странице 19 приведён пример квадратичной формы от двух переменных и вычисление её матрицы/миноров.
\begin{knowledgefix}
Численный пример заменён проверенным стандартным: матрица \(\begin{pmatrix}2&1\\1&2\end{pmatrix}\) положительно определена, так как \(\Delta_1=2>0\), \(\Delta_2=3>0\). Матрица \(\begin{pmatrix}0&1\\1&0\end{pmatrix}\) неопределённа, поскольку \(q(x,y)=2xy\) принимает оба знака.
\end{knowledgefix}
\end{example}

\clearpage

\section{Билет 6. Евклидовы векторные пространства: определение и примеры. Неравенство Коши — Буняковского, неравенство треугольника, теорема Пифагора. Длина и угол между векторами, теорема косинусов.}

\paragraph{Содержание билета.} Евклидовы векторные пространства: определение и примеры. Неравенство Коши — Буняковского, неравенство треугольника, теорема Пифагора. Длина и угол между векторами, теорема косинусов.

\begin{planadd}

это базовый геометрический билет: определение евклидова пространства $\to$ норма $\to$ Коши--Буняковский $\to$ треугольник $\to$ Пифагор и косинусы.

\end{planadd}

\subsubsection{Евклидовы пространства}

\subsubsection*{Определение и примеры}

\begin{definition}[Евклидово пространство]
Вещественное векторное пространство $V$ называется \emph{евклидовым}, если на нём задана положительно определённая симметрическая билинейная форма
\[
    (\cdot,\cdot):V\times V\to\R.
\]
Иначе говоря, для всех $x,y,z\in V$ и $\lambda\in\R$ выполнено:
\begin{enumerate}[label=\arabic*)]
    \item $(x,y)=(y,x)$;
    \item $(\lambda x+z,y)=\lambda(x,y)+(z,y)$;
    \item $(x,x)>0$ при $x\ne0$.
\end{enumerate}
\end{definition}

\begin{example}
В пространстве $\R^n$ стандартное скалярное произведение задаётся формулой
\[
    (x,y)=\sum_{i=1}^n x_i y_i.
\]
Его матрица в стандартном базисе равна единичной матрице.
\end{example}

\begin{example}
В пространстве $C[a,b]$ можно задать скалярное произведение
\[
    (f,g)=\int_a^b f(x)g(x)\,dx.
\]
\begin{knowledgefix}
Для пространства функций стандартный пример скалярного произведения: \(\langle f,g\rangle=\int_a^b f(t)g(t)\,dt\) над \(\R\) и \(\langle f,g\rangle=\int_a^b f(t)\overline{g(t)}\,dt\) над \(\C\), в зависимости от выбранной convention.
\end{knowledgefix}
\end{example}

\begin{example}
На пространстве матриц приведены примеры скалярных произведений вида
\[
    (A,B)=\tr(A^T B),
\]
а также вариант для симметрических матриц.
\begin{knowledgefix}
На пространстве матриц стандартный проверенный пример билинейной формы записывается как
\(\langle A,B\rangle=\tr(A^{T}B)\) над \(\R\) и как \(\tr(A^{*}B)\) над \(\C\). Запись \(\tr(AB)\) тоже билинейна, но не является стандартным положительным скалярным произведением на всех матрицах.
\end{knowledgefix}
\end{example}

\begin{definition}[Длина вектора]
Длина вектора $x\in V$ определяется формулой
\[
    |x|=\sqrt{(x,x)}.
\]
Иногда также пишут
\[
    \|x\|=\sqrt{(x,x)}.
\]
\end{definition}

\subsubsection*{Ортонормированный базис и ортогональное дополнение}

\begin{definition}[Ортонормированная система]
Система векторов $e_1,\ldots,e_n$ называется ортонормированной, если
\[
    (e_i,e_j)=\delta_{ij},
\]
где
\[
    \delta_{ij}=\begin{cases}
        1, & i=j,\\
        0, & i\ne j.
    \end{cases}
\]
\end{definition}

Если $e_1,\ldots,e_n$ --- ортонормированный базис и
\[
    x=\sum_{i=1}^n x_i e_i,
    \qquad
    y=\sum_{i=1}^n y_i e_i,
\]
то
\[
    (x,y)=\sum_{i=1}^n x_i y_i,
    \qquad
    |x|^2=\sum_{i=1}^n x_i^2.
\]
Кроме того,
\[
    x_i=(x,e_i).
\]

\begin{definition}[Ортогональное дополнение]
Пусть $U\le V$. Тогда
\[
    U^\perp=\{x\in V\mid (x,u)=0\quad \forall u\in U\}.
\]
\end{definition}

\begin{theorem}
Для любого подпространства $U$ конечномерного евклидова пространства $V$ выполнено
\[
    V=U\oplus U^\perp.
\]
\end{theorem}

\begin{proof}
Во-первых,
\[
    U\cap U^\perp=0.
\]
Действительно, если $x\in U\cap U^\perp$, то $x\in U$ и $(x,u)=0$ для всех $u\in U$.
В частности, при $u=x$ получаем
\[
    (x,x)=0,
\]
откуда $x=0$.

Во-вторых, если $e_1,\ldots,e_k$ --- базис $U$, то условие $x\in U^\perp$ задаётся системой
\[
    (x,e_i)=0,
    \qquad i=1,\ldots,k.
\]
Эта система даёт
\[
    \dim U^\perp=n-k.
\]
Следовательно,
\[
    \dim U+\dim U^\perp=n=\dim V.
\]
Так как пересечение равно нулю, получаем прямую сумму
\[
    V=U\oplus U^\perp.
\]
\begin{knowledgefix}
Координатная версия ортогональной проекции: если \(U=\langle u_1,
\ldots,u_k\rangle\), ищем \(\pr_U x=\sum_i c_i u_i\). Условия \(x-\pr_U x\perp u_j\) дают систему Грама \(Gc=b\), где \(G_{ij}=(u_i,u_j)\), \(b_j=(x,u_j)\).
\end{knowledgefix}
\end{proof}

\begin{definition}[Ортогональная проекция]
Если
\[
    V=U\oplus U^\perp,
\]
то каждый $x\in V$ единственным образом представляется в виде
\[
    x=u+w,
    \qquad u\in U,
    \quad w\in U^\perp.
\]
Вектор $u$ называется ортогональной проекцией $x$ на $U$ и обозначается
\[
    \pr_U x=u.
\]
Вектор $w=x-\pr_Ux$ является ортогональной составляющей.
\begin{knowledgefix}
Фиксируем запись именно как \(\ort_U(x)\): это ортогональная составляющая в разложении \(x=\pr_U(x)+\ort_U(x)\). Она принадлежит \(U^\perp\), а \(\pr_U(x)\in U\).
\end{knowledgefix}
\end{definition}

\subsubsection{Неравенство Коши--Буняковского и следствия}


\begin{theorem}[Неравенство Коши--Буняковского]
Для любых $x,y\in V$ выполнено
\[
    |(x,y)|\le |x|\,|y|.
\]
Причём равенство достигается тогда и только тогда, когда $x$ и $y$ линейно зависимы.
\end{theorem}

\begin{proof}
Если $y=0$, то утверждение очевидно. Пусть $y\ne0$.
Для любого $t\in\R$ имеем
\[
    0\le |x-ty|^2=(x-ty,x-ty)
    =|x|^2-2t(x,y)+t^2|y|^2.
\]
Это квадратный трёхчлен по $t$, который не принимает отрицательных значений. Следовательно, его дискриминант неположителен:
\[
    4(x,y)^2-4|x|^2|y|^2\le0.
\]
Отсюда
\[
    (x,y)^2\le |x|^2|y|^2,
\]
то есть
\[
    |(x,y)|\le |x|\,|y|.
\]

Равенство возможно тогда и только тогда, когда существует $t$ такой, что
\[
    |x-ty|^2=0,
\]
то есть $x=ty$.
\end{proof}

\begin{corollary}[Неравенство треугольника]
Для любых $x,y\in V$ выполнено
\[
    |x+y|\le |x|+|y|.
\]
\end{corollary}

\begin{proof}
\[
\begin{aligned}
    |x+y|^2
    &=(x+y,x+y)\\
    &=|x|^2+2(x,y)+|y|^2\\
    &\le |x|^2+2|x|\,|y|+|y|^2\\
    &=(|x|+|y|)^2.
\end{aligned}
\]
Так как обе стороны неотрицательны, получаем
\[
    |x+y|\le |x|+|y|.
\]
\end{proof}

\begin{example}
Для $\R^n$ неравенство Коши--Буняковского имеет вид
\[
    \left|\sum_{i=1}^n x_i y_i\right|
    \le
    \sqrt{\sum_{i=1}^n x_i^2}\,
    \sqrt{\sum_{i=1}^n y_i^2}.
\]
\end{example}

\begin{example}
Для непрерывных функций:
\[
    \left|\int_a^b f(x)g(x)\,dx\right|
    \le
    \sqrt{\int_a^b f^2(x)\,dx}\,
    \sqrt{\int_a^b g^2(x)\,dx}.
\]
\begin{knowledgefix}
Для \(C[a,b]\) можно использовать скалярное произведение \(\langle f,g\rangle=\int_a^b f(t)g(t)\,dt\). Тогда неравенство Коши--Буняковского принимает вид \(|\int_a^b f(t)g(t)dt|\le \sqrt{\int_a^b f^2}\sqrt{\int_a^b g^2}\).
\end{knowledgefix}
\end{example}

\clearpage

\section{Билет 7. Ортогональность векторов. Ортогональное дополнение к подпространству, его свойства. Ортогональная проекция и ортогональная составляющая. Ортонормированные базисы и ортогональные матрицы.}

\paragraph{Содержание билета.} Ортогональность векторов. Ортогональное дополнение к подпространству, его свойства. Ортогональная проекция и ортогональная составляющая. Ортонормированные базисы и ортогональные матрицы.

\begin{planadd}

не путай два ортогональных дополнения: в билетах 1--5 оно было относительно произвольной билинейной формы, здесь — относительно скалярного произведения.

\end{planadd}
\VisualProjection


\subsubsection{Евклидовы пространства}

\subsubsection*{Определение и примеры}

\begin{definition}[Евклидово пространство]
Вещественное векторное пространство $V$ называется \emph{евклидовым}, если на нём задана положительно определённая симметрическая билинейная форма
\[
    (\cdot,\cdot):V\times V\to\R.
\]
Иначе говоря, для всех $x,y,z\in V$ и $\lambda\in\R$ выполнено:
\begin{enumerate}[label=\arabic*)]
    \item $(x,y)=(y,x)$;
    \item $(\lambda x+z,y)=\lambda(x,y)+(z,y)$;
    \item $(x,x)>0$ при $x\ne0$.
\end{enumerate}
\end{definition}

\begin{example}
В пространстве $\R^n$ стандартное скалярное произведение задаётся формулой
\[
    (x,y)=\sum_{i=1}^n x_i y_i.
\]
Его матрица в стандартном базисе равна единичной матрице.
\end{example}

\begin{example}
В пространстве $C[a,b]$ можно задать скалярное произведение
\[
    (f,g)=\int_a^b f(x)g(x)\,dx.
\]
\begin{knowledgefix}
Для пространства функций стандартный пример скалярного произведения: \(\langle f,g\rangle=\int_a^b f(t)g(t)\,dt\) над \(\R\) и \(\langle f,g\rangle=\int_a^b f(t)\overline{g(t)}\,dt\) над \(\C\), в зависимости от выбранной convention.
\end{knowledgefix}
\end{example}

\begin{example}
На пространстве матриц приведены примеры скалярных произведений вида
\[
    (A,B)=\tr(A^T B),
\]
а также вариант для симметрических матриц.
\begin{knowledgefix}
На пространстве матриц стандартный проверенный пример билинейной формы записывается как
\(\langle A,B\rangle=\tr(A^{T}B)\) над \(\R\) и как \(\tr(A^{*}B)\) над \(\C\). Запись \(\tr(AB)\) тоже билинейна, но не является стандартным положительным скалярным произведением на всех матрицах.
\end{knowledgefix}
\end{example}

\begin{definition}[Длина вектора]
Длина вектора $x\in V$ определяется формулой
\[
    |x|=\sqrt{(x,x)}.
\]
Иногда также пишут
\[
    \|x\|=\sqrt{(x,x)}.
\]
\end{definition}

\subsubsection*{Ортонормированный базис и ортогональное дополнение}

\begin{definition}[Ортонормированная система]
Система векторов $e_1,\ldots,e_n$ называется ортонормированной, если
\[
    (e_i,e_j)=\delta_{ij},
\]
где
\[
    \delta_{ij}=\begin{cases}
        1, & i=j,\\
        0, & i\ne j.
    \end{cases}
\]
\end{definition}

Если $e_1,\ldots,e_n$ --- ортонормированный базис и
\[
    x=\sum_{i=1}^n x_i e_i,
    \qquad
    y=\sum_{i=1}^n y_i e_i,
\]
то
\[
    (x,y)=\sum_{i=1}^n x_i y_i,
    \qquad
    |x|^2=\sum_{i=1}^n x_i^2.
\]
Кроме того,
\[
    x_i=(x,e_i).
\]

\begin{definition}[Ортогональное дополнение]
Пусть $U\le V$. Тогда
\[
    U^\perp=\{x\in V\mid (x,u)=0\quad \forall u\in U\}.
\]
\end{definition}

\begin{theorem}
Для любого подпространства $U$ конечномерного евклидова пространства $V$ выполнено
\[
    V=U\oplus U^\perp.
\]
\end{theorem}

\begin{proof}
Во-первых,
\[
    U\cap U^\perp=0.
\]
Действительно, если $x\in U\cap U^\perp$, то $x\in U$ и $(x,u)=0$ для всех $u\in U$.
В частности, при $u=x$ получаем
\[
    (x,x)=0,
\]
откуда $x=0$.

Во-вторых, если $e_1,\ldots,e_k$ --- базис $U$, то условие $x\in U^\perp$ задаётся системой
\[
    (x,e_i)=0,
    \qquad i=1,\ldots,k.
\]
Эта система даёт
\[
    \dim U^\perp=n-k.
\]
Следовательно,
\[
    \dim U+\dim U^\perp=n=\dim V.
\]
Так как пересечение равно нулю, получаем прямую сумму
\[
    V=U\oplus U^\perp.
\]
\begin{knowledgefix}
Координатная версия ортогональной проекции: если \(U=\langle u_1,
\ldots,u_k\rangle\), ищем \(\pr_U x=\sum_i c_i u_i\). Условия \(x-\pr_U x\perp u_j\) дают систему Грама \(Gc=b\), где \(G_{ij}=(u_i,u_j)\), \(b_j=(x,u_j)\).
\end{knowledgefix}
\end{proof}

\begin{definition}[Ортогональная проекция]
Если
\[
    V=U\oplus U^\perp,
\]
то каждый $x\in V$ единственным образом представляется в виде
\[
    x=u+w,
    \qquad u\in U,
    \quad w\in U^\perp.
\]
Вектор $u$ называется ортогональной проекцией $x$ на $U$ и обозначается
\[
    \pr_U x=u.
\]
Вектор $w=x-\pr_Ux$ является ортогональной составляющей.
\begin{knowledgefix}
Фиксируем запись именно как \(\ort_U(x)\): это ортогональная составляющая в разложении \(x=\pr_U(x)+\ort_U(x)\). Она принадлежит \(U^\perp\), а \(\pr_U(x)\in U\).
\end{knowledgefix}
\end{definition}

\subsubsection{Ориентация и ортогональные матрицы}


\begin{definition}[Одинаковая ориентация базисов]
Два базиса $e=(e_1,\ldots,e_n)$ и $f=(f_1,\ldots,f_n)$ вещественного пространства задают одинаковую ориентацию, если матрица перехода от одного базиса к другому имеет положительный определитель.
Если определитель отрицателен, то базисы задают противоположные ориентации.
\end{definition}

\begin{remark}
В двумерном случае это соответствует выбору направления обхода: ``по часовой стрелке'' или ``против часовой стрелки''.
На стр. 29 изображены два ориентированных базиса на плоскости.
\end{remark}

\begin{proposition}[Ортогональные матрицы на плоскости]
Если $C\in O(2)$, то
\[
    \det C=\pm1.
\]
Матрицы с $\det C=1$ являются матрицами поворотов вида
\[
    \begin{pmatrix}
        \cos\varphi&-\sin\varphi\\
        \sin\varphi&\cos\varphi
    \end{pmatrix}.
\]
Матрицы с $\det C=-1$ соответствуют отражениям.
\begin{knowledgefix}
Стандартная матрица отражения относительно прямой, образующей угол \(\varphi\) с осью \(Ox\), равна
\[\begin{pmatrix}\cos 2\varphi&\sin 2\varphi\\ \sin 2\varphi&-\cos 2\varphi\end{pmatrix}.\]
Матрица поворота: \(\begin{pmatrix}\cos\varphi&-\sin\varphi\\\sin\varphi&\cos\varphi\end{pmatrix}\).
\end{knowledgefix}
\end{proposition}

\clearpage

\section{Билет 8. Матрица и определитель Грама системы векторов евклидова пространства, их свойства. Процесс ортогонализации Грама — Шмидта, его геометрический смысл. Геометрический смысл определителя. QR-разложение.}

\paragraph{Содержание билета.} Матрица и определитель Грама системы векторов евклидова пространства, их свойства. Процесс ортогонализации Грама — Шмидта, его геометрический смысл. Геометрический смысл определителя. QR-разложение.

\begin{planadd}

отвечай через геометрию: матрица Грама кодирует взаимные скалярные произведения, её определитель равен квадрату объёма параллелепипеда, Грам--Шмидт строит ортонормированный базис, QR — матричная форма этого процесса.

\end{planadd}
\VisualGramSchmidt


\subsubsection{Матрица Грама и определитель Грама}


\begin{definition}[Матрица Грама]
Для векторов $v_1,\ldots,v_k$ евклидова пространства $V$ матрицей Грама называется матрица
\[
    G(v_1,\ldots,v_k)=\bigl((v_i,v_j)\bigr)_{i,j=1}^k.
\]
Её определитель
\[
    \Gamma(v_1,\ldots,v_k)=\det G(v_1,\ldots,v_k)
\]
называется определителем Грама.
\end{definition}

\begin{theorem}[Критерий линейной зависимости через определитель Грама]
Для векторов $v_1,\ldots,v_k$ евклидова пространства выполнено
\[
    \det G(v_1,\ldots,v_k)=0
    \quad\Longleftrightarrow\quad
    v_1,\ldots,v_k \text{ линейно зависимы}.
\]
\end{theorem}

\begin{proof}
Если $v_1,\ldots,v_k$ линейно зависимы, то один из векторов выражается через остальные, поэтому строки и столбцы матрицы Грама линейно зависимы. Следовательно,
\[
    \det G=0.
\]

В обратную сторону рассуждение идёт через квадратичную форму
\[
    \left(\sum_{i=1}^k \lambda_i v_i,
          \sum_{j=1}^k \lambda_j v_j\right)
    =\sum_{i,j=1}^k \lambda_i\lambda_j (v_i,v_j).
\]
Если $\det G=0$, то существует ненулевой столбец $\lambda=(\lambda_1,\ldots,\lambda_k)^T$, для которого
\[
    G\lambda=0.
\]
Тогда
\[
    \left|\sum_{i=1}^k\lambda_i v_i\right|^2
    =\lambda^T G\lambda=0,
\]
а значит
\[
    \sum_{i=1}^k\lambda_i v_i=0.
\]
Следовательно, $v_1,\ldots,v_k$ линейно зависимы.
\begin{knowledgefix}
Грам--Шмидт: \(f_1=e_1\), \(f_k=e_k-\sum_{i=1}^{k-1}\frac{(e_k,f_i)}{(f_i,f_i)}f_i\). Затем \(u_k=f_k/\|f_k\|\). Геометрически на каждом шаге вычитаются проекции на уже построенное подпространство.
\end{knowledgefix}
\end{proof}

\subsubsection{Ортогонализация Грама--Шмидта}


\begin{theorem}[Процесс ортогонализации]
Пусть $e_1,\ldots,e_n$ --- базис евклидова пространства $V$. Тогда из него можно построить ортогональный базис $f_1,\ldots,f_n$ по формулам
\[
    f_1=e_1,
\]
\[
    f_k=e_k-\sum_{i=1}^{k-1}\frac{(e_k,f_i)}{(f_i,f_i)}f_i,
    \qquad k=2,\ldots,n.
\]
Если дополнительно положить
\[
    \varepsilon_k=\frac{f_k}{|f_k|},
\]
то получится ортонормированный базис.
\begin{knowledgefix}
Формула Грама--Шмидта: \(f_k=e_k-\sum\limits_{i=1}^{k-1}\frac{(e_k,f_i)}{(f_i,f_i)}f_i\). Именно она обеспечивает \((f_k,f_j)=0\) для всех \(j<k\).
\end{knowledgefix}
\end{theorem}

\begin{proof}[Идея]
Пусть
\[
    V_k=\langle e_1,\ldots,e_k\rangle.
\]
На каждом шаге берём компоненту $e_k$, ортогональную уже построенному подпространству $V_{k-1}$:
\[
    f_k=e_k-\pr_{V_{k-1}}e_k.
\]
Тогда
\[
    f_k\perp f_1,\ldots,f_{k-1},
\]
а также
\[
    V_k=\langle f_1,\ldots,f_k\rangle.
\]
\end{proof}

\subsubsection{Объём, определитель Грама и высота}


\begin{definition}[Объём параллелепипеда]
Пусть $x_1,\ldots,x_k$ --- векторы евклидова пространства. Объём $k$-мерного параллелепипеда, натянутого на эти векторы, обозначается
\[
    V_k(x_1,\ldots,x_k).
\]
Для ортогональных векторов
\[
    V_k(x_1,\ldots,x_k)=|x_1|\cdots |x_k|.
\]
В общем случае объём определяется через ортогонализацию:
\[
    V_k(x_1,\ldots,x_k)=|f_1|\cdots |f_k|,
\]
где $f_1,\ldots,f_k$ получены из $x_1,\ldots,x_k$ процессом Грама--Шмидта.
\end{definition}

\begin{theorem}[Определитель Грама и объём]
Для любых $x_1,\ldots,x_k\in V$ выполнено
\[
    \det G(x_1,\ldots,x_k)=V_k(x_1,\ldots,x_k)^2.
\]
\end{theorem}

\begin{proof}
Пусть $f_1,\ldots,f_k$ --- ортогонализированная система, полученная из $x_1,\ldots,x_k$.
Переход от $x_i$ к $f_i$ задаётся верхнетреугольной матрицей с единицами на диагонали, поэтому определитель Грама при таком переходе не меняется.
В ортогональной системе матрица Грама диагональна:
\[
    G(f_1,\ldots,f_k)=\diag(|f_1|^2,\ldots,|f_k|^2).
\]
Значит,
\[
    \det G(f_1,\ldots,f_k)=|f_1|^2\cdots |f_k|^2
    =V_k(x_1,\ldots,x_k)^2.
\]
\begin{knowledgefix}
Если столбцы матрицы \(X\) — координаты векторов \(v_1,
\ldots,v_k\) в ортонормированном базисе, то матрица Грама равна \(G=X^TX\). Поэтому \(\det G\ge0\), а при \(k=n\) \(\det G=(\det X)^2\).
\end{knowledgefix}
\end{proof}

\begin{corollary}[Высота]
Если $x_1,\ldots,x_k$ линейно независимы, то высота, опущенная из $x_k$ на подпространство
\[
    \langle x_1,\ldots,x_{k-1}\rangle,
\]
равна
\[
    h=\frac{V_k(x_1,\ldots,x_k)}{V_{k-1}(x_1,\ldots,x_{k-1})}.
\]
В терминах определителей Грама:
\[
    h^2=\frac{\det G(x_1,\ldots,x_k)}{\det G(x_1,\ldots,x_{k-1})}.
\]
\begin{knowledgefix}
Расстояние от вектора \(x\) до подпространства \(U\) равно \(\|\ort_U x\|\). Угол с подпространством задаётся через проекцию: \(\cos\angle(x,U)=\|\pr_Ux\|/\|x\|\) при \(x\ne0\).
\end{knowledgefix}
\end{corollary}

\begin{remark}[Ориентированный объём]
Если $k=n$ и выбран ориентированный ортонормированный базис, то ориентированный объём параллелепипеда равен определителю матрицы координат:
\[
    \Vol(x_1,\ldots,x_n)=\det X.
\]
При этом обычный объём равен
\[
    V_n(x_1,\ldots,x_n)=|\det X|.
\]
\begin{knowledgefix}
Геометрический смысл определителя Грама: \(\sqrt{\det G(v_1,
\ldots,v_k)}\) — \(k\)-мерный объём параллелепипеда, натянутого на эти векторы. Для ориентированного объёма в полной размерности используется \(\det X\), для неориентированного — \(|\det X|\).
\end{knowledgefix}
\end{remark}

\subsubsection{QR-разложение}


\begin{theorem}[QR-разложение]
Пусть $A\in M_{m\times n}(\R)$ имеет линейно независимые столбцы
\[
    a_1,\ldots,a_n.
\]
Тогда существует разложение
\[
    A=QR,
\]
где столбцы матрицы $Q$ образуют ортонормированную систему, а $R$ --- верхнетреугольная матрица с положительными диагональными элементами.
\begin{knowledgefix}
QR-разложение: если \(A\in M_{m,n}(\R)\) имеет линейно независимые столбцы, то \(A=QR\), где столбцы \(Q\) ортонормированы, а \(R\) верхнетреугольная с положительной диагональю. Для квадратной невырожденной матрицы \(m=n\), \(Q\) ортогональна.
\end{knowledgefix}
\end{theorem}

\begin{proof}[Алгоритм]
К столбцам $a_1,\ldots,a_n$ применяем процесс ортогонализации Грама--Шмидта.
Получаем ортонормированные векторы
\[
    q_1,\ldots,q_n.
\]
Каждый столбец $a_j$ выражается через $q_1,\ldots,q_j$:
\[
    a_j=r_{1j}q_1+\cdots+r_{jj}q_j.
\]
В матричном виде это даёт
\[
    A=QR.
\]
Диагональные элементы $r_{jj}$ выбираются положительными.
\end{proof}

\begin{proposition}[Единственность QR-разложения]
Если
\[
    A=Q_1R_1=Q_2R_2,
\]
где $Q_1,Q_2$ ортогональны, а $R_1,R_2$ верхнетреугольны с положительной диагональю, то
\[
    Q_1=Q_2,
    \qquad
    R_1=R_2.
\]
\begin{knowledgefix}
Единственность QR при положительной диагонали \(R\): если \(A=Q_1R_1=Q_2R_2\), то \(Q_2^TQ_1=R_2R_1^{-1}\). Левая часть ортогональна, правая верхнетреугольна с положительной диагональю; значит обе равны \(I\), и \(Q_1=Q_2, R_1=R_2\).
\end{knowledgefix}
\end{proposition}

\clearpage

\section{Билет 9. Метрика (расстояние), её свойства. Расстояние и угол между вектором и подпространством.}

\paragraph{Содержание билета.} Метрика (расстояние), её свойства. Расстояние и угол между вектором и подпространством.

\begin{planadd}

здесь лучше показать, что расстояние — это метрика из нормы, а расстояние до подпространства выражается через ортогональную составляющую.

\end{planadd}

\subsubsection{Метрические пространства, расстояние и угол}


\begin{definition}[Метрическое пространство]
Пусть $X$ --- множество. Функция
\[
    \rho:X\times X\to\R
\]
называется метрикой, если для любых $x,y,z\in X$ выполнено:
\begin{enumerate}[label=\arabic*)]
    \item $\rho(x,y)\ge0$, причём $\rho(x,y)=0\Longleftrightarrow x=y$;
    \item $\rho(x,y)=\rho(y,x)$;
    \item $\rho(x,z)\le \rho(x,y)+\rho(y,z)$.
\end{enumerate}
\end{definition}

В евклидовом пространстве метрика задаётся формулой
\[
    \rho(x,y)=|x-y|.
\]

\begin{definition}[Расстояние от точки до подпространства]
Пусть $U\le V$. Расстояние от $x\in V$ до $U$ определяется как
\[
    \rho(x,U)=\inf_{u\in U}|x-u|.
\]
\end{definition}

\begin{theorem}[Расстояние до подпространства через ортогональную проекцию]
Если $V$ --- евклидово пространство и $U\le V$, то
\[
    \rho(x,U)=|x-\pr_U x|.
\]
Иначе говоря,
\[
    \rho(x,U)=|\ort_U x|.
\]
\end{theorem}

\begin{proof}
Разложим
\[
    x=\pr_U x+\ort_U x,
    \qquad \pr_Ux\in U,
    \quad \ort_Ux\in U^\perp.
\]
Для любого $u\in U$ имеем
\[
    x-u=(\pr_Ux-u)+\ort_Ux,
\]
причём слагаемые ортогональны. Поэтому
\[
    |x-u|^2=|\pr_Ux-u|^2+|\ort_Ux|^2\ge |\ort_Ux|^2.
\]
Минимум достигается при $u=\pr_Ux$, значит
\[
    \rho(x,U)=|\ort_Ux|.
\]
\begin{knowledgefix}
Фиксируем обозначения: \(\pr_U x\) — ортогональная проекция на \(U\), \(\ort_U x=x-\pr_U x\) — ортогональная составляющая. Тогда \(x=\pr_U x+\ort_U x\), где \(\pr_U x\in U\), \(\ort_U x\in U^\perp\).
\end{knowledgefix}
\end{proof}

\begin{definition}[Угол между векторами]
Для ненулевых $x,y\in V$ угол между ними определяется формулой
\[
    \cos\angle(x,y)=\frac{(x,y)}{|x|\,|y|}.
\]
Корректность определения следует из неравенства Коши--Буняковского.
\end{definition}

\begin{definition}[Угол между вектором и подпространством]
Угол между ненулевым вектором $x$ и подпространством $U$ определяется как минимальный угол между $x$ и ненулевыми векторами из $U$.
В конспекте он выражается через ортогональную проекцию:
\[
    \cos\angle(x,U)=\frac{|\pr_Ux|}{|x|},
    \qquad
    \sin\angle(x,U)=\frac{|\ort_Ux|}{|x|}.
\]
\begin{knowledgefix}
Угол между ненулевым вектором \(x\) и подпространством \(U\) удобно определять как угол между \(x\) и его проекцией \(\pr_Ux\): \(\cos\varphi=\|\pr_Ux\|/\|x\|\), если \(\pr_Ux\ne0\).
\end{knowledgefix}
\end{definition}

\clearpage

\section{Билет 10. Изоморфизм евклидовых пространств.}

\paragraph{Содержание билета.} Изоморфизм евклидовых пространств.

\begin{planadd}

в конспекте этот билет дан скорее через свойства ортонормированных базисов и изометрий; стандартный вывод: евклидовы пространства конечной размерности изоморфны тогда и только тогда, когда имеют одинаковую размерность.

\end{planadd}

\subsubsection{Евклидовы пространства}

\subsubsection*{Определение и примеры}

\begin{definition}[Евклидово пространство]
Вещественное векторное пространство $V$ называется \emph{евклидовым}, если на нём задана положительно определённая симметрическая билинейная форма
\[
    (\cdot,\cdot):V\times V\to\R.
\]
Иначе говоря, для всех $x,y,z\in V$ и $\lambda\in\R$ выполнено:
\begin{enumerate}[label=\arabic*)]
    \item $(x,y)=(y,x)$;
    \item $(\lambda x+z,y)=\lambda(x,y)+(z,y)$;
    \item $(x,x)>0$ при $x\ne0$.
\end{enumerate}
\end{definition}

\begin{example}
В пространстве $\R^n$ стандартное скалярное произведение задаётся формулой
\[
    (x,y)=\sum_{i=1}^n x_i y_i.
\]
Его матрица в стандартном базисе равна единичной матрице.
\end{example}

\begin{example}
В пространстве $C[a,b]$ можно задать скалярное произведение
\[
    (f,g)=\int_a^b f(x)g(x)\,dx.
\]
\begin{knowledgefix}
Для пространства функций стандартный пример скалярного произведения: \(\langle f,g\rangle=\int_a^b f(t)g(t)\,dt\) над \(\R\) и \(\langle f,g\rangle=\int_a^b f(t)\overline{g(t)}\,dt\) над \(\C\), в зависимости от выбранной convention.
\end{knowledgefix}
\end{example}

\begin{example}
На пространстве матриц приведены примеры скалярных произведений вида
\[
    (A,B)=\tr(A^T B),
\]
а также вариант для симметрических матриц.
\begin{knowledgefix}
На пространстве матриц стандартный проверенный пример билинейной формы записывается как
\(\langle A,B\rangle=\tr(A^{T}B)\) над \(\R\) и как \(\tr(A^{*}B)\) над \(\C\). Запись \(\tr(AB)\) тоже билинейна, но не является стандартным положительным скалярным произведением на всех матрицах.
\end{knowledgefix}
\end{example}

\begin{definition}[Длина вектора]
Длина вектора $x\in V$ определяется формулой
\[
    |x|=\sqrt{(x,x)}.
\]
Иногда также пишут
\[
    \|x\|=\sqrt{(x,x)}.
\]
\end{definition}

\subsubsection*{Ортонормированный базис и ортогональное дополнение}

\begin{definition}[Ортонормированная система]
Система векторов $e_1,\ldots,e_n$ называется ортонормированной, если
\[
    (e_i,e_j)=\delta_{ij},
\]
где
\[
    \delta_{ij}=\begin{cases}
        1, & i=j,\\
        0, & i\ne j.
    \end{cases}
\]
\end{definition}

Если $e_1,\ldots,e_n$ --- ортонормированный базис и
\[
    x=\sum_{i=1}^n x_i e_i,
    \qquad
    y=\sum_{i=1}^n y_i e_i,
\]
то
\[
    (x,y)=\sum_{i=1}^n x_i y_i,
    \qquad
    |x|^2=\sum_{i=1}^n x_i^2.
\]
Кроме того,
\[
    x_i=(x,e_i).
\]

\begin{definition}[Ортогональное дополнение]
Пусть $U\le V$. Тогда
\[
    U^\perp=\{x\in V\mid (x,u)=0\quad \forall u\in U\}.
\]
\end{definition}

\begin{theorem}
Для любого подпространства $U$ конечномерного евклидова пространства $V$ выполнено
\[
    V=U\oplus U^\perp.
\]
\end{theorem}

\begin{proof}
Во-первых,
\[
    U\cap U^\perp=0.
\]
Действительно, если $x\in U\cap U^\perp$, то $x\in U$ и $(x,u)=0$ для всех $u\in U$.
В частности, при $u=x$ получаем
\[
    (x,x)=0,
\]
откуда $x=0$.

Во-вторых, если $e_1,\ldots,e_k$ --- базис $U$, то условие $x\in U^\perp$ задаётся системой
\[
    (x,e_i)=0,
    \qquad i=1,\ldots,k.
\]
Эта система даёт
\[
    \dim U^\perp=n-k.
\]
Следовательно,
\[
    \dim U+\dim U^\perp=n=\dim V.
\]
Так как пересечение равно нулю, получаем прямую сумму
\[
    V=U\oplus U^\perp.
\]
\begin{knowledgefix}
Координатная версия ортогональной проекции: если \(U=\langle u_1,
\ldots,u_k\rangle\), ищем \(\pr_U x=\sum_i c_i u_i\). Условия \(x-\pr_U x\perp u_j\) дают систему Грама \(Gc=b\), где \(G_{ij}=(u_i,u_j)\), \(b_j=(x,u_j)\).
\end{knowledgefix}
\end{proof}

\begin{definition}[Ортогональная проекция]
Если
\[
    V=U\oplus U^\perp,
\]
то каждый $x\in V$ единственным образом представляется в виде
\[
    x=u+w,
    \qquad u\in U,
    \quad w\in U^\perp.
\]
Вектор $u$ называется ортогональной проекцией $x$ на $U$ и обозначается
\[
    \pr_U x=u.
\]
Вектор $w=x-\pr_Ux$ является ортогональной составляющей.
\begin{knowledgefix}
Фиксируем запись именно как \(\ort_U(x)\): это ортогональная составляющая в разложении \(x=\pr_U(x)+\ort_U(x)\). Она принадлежит \(U^\perp\), а \(\pr_U(x)\in U\).
\end{knowledgefix}
\end{definition}

\subsubsection{Ортогональные и унитарные операторы}

\subsubsection*{Определения и примеры}

\begin{definition}[Ортогональный и унитарный оператор]
Оператор $A$ называется ортогональным в евклидовом пространстве или унитарным в унитарном пространстве, если он сохраняет скалярное произведение:
\[
    (Ax,Ay)=(x,y)
    \qquad \forall x,y\in V.
\]
Эквивалентно,
\[
    A^*A=E.
\]
Если $A$ обратим, то это равносильно
\[
    A^{-1}=A^*.
\]
\end{definition}

\begin{proposition}[Эквивалентные условия]
Для оператора $A$ в евклидовом или унитарном пространстве эквивалентны:
\begin{enumerate}[label=\arabic*)]
    \item $A$ ортогонален/унитарен;
    \item $A$ сохраняет скалярное произведение;
    \item $A$ сохраняет длины: $\|Ax\|=\|x\|$ для всех $x$;
    \item матрица $A$ в ортонормированном базисе ортогональна/унитарна.
\end{enumerate}
\begin{knowledgefix}
Для ортогонального оператора эквивалентны условия: \(A^*A=I\), \((Ax,Ay)=(x,y)\) для всех \(x,y\), \(\|Ax\|=\|x\|\) для всех \(x\), матрица в ортонормированном базисе удовлетворяет \(Q^TQ=I\), а ортонормированный базис переводится в ортонормированный.
\end{knowledgefix}
\end{proposition}

\begin{example}[Повороты и отражения на плоскости]
В $\R^2$ ортогональные операторы задаются матрицами поворотов
\[
    \begin{pmatrix}
        \cos\varphi & -\sin\varphi\\
        \sin\varphi & \cos\varphi
    \end{pmatrix}
\]
и отражений вида
\[
    \begin{pmatrix}
        \cos\varphi & \sin\varphi\\
        \sin\varphi & -\cos\varphi
    \end{pmatrix}.
\]
У поворотов определитель равен $1$, у отражений --- $-1$.
\begin{knowledgefix}
Канонические двумерные блоки ортогонального оператора: поворот \(R_\varphi=\begin{pmatrix}\cos\varphi&-\sin\varphi\\\sin\varphi&\cos\varphi\end{pmatrix}\). Отражательный блок можно записать как \(\begin{pmatrix}\cos\varphi&\sin\varphi\\\sin\varphi&-\cos\varphi\end{pmatrix}\) после выбора осей.
\end{knowledgefix}
\end{example}

\begin{theorem}[Теорема Эйлера о повороте в трёхмерном пространстве]
В трёхмерном ориентированном евклидовом пространстве всякий ортогональный оператор с определителем $1$ имеет собственную ось, то есть существует ненулевой вектор $e$ такой, что
\[
    Ae=e.
\]
\begin{knowledgefix}
Теорема Эйлера в \(\R^3\): всякое ортогональное преобразование, сохраняющее ориентацию, имеет собственный вектор с собственным значением \(1\). Поэтому оно является вращением вокруг некоторой оси.
\end{knowledgefix}
\end{theorem}
\begin{planadd}
\textbf{Стандартное следствие из материала об ортонормированных базисах.}
Пусть $V$ и $W$ — конечномерные евклидовы пространства. Линейное отображение
$f:V\to W$ называется изоморфизмом евклидовых пространств, если оно биективно и
сохраняет скалярное произведение:
\[
    (f(x),f(y))_W=(x,y)_V \quad \forall x,y\in V.
\]
Если $\dim V=\dim W=n$, выбираем ортонормированные базисы $e_1,\ldots,e_n$ в $V$ и
$u_1,\ldots,u_n$ в $W$, задаём $f(e_i)=u_i$ и продолжаем линейно. Тогда $f$ сохраняет
скалярное произведение. Обратно, изоморфизм евклидовых пространств является линейной
биекцией, поэтому размеры совпадают. Значит, конечномерные евклидовы пространства
изоморфны тогда и только тогда, когда имеют одинаковую размерность.
\end{planadd}

\clearpage

\section{Билет 11. Отождествление $V$ и $V^*$ при наличии невырожденной билинейной формы (например, скалярного произведения). Соответствие между операторами и билинейными формами в евклидовом пространстве. Сопряжённый оператор и его свойства.}

\paragraph{Содержание билета.} Отождествление $V$ и $V^*$ при наличии невырожденной билинейной формы (например, скалярного произведения). Соответствие между операторами и билинейными формами в евклидовом пространстве. Сопряжённый оператор и его свойства.

\begin{planadd}

центральная мысль — невырожденная билинейная форма даёт изоморфизм $V\simeq V^*$; скалярное произведение переводит билинейные формы в операторы; сопряжённый оператор определяется переносом действия с одного аргумента скалярного произведения на другой.

\end{planadd}

\subsubsection{Связь билинейных форм, двойственного пространства и операторов}

\subsubsection*{Невырожденная форма и двойственное пространство}

Пусть $V$ --- конечномерное пространство и $\beta\in\BL(V)$.
Форма $\beta$ задаёт отображение
\[
    V\to V^*,
    \qquad
    y\mapsto \beta(\,\cdot\,,y),
\]
где функционал $\beta(\,\cdot\,,y)$ действует по правилу
\[
    x\mapsto \beta(x,y).
\]
\begin{knowledgefix}
Для билинейной формы есть два естественных отображения: \(x\mapsto\beta(x,\cdot)\) и \(y\mapsto\beta(\cdot,y)\). В евклидовом пространстве оператор обычно задают условием \(\beta(x,y)=(Ax,y)\) или \((x,Ay)\); выбранное соглашение надо держать единым.
\end{knowledgefix}

\begin{proposition}
Если $\beta$ невырождена, то соответствующее отображение
\[
    V\to V^*
\]
является изоморфизмом.
\end{proposition}

\begin{proof}
Ядро этого отображения состоит из таких $y\in V$, что
\[
    \beta(x,y)=0
    \qquad \forall x\in V.
\]
То есть ядро совпадает с правым ядром формы. Если форма невырождена, то ядро равно нулю.
Так как
\[
    \dim V=\dim V^*,
\]
инъективное линейное отображение $V\to V^*$ является изоморфизмом.
\end{proof}

\begin{corollary}[Евклидово пространство отождествляется с двойственным]
В конечномерном евклидовом пространстве каждый линейный функционал $\omega\in V^*$ единственным образом представляется в виде
\[
    \omega(x)=(y,x)
\]
для некоторого $y\in V$.
\begin{knowledgefix}
В вещественном евклидовом случае \((x,y)=(y,x)\), поэтому порядок аргументов не влияет. Для эрмитова случая в этом конспекте фиксируем convention: линейность по первому аргументу и полулинейность по второму, \((y,x)=\overline{(x,y)}\).
\end{knowledgefix}
\end{corollary}

Если $e_1,\ldots,e_n$ --- ортонормированный базис, то двойственный базис $e_1^*,\ldots,e_n^*$ связан с ним так:
\[
    e_i^*(e_j)=\delta_{ij}=(e_i,e_j).
\]
Следовательно, в евклидовом пространстве ортонормированный базис можно отождествлять со своим двойственным базисом.

\subsubsection*{Билинейные формы и линейные операторы}

Пусть $V$ --- конечномерное евклидово пространство.
Каждому линейному оператору $A\in\End(V)$ можно сопоставить билинейную форму
\[
    \alpha_A(x,y)=(x,Ay).
\]
Это задаёт отображение
\[
    \End(V)\to\BL(V),
    \qquad
    A\mapsto \alpha_A.
\]

\begin{proposition}
Отображение
\[
    A\mapsto \alpha_A,
    \qquad
    \alpha_A(x,y)=(x,Ay),
\]
является изоморфизмом векторных пространств
\[
    \End(V)\cong\BL(V).
\]
\end{proposition}

\begin{proof}
Линейность по $A$ очевидна:
\[
    \alpha_{A+B}(x,y)=\alpha_A(x,y)+\alpha_B(x,y),
\]
\[
    \alpha_{\lambda A}(x,y)=\lambda\alpha_A(x,y).
\]

Покажем инъективность. Если $\alpha_A=0$, то
\[
    (x,Ay)=0
    \qquad \forall x,y\in V.
\]
Для каждого фиксированного $y$ вектор $Ay$ ортогонален всем $x\in V$.
Значит,
\[
    Ay=0
    \qquad \forall y\in V,
\]
то есть
\[
    A=0.
\]

Так как
\[
    \dim\End(V)=n^2,
    \qquad
    \dim\BL(V)=n^2,
\]
инъективное линейное отображение между пространствами одинаковой размерности является изоморфизмом.
\end{proof}

В ортонормированном базисе связь матрицы формы и матрицы оператора записывается особенно просто.
Если
\[
    A e_j=\sum_i a_{ij}e_i,
\]
то
\[
    \alpha_A(e_i,e_j)=(e_i,Ae_j)=a_{ij}.
\]
Следовательно, в ортонормированном базисе матрица билинейной формы $\alpha_A$ совпадает с матрицей оператора $A$.
\begin{knowledgefix}
Сопряжённый оператор \(A^*\) задаётся равенством \((Ax,y)=(x,A^*y)\) в вещественном случае и при выбранной эрмитовой convention. В ортонормированном базисе его матрица равна \(A^T\) над \(\R\) и \(\overline{A}^{T}\) над \(\C\).
\end{knowledgefix}

\subsubsection{Сопряжённый оператор и инвариантные подпространства}

\subsubsection*{Связь билинейных форм и операторов}

Пусть $V$ --- евклидово или унитарное пространство, а $A\in\End(V)$.
Сопряжённый оператор $A^*$ задаётся условием
\[
    (Ax,y)=(x,A^*y)
    \qquad \forall x,y\in V.
\]
\begin{knowledgefix}
Фиксируем convention для всего файла: \(\langle \lambda x,y\rangle=\lambda\langle x,y\rangle\), \(\langle x,\lambda y\rangle=\overline{\lambda}\langle x,y\rangle\), \(\langle y,x\rangle=\overline{\langle x,y\rangle}\). Тогда \(\langle Ax,y\rangle=\langle x,A^*y\rangle\).
\end{knowledgefix}

В ортонормированном базисе матрица сопряжённого оператора получается из матрицы $A$ транспонированием и, в комплексном случае, комплексным сопряжением:
\[
    [A^*]=[A]^*.
\]
В вещественном евклидовом случае это просто транспонирование:
\[
    [A^*]=[A]^T.
\]

\begin{lemma}[Инвариантность ортогонального дополнения]
Пусть $U\subset V$ --- подпространство. Тогда
\[
    U \text{ инвариантно относительно } A
    \quad\Longleftrightarrow\quad
    U^\perp \text{ инвариантно относительно } A^*.
\]
В частности, если $A=A^*$, то из $A$-инвариантности $U$ следует $A$-инвариантность $U^\perp$.
\end{lemma}

\begin{proof}
Пусть $U$ инвариантно относительно $A$, то есть
\[
    Au\in U \qquad \forall u\in U.
\]
Возьмём $x\in U^\perp$ и $u\in U$. Тогда
\[
    (u,A^*x)=(Au,x)=0,
\]
так как $Au\in U$, а $x\perp U$. Следовательно, $A^*x\in U^\perp$.
Обратное направление доказывается аналогично, применяя рассуждение к $A^*$.
\begin{knowledgefix}
Правильная форма определения сопряжённого: \((Ax,y)=(x,A^*y)\). Эквивалентная запись с переименованными переменными: \((A^*u,x)=(u,Ax)\) в вещественном случае; в эрмитовом случае нужно учитывать комплексное сопряжение по выбранной convention.
\end{knowledgefix}
\end{proof}

\subsection*{Дополнительные уточнения}

\subsubsection*{1. Билинейные формы, сопряжённое пространство и невырожденность}
Пусть $V$ --- конечномерное пространство над полем $\F$, а
\[
    \beta:V\times V\to \F
\]
--- билинейная форма. При фиксации второго аргумента $v\in V$ получаем линейный функционал
\[
    x\mapsto \beta(x,v),
\]
поэтому билинейной форме соответствует линейное отображение
\[
    \Phi_\beta:V\to V^*,\qquad
    \Phi_\beta(v)(x)=\beta(x,v).
\]
Это даёт изоморфизм
\[
    \operatorname{Bil}(V)\simeq \Hom(V,V^*).
\]

В этой конвенции правое ядро формы записывается как
\[
    \Ker_R\beta=\{v\in V\mid \beta(x,v)=0\ \forall x\in V\}=\Ker \Phi_\beta.
\]
Аналогично, если фиксировать первый аргумент, возникает отображение
\[
    \Psi_\beta:V\to V^*,\qquad \Psi_\beta(u)(y)=\beta(u,y),
\]
и левое ядро
\[
    \Ker_L\beta=\{u\in V\mid \beta(u,y)=0\ \forall y\in V\}=\Ker\Psi_\beta.
\]
Для конечномерного пространства следующие условия эквивалентны:
\[
    \beta\text{ невырождена}
    \quad\Longleftrightarrow\quad
    \Phi_\beta\text{ --- изоморфизм}
    \quad\Longleftrightarrow\quad
    \det B\ne 0
    \quad\Longleftrightarrow\quad
    \Ker_L\beta=\Ker_R\beta=0.
\]

Если на $V$ задано скалярное произведение $(\cdot,\cdot)$, то оно задаёт изоморфизм
\[
    V\xrightarrow{\sim} V^*,\qquad v\mapsto (\cdot,v).
\]
В ортонормированном базисе $e_1,\ldots,e_n$ вектор $e_i$ переходит в сопряжённый ковектор $e^i$, поскольку
\[
    e^i(e_j)=\delta^i_j,
    \qquad
    (e_j,e_i)=\delta_{ji}.
\]

В евклидовом пространстве билинейной форме $\beta$ соответствует единственный оператор $A\in\End(V)$ такой, что
\[
    \beta(x,y)=(x,Ay).
\]
В ортонормированном базисе матрица оператора $A$ совпадает с матрицей формы $\beta$.
Сопряжённый оператор задаётся условием
\[
    (Ax,y)=(x,A^*y).
\]
Тогда
\[
    \beta(x,y)=\beta(y,x)\Longleftrightarrow A^*=A,
    \qquad
    \beta(x,y)=-\beta(y,x)\Longleftrightarrow A^*=-A.
\]

\clearpage

\section{Билет 12. Комплексификация евклидова пространства — эрмитово (унитарное) пространство. Эрмитово скалярное произведение, его свойства. Неравенства Коши — Буняковского и треугольника в эрмитовом пространстве. Эрмитово сопряжённый оператор и его свойства.}

\paragraph{Содержание билета.} Комплексификация евклидова пространства — эрмитово (унитарное) пространство. Эрмитово скалярное произведение, его свойства. Неравенства Коши — Буняковского и треугольника в эрмитовом пространстве. Эрмитово сопряжённый оператор и его свойства.

\begin{planadd}

аккуратно проговори convention: какой аргумент линейный, какой сопряжённо-линейный. Это место часто валит знаками и комплексным сопряжением.

\end{planadd}

\subsubsection{Эрмитовы пространства}

\subsubsection*{Эрмитова форма и неравенство Коши--Буняковского}

\begin{definition}[Эрмитова форма]
Пусть $V$ --- комплексное векторное пространство. Отображение
\[
    (\cdot,\cdot):V\times V\to\C
\]
называется эрмитовой формой, если выполнены свойства:
\begin{enumerate}[label=\arabic*)]
    \item $(x+y,z)=(x,z)+(y,z)$;
    \item $(\lambda x,z)=\lambda(x,z)$;
    \item $(x,y)=\overline{(y,x)}$;
    \item $(x,x)>0$ при $x\ne0$.
\end{enumerate}
\begin{knowledgefix}
Convention зафиксирована: эрмитово произведение линейно по первому аргументу и сопряжённо-линейно по второму. Тогда \((x,\alpha y)=\overline{\alpha}(x,y)\), \((\alpha x,y)=\alpha(x,y)\).
\end{knowledgefix}
\end{definition}

Из свойства эрмитовости следует полулинейность по другому аргументу:
\[
    (x,\lambda y)=\overline{\lambda}(x,y).
\]

\begin{definition}[Унитарное пространство]
Комплексное векторное пространство с положительно определённой эрмитовой формой называется \emph{унитарным пространством}.
\end{definition}

\begin{theorem}[Неравенство Коши--Буняковского]
Для любых $x,y$ в евклидовом или унитарном пространстве выполнено
\[
    |(x,y)|^2\le (x,x)(y,y).
\]
Иначе говоря,
\[
    |(x,y)|\le \|x\|\,\|y\|.
\]
\end{theorem}

\begin{proof}
Если $y=0$, то утверждение очевидно. Пусть $y\ne0$.
Рассмотрим вектор
\[
    z=x-\frac{(x,y)}{(y,y)}y.
\]
Тогда
\[
    (z,y)=0.
\]
Так как $(z,z)\ge0$, получаем
\[
\begin{aligned}
    0\le (z,z)
    &=\left(x-\frac{(x,y)}{(y,y)}y,\ x-\frac{(x,y)}{(y,y)}y\right)\\
    &=(x,x)-\frac{|(x,y)|^2}{(y,y)}.
\end{aligned}
\]
Отсюда
\[
    |(x,y)|^2\le (x,x)(y,y).
\]
\begin{knowledgefix}
Коши--Буняковский доказывается через неотрицательность \(\|x-\lambda y\|^2\). При \(y\ne0\) берём \(\lambda=\frac{(x,y)}{(y,y)}\) для выбранной convention и получаем \(|(x,y)|^2\le (x,x)(y,y)\).
\end{knowledgefix}
\end{proof}

\begin{theorem}[Неравенство треугольника]
Для нормы, порождённой скалярным произведением,
\[
    \|x+y\|\le \|x\|+\|y\|.
\]
\end{theorem}

\begin{proof}
Используя неравенство Коши--Буняковского, имеем
\[
\begin{aligned}
    \|x+y\|^2
    &=(x+y,x+y)\\
    &=\|x\|^2+2\operatorname{Re}(x,y)+\|y\|^2\\
    &\le \|x\|^2+2|(x,y)|+\|y\|^2\\
    &\le \|x\|^2+2\|x\|\|y\|+\|y\|^2\\
    &=(\|x\|+\|y\|)^2.
\end{aligned}
\]
\end{proof}

\subsubsection{Сопряжённые, самосопряжённые и нормальные операторы}

\subsubsection*{Основные свойства сопряжения}

\begin{definition}[Сопряжённый оператор]
Пусть $A\in\End(V)$, где $V$ --- евклидово или унитарное пространство.
Оператор $A^*$ называется сопряжённым к $A$, если
\[
    (Ax,y)=(x,A^*y)
    \qquad \forall x,y\in V.
\]
\end{definition}

Для сопряжения выполнены свойства:
\begin{enumerate}[label=\arabic*)]
    \item $(A^*)^*=A$;
    \item $(A+B)^*=A^*+B^*$;
    \item $(\lambda A)^*=\overline{\lambda}A^*$ в комплексном случае и $(\lambda A)^*=\lambda A^*$ в вещественном случае;
    \item $(AB)^*=B^*A^*$.
\end{enumerate}

\begin{proof}[Проверка свойства для произведения]
Для любых $x,y\in V$ имеем
\[
    (ABx,y)=(Bx,A^*y)=(x,B^*A^*y).
\]
Следовательно,
\[
    (AB)^*=B^*A^*.
\]
\end{proof}

\begin{definition}
Оператор $A$ называется:
\begin{enumerate}[label=\arabic*)]
    \item \emph{самосопряжённым}, если $A=A^*$;
    \item \emph{кососамосопряжённым}, если $A^*=-A$;
    \item \emph{нормальным}, если $AA^*=A^*A$;
    \item \emph{унитарным} в унитарном пространстве, если $A^*A=AA^*=E$;
    \item \emph{ортогональным} в евклидовом пространстве, если $A^TA=AA^T=E$.
\end{enumerate}
\begin{knowledgefix}
Стандартный набор примеров нормальных операторов: самосопряжённые \(A^*=A\), кососамосопряжённые \(A^*=-A\), унитарные/ортогональные \(A^*A=I\). Все они нормальны, то есть \(A^*A=AA^*\).
\end{knowledgefix}
\end{definition}

\begin{lemma}
Если $A$ --- нормальный оператор в унитарном пространстве и $Ax=\lambda x$, то
\[
    A^*x=\overline{\lambda}x.
\]
\end{lemma}

\begin{proof}
Пусть $Ax=\lambda x$. Тогда
\[
    (A-\lambda E)x=0.
\]
Так как $A$ нормален, оператор $A-\lambda E$ тоже нормален. Для нормального оператора $B$ выполнено
\[
    \|Bx\|=\|B^*x\|.
\]
Отсюда
\[
    0=\|(A-\lambda E)x\|=\|(A^*-\overline\lambda E)x\|,
\]
значит
\[
    A^*x=\overline\lambda x.
\]
\begin{knowledgefix}
Если \(A\) нормален, то \(A-\lambda I\) также нормален, поскольку \((A-\lambda I)^*(A-\lambda I)=(A^*-\overline\lambda I)(A-\lambda I)\) и множители переставляются благодаря нормальности \(A\).
\end{knowledgefix}
\end{proof}

\begin{corollary}
Собственные векторы нормального оператора, отвечающие различным собственным значениям, ортогональны.
\end{corollary}

\begin{proof}
Пусть
\[
    Ax=\lambda x,
    \qquad
    Ay=\mu y,
    \qquad
    \lambda\ne\mu.
\]
Тогда
\[
    \lambda(x,y)=(Ax,y)=(x,A^*y)=(x,\overline\mu y)=\mu(x,y),
\]
поэтому
\[
    (\lambda-\mu)(x,y)=0.
\]
Так как $\lambda\ne\mu$, получаем $(x,y)=0$.
\begin{knowledgefix}
При выбранной convention собственные значения в формулах с сопряжённым оператором комплексно сопрягаются: если \(Ax=\lambda x\), то в скалярных произведениях появляется \(\overline\lambda\) при переходе через второй аргумент.
\end{knowledgefix}
\end{proof}

\begin{corollary}
Если $A=A^*$, то все собственные значения $A$ вещественны.
\end{corollary}

\begin{proof}
Если $Ax=\lambda x$, $x\ne0$, то
\[
    \lambda(x,x)=(Ax,x)=(x,Ax)=\overline{\lambda}(x,x).
\]
Так как $(x,x)>0$, имеем $\lambda=\overline\lambda$, то есть $\lambda\in\R$.
\end{proof}

\subsection*{Дополнительные уточнения}

\subsubsection*{2. Эрмитовы пространства и convention линейности}
Будем использовать следующее соглашение:
\[
    \text{эрмитово произведение линейно по первому аргументу и полулинейно по второму.}
\]
То есть
\[
    (\alpha x_1+\beta x_2,y)=\alpha(x_1,y)+\beta(x_2,y),
\]
а по второму аргументу
\[
    (x,\alpha y_1+\beta y_2)=\overline{\alpha}(x,y_1)+\overline{\beta}(x,y_2).
\]
Эрмитова симметричность:
\[
    (x,y)=\overline{(y,x)}.
\]
Положительная определённость:
\[
    (x,x)\ge 0,
    \qquad
    (x,x)=0\Longleftrightarrow x=0.
\]

Для доказательства неравенства Коши--Буняковского при $y\ne0$ удобно рассмотреть
\[
    x-\lambda y,
    \qquad
    \lambda=\frac{(x,y)}{(y,y)}.
\]
Из положительности
\[
    (x-\lambda y,x-\lambda y)\ge0
\]
после раскрытия получается
\[
    |(x,y)|^2\le (x,x)(y,y).
\]
Равенство достигается тогда и только тогда, когда $x$ и $y$ линейно зависимы.

Сопряжённый оператор в эрмитовом пространстве при выбранной convention задаётся формулой
\[
    (Ax,y)=(x,A^*y).
\]
Нормальный оператор удовлетворяет
\[
    AA^*=A^*A.
\]
В конечномерном эрмитовом пространстве нормальный оператор диагонализуется в ортонормированном базисе.
Для самосопряжённого оператора собственные значения вещественны, а собственные подпространства, отвечающие различным собственным значениям, ортогональны.

\clearpage

\section{Билет 13. Нормальные операторы в эрмитовом пространстве, их собственные значения и собственные векторы. Теорема о каноническом виде матрицы нормального оператора в эрмитовом пространстве.}

\paragraph{Содержание билета.} Нормальные операторы в эрмитовом пространстве, их собственные значения и собственные векторы. Теорема о каноническом виде матрицы нормального оператора в эрмитовом пространстве.

\begin{planadd}

для нормального оператора ключевое равенство $AA^*=A^*A$; из него идут ортогональность собственных подпространств и диагональный вид в ортонормированном базисе.

\end{planadd}

\subsubsection{Сопряжённые, самосопряжённые и нормальные операторы}

\subsubsection*{Основные свойства сопряжения}

\begin{definition}[Сопряжённый оператор]
Пусть $A\in\End(V)$, где $V$ --- евклидово или унитарное пространство.
Оператор $A^*$ называется сопряжённым к $A$, если
\[
    (Ax,y)=(x,A^*y)
    \qquad \forall x,y\in V.
\]
\end{definition}

Для сопряжения выполнены свойства:
\begin{enumerate}[label=\arabic*)]
    \item $(A^*)^*=A$;
    \item $(A+B)^*=A^*+B^*$;
    \item $(\lambda A)^*=\overline{\lambda}A^*$ в комплексном случае и $(\lambda A)^*=\lambda A^*$ в вещественном случае;
    \item $(AB)^*=B^*A^*$.
\end{enumerate}

\begin{proof}[Проверка свойства для произведения]
Для любых $x,y\in V$ имеем
\[
    (ABx,y)=(Bx,A^*y)=(x,B^*A^*y).
\]
Следовательно,
\[
    (AB)^*=B^*A^*.
\]
\end{proof}

\begin{definition}
Оператор $A$ называется:
\begin{enumerate}[label=\arabic*)]
    \item \emph{самосопряжённым}, если $A=A^*$;
    \item \emph{кососамосопряжённым}, если $A^*=-A$;
    \item \emph{нормальным}, если $AA^*=A^*A$;
    \item \emph{унитарным} в унитарном пространстве, если $A^*A=AA^*=E$;
    \item \emph{ортогональным} в евклидовом пространстве, если $A^TA=AA^T=E$.
\end{enumerate}
\begin{knowledgefix}
Стандартный набор примеров нормальных операторов: самосопряжённые \(A^*=A\), кососамосопряжённые \(A^*=-A\), унитарные/ортогональные \(A^*A=I\). Все они нормальны, то есть \(A^*A=AA^*\).
\end{knowledgefix}
\end{definition}

\begin{lemma}
Если $A$ --- нормальный оператор в унитарном пространстве и $Ax=\lambda x$, то
\[
    A^*x=\overline{\lambda}x.
\]
\end{lemma}

\begin{proof}
Пусть $Ax=\lambda x$. Тогда
\[
    (A-\lambda E)x=0.
\]
Так как $A$ нормален, оператор $A-\lambda E$ тоже нормален. Для нормального оператора $B$ выполнено
\[
    \|Bx\|=\|B^*x\|.
\]
Отсюда
\[
    0=\|(A-\lambda E)x\|=\|(A^*-\overline\lambda E)x\|,
\]
значит
\[
    A^*x=\overline\lambda x.
\]
\begin{knowledgefix}
Если \(A\) нормален, то \(A-\lambda I\) также нормален, поскольку \((A-\lambda I)^*(A-\lambda I)=(A^*-\overline\lambda I)(A-\lambda I)\) и множители переставляются благодаря нормальности \(A\).
\end{knowledgefix}
\end{proof}

\begin{corollary}
Собственные векторы нормального оператора, отвечающие различным собственным значениям, ортогональны.
\end{corollary}

\begin{proof}
Пусть
\[
    Ax=\lambda x,
    \qquad
    Ay=\mu y,
    \qquad
    \lambda\ne\mu.
\]
Тогда
\[
    \lambda(x,y)=(Ax,y)=(x,A^*y)=(x,\overline\mu y)=\mu(x,y),
\]
поэтому
\[
    (\lambda-\mu)(x,y)=0.
\]
Так как $\lambda\ne\mu$, получаем $(x,y)=0$.
\begin{knowledgefix}
При выбранной convention собственные значения в формулах с сопряжённым оператором комплексно сопрягаются: если \(Ax=\lambda x\), то в скалярных произведениях появляется \(\overline\lambda\) при переходе через второй аргумент.
\end{knowledgefix}
\end{proof}

\begin{corollary}
Если $A=A^*$, то все собственные значения $A$ вещественны.
\end{corollary}

\begin{proof}
Если $Ax=\lambda x$, $x\ne0$, то
\[
    \lambda(x,x)=(Ax,x)=(x,Ax)=\overline{\lambda}(x,x).
\]
Так как $(x,x)>0$, имеем $\lambda=\overline\lambda$, то есть $\lambda\in\R$.
\end{proof}

\subsubsection{Спектральная теорема}

\subsubsection*{Нормальные и самосопряжённые операторы}

\begin{theorem}[Спектральная теорема для нормального оператора]
Пусть $V$ --- конечномерное унитарное пространство и $A\in\End(V)$ --- нормальный оператор.
Тогда в $V$ существует ортонормированный базис из собственных векторов оператора $A$.
В этом базисе матрица $A$ диагональна.
\end{theorem}

\begin{proof}[Схема доказательства из конспекта]
Доказательство проводится индукцией по $n=\dim V$.

Так как поле комплексное, у оператора $A$ есть собственный вектор $e_1$:
\[
    Ae_1=\lambda e_1.
\]
Положим
\[
    U=\langle e_1\rangle.
\]
Для нормального оператора из леммы следует, что
\[
    A^*e_1=\overline\lambda e_1,
\]
поэтому $U$ инвариантно и относительно $A$, и относительно $A^*$.
Следовательно, $U^\perp$ инвариантно относительно $A$.

Ограничение $A|_{U^\perp}$ снова нормально. По предположению индукции в $U^\perp$ существует ортонормированный базис из собственных векторов.
Добавляя к нему $e_1$, получаем ортонормированный базис всего пространства $V$.
\begin{knowledgefix}
Если \(U\) инвариантно относительно нормального/самосопряжённого оператора, то при подходящих условиях \(U^\perp\) тоже инвариантно. В ортогональном разложении \(V=U\oplus U^\perp\) матрица оператора становится блочно-диагональной.
\end{knowledgefix}
\end{proof}

\begin{corollary}[Самосопряжённый случай]
Пусть $V$ --- конечномерное евклидово или унитарное пространство, $A=A^*$.
Тогда существует ортонормированный базис из собственных векторов $A$.
В этом базисе матрица $A$ диагональна, а собственные значения вещественны.
\end{corollary}

\begin{remark}
В вещественном евклидовом случае самосопряжённый оператор соответствует симметрической матрице в ортонормированном базисе.
Поэтому всякая вещественная симметрическая матрица ортогонально диагонализуема.
\begin{knowledgefix}
Самосопряжённый оператор в эрмитовом пространстве имеет вещественные собственные значения и ортонормированный базис из собственных векторов. Над вещественным евклидовым пространством это даёт симметрический оператор и диагональную матрицу в ортонормированном базисе.
\end{knowledgefix}
\end{remark}

\clearpage

\section{Билет 14. Самосопряжённые (симметрические), кососимметрические, ортогональные операторы и их комплексификации — эрмитовы, косоэрмитовы, унитарные операторы соответственно. Собственные значения и собственные векторы этих операторов. Канонические виды их матриц.}

\paragraph{Содержание билета.} Самосопряжённые (симметрические), кососимметрические, ортогональные операторы и их комплексификации — эрмитовы, косоэрмитовы, унитарные операторы соответственно. Собственные значения и собственные векторы этих операторов. Канонические виды их матриц.

\begin{planadd}

это сводный билет по классам операторов; удобно дать таблицу: самосопряжённый/эрмитов, кососимметрический/косоэрмитов, ортогональный/унитарный и их спектральные свойства.

\end{planadd}

\subsubsection{Сопряжённые, самосопряжённые и нормальные операторы}

\subsubsection*{Основные свойства сопряжения}

\begin{definition}[Сопряжённый оператор]
Пусть $A\in\End(V)$, где $V$ --- евклидово или унитарное пространство.
Оператор $A^*$ называется сопряжённым к $A$, если
\[
    (Ax,y)=(x,A^*y)
    \qquad \forall x,y\in V.
\]
\end{definition}

Для сопряжения выполнены свойства:
\begin{enumerate}[label=\arabic*)]
    \item $(A^*)^*=A$;
    \item $(A+B)^*=A^*+B^*$;
    \item $(\lambda A)^*=\overline{\lambda}A^*$ в комплексном случае и $(\lambda A)^*=\lambda A^*$ в вещественном случае;
    \item $(AB)^*=B^*A^*$.
\end{enumerate}

\begin{proof}[Проверка свойства для произведения]
Для любых $x,y\in V$ имеем
\[
    (ABx,y)=(Bx,A^*y)=(x,B^*A^*y).
\]
Следовательно,
\[
    (AB)^*=B^*A^*.
\]
\end{proof}

\begin{definition}
Оператор $A$ называется:
\begin{enumerate}[label=\arabic*)]
    \item \emph{самосопряжённым}, если $A=A^*$;
    \item \emph{кососамосопряжённым}, если $A^*=-A$;
    \item \emph{нормальным}, если $AA^*=A^*A$;
    \item \emph{унитарным} в унитарном пространстве, если $A^*A=AA^*=E$;
    \item \emph{ортогональным} в евклидовом пространстве, если $A^TA=AA^T=E$.
\end{enumerate}
\begin{knowledgefix}
Стандартный набор примеров нормальных операторов: самосопряжённые \(A^*=A\), кососамосопряжённые \(A^*=-A\), унитарные/ортогональные \(A^*A=I\). Все они нормальны, то есть \(A^*A=AA^*\).
\end{knowledgefix}
\end{definition}

\begin{lemma}
Если $A$ --- нормальный оператор в унитарном пространстве и $Ax=\lambda x$, то
\[
    A^*x=\overline{\lambda}x.
\]
\end{lemma}

\begin{proof}
Пусть $Ax=\lambda x$. Тогда
\[
    (A-\lambda E)x=0.
\]
Так как $A$ нормален, оператор $A-\lambda E$ тоже нормален. Для нормального оператора $B$ выполнено
\[
    \|Bx\|=\|B^*x\|.
\]
Отсюда
\[
    0=\|(A-\lambda E)x\|=\|(A^*-\overline\lambda E)x\|,
\]
значит
\[
    A^*x=\overline\lambda x.
\]
\begin{knowledgefix}
Если \(A\) нормален, то \(A-\lambda I\) также нормален, поскольку \((A-\lambda I)^*(A-\lambda I)=(A^*-\overline\lambda I)(A-\lambda I)\) и множители переставляются благодаря нормальности \(A\).
\end{knowledgefix}
\end{proof}

\begin{corollary}
Собственные векторы нормального оператора, отвечающие различным собственным значениям, ортогональны.
\end{corollary}

\begin{proof}
Пусть
\[
    Ax=\lambda x,
    \qquad
    Ay=\mu y,
    \qquad
    \lambda\ne\mu.
\]
Тогда
\[
    \lambda(x,y)=(Ax,y)=(x,A^*y)=(x,\overline\mu y)=\mu(x,y),
\]
поэтому
\[
    (\lambda-\mu)(x,y)=0.
\]
Так как $\lambda\ne\mu$, получаем $(x,y)=0$.
\begin{knowledgefix}
При выбранной convention собственные значения в формулах с сопряжённым оператором комплексно сопрягаются: если \(Ax=\lambda x\), то в скалярных произведениях появляется \(\overline\lambda\) при переходе через второй аргумент.
\end{knowledgefix}
\end{proof}

\begin{corollary}
Если $A=A^*$, то все собственные значения $A$ вещественны.
\end{corollary}

\begin{proof}
Если $Ax=\lambda x$, $x\ne0$, то
\[
    \lambda(x,x)=(Ax,x)=(x,Ax)=\overline{\lambda}(x,x).
\]
Так как $(x,x)>0$, имеем $\lambda=\overline\lambda$, то есть $\lambda\in\R$.
\end{proof}

\subsubsection{Спектральная теорема}

\subsubsection*{Нормальные и самосопряжённые операторы}

\begin{theorem}[Спектральная теорема для нормального оператора]
Пусть $V$ --- конечномерное унитарное пространство и $A\in\End(V)$ --- нормальный оператор.
Тогда в $V$ существует ортонормированный базис из собственных векторов оператора $A$.
В этом базисе матрица $A$ диагональна.
\end{theorem}

\begin{proof}[Схема доказательства из конспекта]
Доказательство проводится индукцией по $n=\dim V$.

Так как поле комплексное, у оператора $A$ есть собственный вектор $e_1$:
\[
    Ae_1=\lambda e_1.
\]
Положим
\[
    U=\langle e_1\rangle.
\]
Для нормального оператора из леммы следует, что
\[
    A^*e_1=\overline\lambda e_1,
\]
поэтому $U$ инвариантно и относительно $A$, и относительно $A^*$.
Следовательно, $U^\perp$ инвариантно относительно $A$.

Ограничение $A|_{U^\perp}$ снова нормально. По предположению индукции в $U^\perp$ существует ортонормированный базис из собственных векторов.
Добавляя к нему $e_1$, получаем ортонормированный базис всего пространства $V$.
\begin{knowledgefix}
Если \(U\) инвариантно относительно нормального/самосопряжённого оператора, то при подходящих условиях \(U^\perp\) тоже инвариантно. В ортогональном разложении \(V=U\oplus U^\perp\) матрица оператора становится блочно-диагональной.
\end{knowledgefix}
\end{proof}

\begin{corollary}[Самосопряжённый случай]
Пусть $V$ --- конечномерное евклидово или унитарное пространство, $A=A^*$.
Тогда существует ортонормированный базис из собственных векторов $A$.
В этом базисе матрица $A$ диагональна, а собственные значения вещественны.
\end{corollary}

\begin{remark}
В вещественном евклидовом случае самосопряжённый оператор соответствует симметрической матрице в ортонормированном базисе.
Поэтому всякая вещественная симметрическая матрица ортогонально диагонализуема.
\begin{knowledgefix}
Самосопряжённый оператор в эрмитовом пространстве имеет вещественные собственные значения и ортонормированный базис из собственных векторов. Над вещественным евклидовым пространством это даёт симметрический оператор и диагональную матрицу в ортонормированном базисе.
\end{knowledgefix}
\end{remark}

\subsubsection{Ортогональные и унитарные операторы}

\subsubsection*{Определения и примеры}

\begin{definition}[Ортогональный и унитарный оператор]
Оператор $A$ называется ортогональным в евклидовом пространстве или унитарным в унитарном пространстве, если он сохраняет скалярное произведение:
\[
    (Ax,Ay)=(x,y)
    \qquad \forall x,y\in V.
\]
Эквивалентно,
\[
    A^*A=E.
\]
Если $A$ обратим, то это равносильно
\[
    A^{-1}=A^*.
\]
\end{definition}

\begin{proposition}[Эквивалентные условия]
Для оператора $A$ в евклидовом или унитарном пространстве эквивалентны:
\begin{enumerate}[label=\arabic*)]
    \item $A$ ортогонален/унитарен;
    \item $A$ сохраняет скалярное произведение;
    \item $A$ сохраняет длины: $\|Ax\|=\|x\|$ для всех $x$;
    \item матрица $A$ в ортонормированном базисе ортогональна/унитарна.
\end{enumerate}
\begin{knowledgefix}
Для ортогонального оператора эквивалентны условия: \(A^*A=I\), \((Ax,Ay)=(x,y)\) для всех \(x,y\), \(\|Ax\|=\|x\|\) для всех \(x\), матрица в ортонормированном базисе удовлетворяет \(Q^TQ=I\), а ортонормированный базис переводится в ортонормированный.
\end{knowledgefix}
\end{proposition}

\begin{example}[Повороты и отражения на плоскости]
В $\R^2$ ортогональные операторы задаются матрицами поворотов
\[
    \begin{pmatrix}
        \cos\varphi & -\sin\varphi\\
        \sin\varphi & \cos\varphi
    \end{pmatrix}
\]
и отражений вида
\[
    \begin{pmatrix}
        \cos\varphi & \sin\varphi\\
        \sin\varphi & -\cos\varphi
    \end{pmatrix}.
\]
У поворотов определитель равен $1$, у отражений --- $-1$.
\begin{knowledgefix}
Канонические двумерные блоки ортогонального оператора: поворот \(R_\varphi=\begin{pmatrix}\cos\varphi&-\sin\varphi\\\sin\varphi&\cos\varphi\end{pmatrix}\). Отражательный блок можно записать как \(\begin{pmatrix}\cos\varphi&\sin\varphi\\\sin\varphi&-\cos\varphi\end{pmatrix}\) после выбора осей.
\end{knowledgefix}
\end{example}

\begin{theorem}[Теорема Эйлера о повороте в трёхмерном пространстве]
В трёхмерном ориентированном евклидовом пространстве всякий ортогональный оператор с определителем $1$ имеет собственную ось, то есть существует ненулевой вектор $e$ такой, что
\[
    Ae=e.
\]
\begin{knowledgefix}
Теорема Эйлера в \(\R^3\): всякое ортогональное преобразование, сохраняющее ориентацию, имеет собственный вектор с собственным значением \(1\). Поэтому оно является вращением вокруг некоторой оси.
\end{knowledgefix}
\end{theorem}

\subsubsection{Канонический вид ортогонального оператора}


\begin{theorem}[Канонический вид ортогонального оператора]
Пусть $A$ --- ортогональный оператор конечномерного евклидова пространства $V$.
Тогда в $V$ существует ортонормированный базис, в котором матрица $A$ имеет блочно-диагональный вид с блоками:
\[
    (1),
    \qquad
    (-1),
    \qquad
    \begin{pmatrix}
        \cos\varphi & -\sin\varphi\\
        \sin\varphi & \cos\varphi
    \end{pmatrix}.
\]
\end{theorem}

\begin{proof}[Идея доказательства]
Доказательство строится по индукции по $\dim V$.
Если у $A$ есть вещественный собственный вектор $e$, то из ортогональности следует, что соответствующее собственное значение равно $1$ или $-1$.
Тогда одномерное пространство $\langle e\rangle$ инвариантно, а его ортогональное дополнение также инвариантно.

Если вещественного собственного вектора нет, то рассматривается двумерное инвариантное подпространство, на котором оператор имеет вид поворота.
Далее применяется индукционное предположение к ортогональному дополнению.
\begin{knowledgefix}
Для вещественного кососимметрического оператора комплексные собственные значения чисто мнимые и попарно сопряжённые. Из собственного вектора \(z=x+iy\) получают вещественное двумерное инвариантное подпространство \(\langle x,y\rangle\), на котором матрица имеет блок \(\begin{pmatrix}0&-a\\a&0\end{pmatrix}\).
\end{knowledgefix}
\end{proof}

\subsection*{Дополнительные уточнения}

\subsubsection*{3. Ортогональные операторы, канонический вид и теорема Эйлера}
Оператор $A$ евклидова пространства называется ортогональным, если он сохраняет скалярное произведение:
\[
    (Ax,Ay)=(x,y)\qquad \forall x,y\in V.
\]
Эквивалентные условия:
\[
    A^*A=E,
    \qquad
    AA^*=E,
    \qquad
    A^{-1}=A^*,
\]
а в ортонормированном базисе:
\[
    A^TA=E.
\]
Ортогональный оператор сохраняет длины, углы и ортонормированные системы.

Канонически вещественный ортогональный оператор раскладывается в ортонормированном базисе на блоки вида
\[
    (1),\qquad (-1),\qquad
    R_\varphi=\begin{pmatrix}
        \cos\varphi & -\sin\varphi\\
        \sin\varphi & \cos\varphi
    \end{pmatrix}.
\]
В трёхмерном ориентированном евклидовом пространстве всякий ортогональный оператор с определителем $1$ имеет собственное направление с собственным значением $1$; геометрически это ось вращения. Это и есть содержание теоремы Эйлера о вращении.

\clearpage

\section{Билет 15. Приведение квадратичной формы к главным осям с использованием самосопряжённых операторов.}

\paragraph{Содержание билета.} Приведение квадратичной формы к главным осям с использованием самосопряжённых операторов.

\begin{planadd}

смысл главных осей — выбрать ортонормированный базис из собственных векторов самосопряжённого оператора, соответствующего квадратичной форме.

\end{planadd}

\subsubsection{Канонический вид квадратичной формы}


\begin{theorem}[Существование ортогонального базиса]
Пусть $\operatorname{char}\F\ne2$ и $\beta$ --- симметрическая билинейная форма на конечномерном пространстве $V$.
Тогда существует ортогональный базис пространства $V$ относительно формы $\beta$.
\end{theorem}

\begin{proof}
Доказательство ведётся индукцией по $n=\dim V$.

Если $n=1$, утверждение очевидно.

Пусть $n>1$ и утверждение уже доказано для пространств меньшей размерности.
Рассмотрим квадратичную форму
\[
    q(x)=\beta(x,x).
\]
Если $q(x)=0$ для всех $x\in V$, то по формуле поляризации
\[
    \beta(x,y)=\frac{q(x+y)-q(x)-q(y)}{2}=0,
\]
то есть форма нулевая, и любой базис является ортогональным.

Если же существует $e_1\in V$ такой, что
\[
    q(e_1)=\beta(e_1,e_1)\ne0,
\]
то подпространство $\langle e_1\rangle$ невырождено. Тогда
\[
    V=\langle e_1\rangle\oplus \langle e_1\rangle^\perp.
\]
К подпространству $\langle e_1\rangle^\perp$ применяем индукционное предположение и получаем в нём ортогональный базис
\[
    e_2,\ldots,e_n.
\]
Тогда
\[
    e_1,e_2,\ldots,e_n
\]
--- ортогональный базис всего $V$.
\end{proof}

\begin{corollary}
Любую квадратичную форму можно привести к каноническому виду.
\end{corollary}

\subsubsection*{Метод Лагранжа}

Рассмотрим квадратичную форму
\[
    q(x_1,\ldots,x_n)
    =\sum_{i=1}^n b_{ii}x_i^2
    +2\sum_{1\le i<j\le n} b_{ij}x_ix_j.
\]
Если $b_{11}\ne0$, то выделяем полный квадрат:
\[
\begin{aligned}
    q(x_1,\ldots,x_n)
    &=b_{11}\left(x_1+\frac{b_{12}}{b_{11}}x_2+\cdots+\frac{b_{1n}}{b_{11}}x_n\right)^2 \\
    &\quad + q_1(x_2,\ldots,x_n),
\end{aligned}
\]
где $q_1$ --- квадратичная форма от $n-1$ переменных.
Далее применяем тот же процесс к $q_1$.

Если $b_{11}=0$, но есть ненулевой коэффициент при смешанном члене, то выполняют линейную замену переменных, чтобы появился ненулевой квадрат.
\begin{knowledgefix}
Стандартный ход метода Лагранжа: если при \(x_1^2\) коэффициент \(a_{11}\ne0\), выделяют полный квадрат
\[q=a_{11}\left(x_1+\frac{a_{12}}{a_{11}}x_2+\cdots+\frac{a_{1n}}{a_{11}}x_n\right)^2+q_1(x_2,
\ldots,x_n),\]
после чего повторяют процедуру для \(q_1\). Если подходящего квадратного коэффициента нет, сначала делают линейную замену, создающую ненулевой квадрат.
\end{knowledgefix}

\begin{remark}
В результате метода Лагранжа получается новый базис, в котором матрица квадратичной формы диагональна.
\end{remark}

\subsubsection*{Метод Якоби}

Пусть $B$ --- матрица симметрической билинейной формы в базисе $e_1,\ldots,e_n$.
Обозначим угловые миноры:
\[
    \Delta_k=\det B_k,
\]
где $B_k$ --- левый верхний $k\times k$ блок матрицы $B$, и положим
\[
    \Delta_0=1.
\]

\begin{theorem}[Формула Якоби]
Если
\[
    \Delta_k\ne0,
    \qquad k=1,\ldots,n,
\]
то существует ортогональный базис $f_1,\ldots,f_n$, в котором квадратичная форма имеет вид
\[
    q(x_1,\ldots,x_n)
    =\frac{\Delta_1}{\Delta_0}x_1^2
    +\frac{\Delta_2}{\Delta_1}x_2^2
    +\cdots+
    \frac{\Delta_n}{\Delta_{n-1}}x_n^2.
\]
\begin{knowledgefix}
Точная стандартная формулировка — метод Якоби: если ведущие главные миноры \(\Delta_k\ne0\), то квадратичная форма приводится треугольной заменой к виду
\[q=\Delta_1 y_1^2+\frac{\Delta_2}{\Delta_1}y_2^2+\cdots+\frac{\Delta_n}{\Delta_{n-1}}y_n^2.\]
Эта же формула используется для быстрого определения инерции при невырожденных ведущих минорах.
\end{knowledgefix}
\end{theorem}

\begin{proof}[Идея доказательства]
Пусть
\[
    V_k=\langle e_1,\ldots,e_k\rangle.
\]
Условие $\Delta_k\ne0$ означает, что ограничение формы на $V_k$ невырождено.
Тогда можно строить ортогональный базис последовательно:
\[
    V_1\subset V_2\subset\cdots\subset V_n=V,
\]
выбирая
\[
    f_k\in V_k\cap V_{k-1}^{\perp}
\]
так, чтобы
\[
    V_k=V_{k-1}\oplus\langle f_k\rangle.
\]
В этом базисе матрица формы диагональна.
Диагональные элементы определяются из соотношения для определителей:
\[
    \Delta_k=\prod_{i=1}^k \beta(f_i,f_i).
\]
Поэтому
\[
    \beta(f_k,f_k)=\frac{\Delta_k}{\Delta_{k-1}}.
\]
\begin{knowledgefix}
Смысл геометрического доказательства: симметрическая билинейная форма задаётся самосопряжённым оператором, поэтому в ортонормированном базисе из собственных векторов её матрица диагональна. Геометрически это переход к главным осям.
\end{knowledgefix}
\end{proof}

\subsubsection*{Положительная определённость и критерий Сильвестра}

\begin{definition}[Положительно определённая квадратичная форма]
Квадратичная форма $q$ называется \emph{положительно определённой}, если
\[
    q(x)>0
    \qquad \forall x\ne0.
\]
\end{definition}

\begin{definition}[Отрицательно определённая квадратичная форма]
Квадратичная форма $q$ называется \emph{отрицательно определённой}, если
\[
    q(x)<0
    \qquad \forall x\ne0.
\]
\end{definition}

\begin{definition}[Полуопределённые формы]
Форма называется положительно полуопределённой, если
\[
    q(x)\ge0
    \qquad \forall x\in V,
\]
и отрицательно полуопределённой, если
\[
    q(x)\le0
    \qquad \forall x\in V.
\]
\begin{knowledgefix}
Классификация по знаку: положительно определённая \(q(x)>0\) при \(x\ne0\); отрицательно определённая \(q(x)<0\); положительно полуопределённая \(q(x)\ge0\); отрицательно полуопределённая \(q(x)\le0\); неопределённая — принимает и положительные, и отрицательные значения.
\end{knowledgefix}
\end{definition}

\begin{theorem}[Критерий через канонический вид]
Пусть квадратичная форма приведена к каноническому виду
\[
    q(x_1,
    \ldots,x_n)=a_1x_1^2+
    \cdots+a_nx_n^2.
\]
Тогда:
\begin{enumerate}[label=\arabic*)]
    \item $q$ положительно определена тогда и только тогда, когда $a_i>0$ для всех $i$;
    \item $q$ отрицательно определена тогда и только тогда, когда $a_i<0$ для всех $i$;
    \item если среди $a_i$ есть положительные и отрицательные, то форма неопределённая.
\end{enumerate}
\end{theorem}

\begin{theorem}[Критерий Сильвестра]
Пусть $q$ --- вещественная квадратичная форма с симметрической матрицей $B$, а
\[
    \Delta_k
\]
--- её угловые миноры. Тогда $q$ положительно определена тогда и только тогда, когда
\[
    \Delta_1>0,
    \quad
    \Delta_2>0,
    \quad
    \ldots,
    \quad
    \Delta_n>0.
\]
\end{theorem}

\begin{proof}
Если все угловые миноры положительны, то по формуле Якоби
\[
    q(x_1,
    \ldots,x_n)=
    \frac{\Delta_1}{\Delta_0}x_1^2+
    \frac{\Delta_2}{\Delta_1}x_2^2+
    \cdots+
    \frac{\Delta_n}{\Delta_{n-1}}x_n^2.
\]
Так как все отношения
\[
    \frac{\Delta_k}{\Delta_{k-1}}
\]
положительны, форма положительно определена.
Обратное направление следует из того, что ограничение положительно определённой формы на любое подпространство также положительно определено, поэтому все соответствующие угловые миноры положительны.
\begin{knowledgefix}
Доказательство закона инерции через переходную матрицу: если две диагональные записи одной формы связаны невырожденной заменой \(x=Cy\), то максимальная размерность подпространства, на котором форма положительно определена, не меняется. Эта максимальная размерность равна числу положительных квадратов; аналогично для отрицательных. Поэтому числа положительных, отрицательных и нулевых коэффициентов не зависят от приведения.
\end{knowledgefix}
\end{proof}

\begin{corollary}[Критерий отрицательной определённости]
Форма $q$ отрицательно определена тогда и только тогда, когда
\[
    (-1)^k\Delta_k>0,
    \qquad k=1,
    \ldots,n.
\]
Иначе говоря,
\[
    \Delta_1<0,
    \quad
    \Delta_2>0,
    \quad
    \Delta_3<0,
    \quad
    \ldots
\]
\begin{knowledgefix}
Для отрицательной определённости критерий записывается как \(\Delta_1<0,\Delta_2>0,\Delta_3<0,\ldots\), то есть \((-1)^k\Delta_k>0\). Для положительной определённости: \(\Delta_k>0\) для всех \(k\).
\end{knowledgefix}
\end{corollary}

\begin{example}
На странице 19 приведён пример квадратичной формы от двух переменных и вычисление её матрицы/миноров.
\begin{knowledgefix}
Численный пример заменён проверенным стандартным: матрица \(\begin{pmatrix}2&1\\1&2\end{pmatrix}\) положительно определена, так как \(\Delta_1=2>0\), \(\Delta_2=3>0\). Матрица \(\begin{pmatrix}0&1\\1&0\end{pmatrix}\) неопределённа, поскольку \(q(x,y)=2xy\) принимает оба знака.
\end{knowledgefix}
\end{example}

\subsubsection{Спектральная теорема}

\subsubsection*{Нормальные и самосопряжённые операторы}

\begin{theorem}[Спектральная теорема для нормального оператора]
Пусть $V$ --- конечномерное унитарное пространство и $A\in\End(V)$ --- нормальный оператор.
Тогда в $V$ существует ортонормированный базис из собственных векторов оператора $A$.
В этом базисе матрица $A$ диагональна.
\end{theorem}

\begin{proof}[Схема доказательства из конспекта]
Доказательство проводится индукцией по $n=\dim V$.

Так как поле комплексное, у оператора $A$ есть собственный вектор $e_1$:
\[
    Ae_1=\lambda e_1.
\]
Положим
\[
    U=\langle e_1\rangle.
\]
Для нормального оператора из леммы следует, что
\[
    A^*e_1=\overline\lambda e_1,
\]
поэтому $U$ инвариантно и относительно $A$, и относительно $A^*$.
Следовательно, $U^\perp$ инвариантно относительно $A$.

Ограничение $A|_{U^\perp}$ снова нормально. По предположению индукции в $U^\perp$ существует ортонормированный базис из собственных векторов.
Добавляя к нему $e_1$, получаем ортонормированный базис всего пространства $V$.
\begin{knowledgefix}
Если \(U\) инвариантно относительно нормального/самосопряжённого оператора, то при подходящих условиях \(U^\perp\) тоже инвариантно. В ортогональном разложении \(V=U\oplus U^\perp\) матрица оператора становится блочно-диагональной.
\end{knowledgefix}
\end{proof}

\begin{corollary}[Самосопряжённый случай]
Пусть $V$ --- конечномерное евклидово или унитарное пространство, $A=A^*$.
Тогда существует ортонормированный базис из собственных векторов $A$.
В этом базисе матрица $A$ диагональна, а собственные значения вещественны.
\end{corollary}

\begin{remark}
В вещественном евклидовом случае самосопряжённый оператор соответствует симметрической матрице в ортонормированном базисе.
Поэтому всякая вещественная симметрическая матрица ортогонально диагонализуема.
\begin{knowledgefix}
Самосопряжённый оператор в эрмитовом пространстве имеет вещественные собственные значения и ортонормированный базис из собственных векторов. Над вещественным евклидовым пространством это даёт симметрический оператор и диагональную матрицу в ортонормированном базисе.
\end{knowledgefix}
\end{remark}

\clearpage

\section{Билет 16. Неотрицательно и положительно определённые самосопряжённые операторы, извлечение корней.}

\paragraph{Содержание билета.} Неотрицательно и положительно определённые самосопряжённые операторы, извлечение корней.

\begin{planadd}

держи в голове функциональное исчисление на диагональном виде: если $A=Q\Lambda Q^*$ и $\Lambda\ge0$, то $\sqrt A=Q\sqrt\Lambda Q^*$.

\end{planadd}

\subsubsection{Самосопряжённые операторы и положительные операторы}


\begin{definition}[Положительный оператор]
Самосопряжённый оператор $A=A^*$ называется положительно определённым, если
\[
    (Ax,x)>0
    \qquad \forall x\ne0.
\]
Если вместо строгого неравенства стоит
\[
    (Ax,x)\ge0,
\]
то оператор называется неотрицательным.
\begin{knowledgefix}
Для положительного корня обозначаем исходный положительный самосопряжённый оператор через \(A\), а его корень через \(B\): \(B\ge0\), \(B^*=B\), \(B^2=A\). Если \(A=Q\diag(\lambda_i)Q^*\), то \(B=Q\diag(\sqrt{\lambda_i})Q^*\).
\end{knowledgefix}
\end{definition}

\begin{theorem}[Квадратный корень из положительного оператора]
Пусть $A$ --- положительно определённый самосопряжённый оператор.
Тогда существует единственный положительно определённый самосопряжённый оператор $B$ такой, что
\[
    B^2=A.
\]
\end{theorem}

\begin{proof}
Так как $A$ самосопряжён, существует ортонормированный базис, в котором
\[
    [A]=\diag(\lambda_1,\ldots,\lambda_n),
    \qquad
    \lambda_i>0.
\]
Положим
\[
    [B]=\diag(\sqrt{\lambda_1},\ldots,\sqrt{\lambda_n}).
\]
Тогда $B$ самосопряжён, положительно определён и
\[
    B^2=A.
\]

Единственность следует из того, что любой положительно определённый корень должен иметь в том же собственном базисе собственные значения $\sqrt{\lambda_i}$.
\begin{knowledgefix}
Единственность QR при положительной диагонали \(R\): если \(A=Q_1R_1=Q_2R_2\), то \(Q_2^TQ_1=R_2R_1^{-1}\). Левая часть ортогональна, правая верхнетреугольна с положительной диагональю; значит обе равны \(I\), и \(Q_1=Q_2, R_1=R_2\).
\end{knowledgefix}
\end{proof}

\begin{theorem}[Полярное разложение]
Пусть $A$ --- невырожденный линейный оператор евклидова или унитарного пространства $V$.
Тогда существует единственное представление
\[
    A=UB,
\]
где $U$ --- ортогональный/унитарный оператор, а $B$ --- положительно определённый самосопряжённый оператор.
\begin{knowledgefix}
Полярное разложение для невырожденного \(A\): \(A=UP\), где \(P=(A^*A)^{1/2}\) положительно определён, а \(U=AP^{-1}\) ортогонален/унитарен. Также можно писать \(A=Q U\) с \(Q=(AA^*)^{1/2}\); это другая, левая версия разложения.
\end{knowledgefix}
\end{theorem}

\begin{proof}[Схема]
Рассмотрим оператор
\[
    A^*A.
\]
Он самосопряжён и положительно определён, так как
\[
    (A^*Ax,x)=(Ax,Ax)=\|Ax\|^2>0
\]
при $x\ne0$ и невырожденном $A$.
Пусть
\[
    B=(A^*A)^{1/2}.
\]
Тогда $B$ положительно определён и $B^2=A^*A$.
Положим
\[
    U=AB^{-1}.
\]
Проверим, что $U$ ортогонален/унитарен:
\[
\begin{aligned}
    U^*U
    &=(AB^{-1})^*(AB^{-1})\\
    &=(B^{-1})^*A^*AB^{-1}\\
    &=B^{-1}B^2B^{-1}=E.
\end{aligned}
\]
Следовательно, $A=UB$.

Единственность следует из единственности положительного квадратного корня.
\begin{knowledgefix}
Единственность полярного разложения: из \(A=UP\), где \(P>0\), следует \(A^*A=P U^* U P=P^2\). Значит \(P=(A^*A)^{1/2}\) единственен, а затем \(U=AP^{-1}\) тоже единственен.
\end{knowledgefix}
\end{proof}

\subsection*{Дополнительные уточнения}

\subsubsection*{4. Положительный корень оператора и нормированные пространства}
Самосопряжённый оператор $A$ называется положительно определённым, если
\[
    (Ax,x)>0\qquad \forall x\ne0.
\]
Если $A$ положительно определённый самосопряжённый оператор, то существует единственный положительно определённый самосопряжённый оператор $B$ такой, что
\[
    B^2=A.
\]
Если в ортонормированном базисе
\[
    A=\diag(\lambda_1,\ldots,\lambda_n),\qquad \lambda_i>0,
\]
то
\[
    \sqrt A=\diag(\sqrt{\lambda_1},\ldots,\sqrt{\lambda_n}).
\]
Единственность следует из того, что любой положительный корень сохраняет собственные подпространства $A$ и на каждом таком подпространстве действует как положительный корень соответствующего собственного значения.

Нормированное пространство --- это пространство с функцией $\|\cdot\|$, для которой выполнены:
\[
    \|x\|\ge0,
    \qquad
    \|x\|=0\Longleftrightarrow x=0,
\]
\[
    \|\alpha x\|=|\alpha|\|x\|,
    \qquad
    \|x+y\|\le \|x\|+\|y\|.
\]
Примеры:
\[
    \|x\|_2=\sqrt{x_1^2+\cdots+x_n^2},
    \qquad
    \|x\|_1=|x_1|+\cdots+|x_n|,
    \qquad
    \|x\|_\infty=\max_i |x_i|.
\]
Для нормы, порождённой скалярным произведением, выполнено тождество параллелограмма:
\[
    \|x+y\|^2+\|x-y\|^2=2\|x\|^2+2\|y\|^2.
\]
Теорема Йордана--фон Неймана: вещественная норма порождена скалярным произведением тогда и только тогда, когда выполнено тождество параллелограмма. Скалярное произведение восстанавливается формулой поляризации
\[
    (x,y)=\frac{1}{4}\bigl(\|x+y\|^2-\|x-y\|^2\bigr).
\]

\clearpage

\section{Билет 17. Полярное разложение невырожденного линейного оператора в евклидовом пространстве.}

\paragraph{Содержание билета.} Полярное разложение невырожденного линейного оператора в евклидовом пространстве.

\begin{planadd}

полярное разложение надо объяснять как аналог $z=re^{i\varphi}$: оператор = положительная часть $\times$ ортогональная/унитарная часть.

\end{planadd}

\subsubsection{Самосопряжённые операторы и положительные операторы}


\begin{definition}[Положительный оператор]
Самосопряжённый оператор $A=A^*$ называется положительно определённым, если
\[
    (Ax,x)>0
    \qquad \forall x\ne0.
\]
Если вместо строгого неравенства стоит
\[
    (Ax,x)\ge0,
\]
то оператор называется неотрицательным.
\begin{knowledgefix}
Для положительного корня обозначаем исходный положительный самосопряжённый оператор через \(A\), а его корень через \(B\): \(B\ge0\), \(B^*=B\), \(B^2=A\). Если \(A=Q\diag(\lambda_i)Q^*\), то \(B=Q\diag(\sqrt{\lambda_i})Q^*\).
\end{knowledgefix}
\end{definition}

\begin{theorem}[Квадратный корень из положительного оператора]
Пусть $A$ --- положительно определённый самосопряжённый оператор.
Тогда существует единственный положительно определённый самосопряжённый оператор $B$ такой, что
\[
    B^2=A.
\]
\end{theorem}

\begin{proof}
Так как $A$ самосопряжён, существует ортонормированный базис, в котором
\[
    [A]=\diag(\lambda_1,\ldots,\lambda_n),
    \qquad
    \lambda_i>0.
\]
Положим
\[
    [B]=\diag(\sqrt{\lambda_1},\ldots,\sqrt{\lambda_n}).
\]
Тогда $B$ самосопряжён, положительно определён и
\[
    B^2=A.
\]

Единственность следует из того, что любой положительно определённый корень должен иметь в том же собственном базисе собственные значения $\sqrt{\lambda_i}$.
\begin{knowledgefix}
Единственность QR при положительной диагонали \(R\): если \(A=Q_1R_1=Q_2R_2\), то \(Q_2^TQ_1=R_2R_1^{-1}\). Левая часть ортогональна, правая верхнетреугольна с положительной диагональю; значит обе равны \(I\), и \(Q_1=Q_2, R_1=R_2\).
\end{knowledgefix}
\end{proof}

\begin{theorem}[Полярное разложение]
Пусть $A$ --- невырожденный линейный оператор евклидова или унитарного пространства $V$.
Тогда существует единственное представление
\[
    A=UB,
\]
где $U$ --- ортогональный/унитарный оператор, а $B$ --- положительно определённый самосопряжённый оператор.
\begin{knowledgefix}
Полярное разложение для невырожденного \(A\): \(A=UP\), где \(P=(A^*A)^{1/2}\) положительно определён, а \(U=AP^{-1}\) ортогонален/унитарен. Также можно писать \(A=Q U\) с \(Q=(AA^*)^{1/2}\); это другая, левая версия разложения.
\end{knowledgefix}
\end{theorem}

\begin{proof}[Схема]
Рассмотрим оператор
\[
    A^*A.
\]
Он самосопряжён и положительно определён, так как
\[
    (A^*Ax,x)=(Ax,Ax)=\|Ax\|^2>0
\]
при $x\ne0$ и невырожденном $A$.
Пусть
\[
    B=(A^*A)^{1/2}.
\]
Тогда $B$ положительно определён и $B^2=A^*A$.
Положим
\[
    U=AB^{-1}.
\]
Проверим, что $U$ ортогонален/унитарен:
\[
\begin{aligned}
    U^*U
    &=(AB^{-1})^*(AB^{-1})\\
    &=(B^{-1})^*A^*AB^{-1}\\
    &=B^{-1}B^2B^{-1}=E.
\end{aligned}
\]
Следовательно, $A=UB$.

Единственность следует из единственности положительного квадратного корня.
\begin{knowledgefix}
Единственность полярного разложения: из \(A=UP\), где \(P>0\), следует \(A^*A=P U^* U P=P^2\). Значит \(P=(A^*A)^{1/2}\) единственен, а затем \(U=AP^{-1}\) тоже единственен.
\end{knowledgefix}
\end{proof}

\clearpage

\section{Билет 18. Нормированные пространства. Пространство со скалярным произведением, тождество параллелограмма. Теорема Йордана — фон Неймана.}

\paragraph{Содержание билета.} Нормированные пространства. Пространство со скалярным произведением, тождество параллелограмма. Теорема Йордана — фон Неймана.

\begin{planadd}

билет про отличие произвольной нормы от нормы из скалярного произведения. Критерий — тождество параллелограмма, восстановление — формула поляризации.

\end{planadd}

\subsubsection{Нормированные пространства}


\begin{definition}[Норма]
Пусть $V$ --- векторное пространство над $\R$ или $\C$.
Отображение
\[
    \|\cdot\|:V\to\R
\]
называется нормой, если для всех $x,y\in V$ и $\lambda\in\F$ выполнено:
\begin{enumerate}[label=\arabic*)]
    \item $\|x\|\ge0$, причём $\|x\|=0\Longleftrightarrow x=0$;
    \item $\|\lambda x\|=|\lambda|\,\|x\|$;
    \item $\|x+y\|\le \|x\|+\|y\|$.
\end{enumerate}
\end{definition}

\begin{example}
В евклидовом или унитарном пространстве скалярное произведение задаёт норму
\[
    \|x\|=\sqrt{(x,x)}.
\]
\end{example}

\begin{example}
В $\R^n$ можно рассматривать нормы
\[
    \|x\|_1=|x_1|+\cdots+|x_n|,
\]
\[
    \|x\|_2=\sqrt{|x_1|^2+\cdots+|x_n|^2},
\]
\[
    \|x\|_\infty=\max_i |x_i|.
\]
\begin{knowledgefix}
Примеры норм: \(\|x\|_2=\sqrt{\sum x_i^2}\), \(\|x\|_1=\sum |x_i|\), \(\|x\|_\infty=\max_i|x_i|\); на \(C[a,b]\): \(\|f\|_\infty=\max_{[a,b]}|f|\) и \(\|f\|_2=(\int_a^b |f(t)|^2dt)^{1/2}\).
\end{knowledgefix}
\end{example}

\begin{example}
Для пространства непрерывных функций можно задать нормы вида
\[
    \|f\|_\infty=\max_{x\in[a,b]}|f(x)|,
\]
а также интегральные нормы, например
\[
    \|f\|_1=\int_a^b |f(x)|\,dx.
\]
\begin{knowledgefix}
Для пространства функций стандартный пример скалярного произведения: \(\langle f,g\rangle=\int_a^b f(t)g(t)\,dt\) над \(\R\) и \(\langle f,g\rangle=\int_a^b f(t)\overline{g(t)}\,dt\) над \(\C\), в зависимости от выбранной convention.
\end{knowledgefix}
\end{example}

\begin{lemma}[Тождество параллелограмма]
Если норма порождена скалярным произведением, то для любых $x,y\in V$ выполнено
\[
    \|x+y\|^2+\|x-y\|^2=2\|x\|^2+2\|y\|^2.
\]
\end{lemma}

\begin{proof}
Раскроем обе нормы через скалярное произведение:
\[
\begin{aligned}
    \|x+y\|^2&=(x+y,x+y)=\|x\|^2+2\operatorname{Re}(x,y)+\|y\|^2,\\
    \|x-y\|^2&=(x-y,x-y)=\|x\|^2-2\operatorname{Re}(x,y)+\|y\|^2.
\end{aligned}
\]
Складывая, получаем
\[
    \|x+y\|^2+\|x-y\|^2=2\|x\|^2+2\|y\|^2.
\]
\end{proof}

\begin{theorem}[Теорема Йордана--фон Неймана]
Норма на вещественном векторном пространстве порождена некоторым скалярным произведением тогда и только тогда, когда она удовлетворяет тождеству параллелограмма.
\begin{knowledgefix}
Теорема Йордана--фон Неймана: норма порождена скалярным произведением тогда и только тогда, когда выполнено тождество параллелограмма \(\|x+y\|^2+\|x-y\|^2=2\|x\|^2+2\|y\|^2\).
\end{knowledgefix}
\end{theorem}

В этом случае скалярное произведение восстанавливается по норме формулой поляризации. В вещественном случае:
\[
    (x,y)=\frac{1}{4}\left(\|x+y\|^2-\|x-y\|^2\right).
\]
В комплексном случае формула поляризации имеет вид
\[
    (x,y)=\frac14\left(\|x+y\|^2-\|x-y\|^2+i\|x+iy\|^2-i\|x-iy\|^2\right).
\]
\begin{knowledgefix}
Формула поляризации над \(\R\): \((x,y)=\frac14(\|x+y\|^2-\|x-y\|^2)\). Над \(\C\): \((x,y)=\frac14\sum_{k=0}^{3} i^k\|x+i^k y\|^2\) при соответствующей convention.
\end{knowledgefix}

\subsection*{Дополнительные уточнения}

\subsubsection*{4. Положительный корень оператора и нормированные пространства}
Самосопряжённый оператор $A$ называется положительно определённым, если
\[
    (Ax,x)>0\qquad \forall x\ne0.
\]
Если $A$ положительно определённый самосопряжённый оператор, то существует единственный положительно определённый самосопряжённый оператор $B$ такой, что
\[
    B^2=A.
\]
Если в ортонормированном базисе
\[
    A=\diag(\lambda_1,\ldots,\lambda_n),\qquad \lambda_i>0,
\]
то
\[
    \sqrt A=\diag(\sqrt{\lambda_1},\ldots,\sqrt{\lambda_n}).
\]
Единственность следует из того, что любой положительный корень сохраняет собственные подпространства $A$ и на каждом таком подпространстве действует как положительный корень соответствующего собственного значения.

Нормированное пространство --- это пространство с функцией $\|\cdot\|$, для которой выполнены:
\[
    \|x\|\ge0,
    \qquad
    \|x\|=0\Longleftrightarrow x=0,
\]
\[
    \|\alpha x\|=|\alpha|\|x\|,
    \qquad
    \|x+y\|\le \|x\|+\|y\|.
\]
Примеры:
\[
    \|x\|_2=\sqrt{x_1^2+\cdots+x_n^2},
    \qquad
    \|x\|_1=|x_1|+\cdots+|x_n|,
    \qquad
    \|x\|_\infty=\max_i |x_i|.
\]
Для нормы, порождённой скалярным произведением, выполнено тождество параллелограмма:
\[
    \|x+y\|^2+\|x-y\|^2=2\|x\|^2+2\|y\|^2.
\]
Теорема Йордана--фон Неймана: вещественная норма порождена скалярным произведением тогда и только тогда, когда выполнено тождество параллелограмма. Скалярное произведение восстанавливается формулой поляризации
\[
    (x,y)=\frac{1}{4}\bigl(\|x+y\|^2-\|x-y\|^2\bigr).
\]

\clearpage

\section{Билет 19. Евклидовы аффинные пространства. Декартова прямоугольная система координат. Расстояние между аффинными плоскостями. Расстояние от точки до аффинной плоскости. Расстояние от точки до гиперплоскости.}

\paragraph{Содержание билета.} Евклидовы аффинные пространства. Декартова прямоугольная система координат. Расстояние между аффинными плоскостями. Расстояние от точки до аффинной плоскости. Расстояние от точки до гиперплоскости.

\begin{planadd}

в ответе раздели: аффинное пространство как точки + векторы; декартова прямоугольная система; расстояние между плоскостями через ортогональную составляющую; частные случаи для точки и гиперплоскости.

\end{planadd}

\subsubsection{Переход к аффинным евклидовым пространствам}

\subsubsection*{Конец страницы 60}

В конце страницы 60 начинается новая тема: \emph{аффинные евклидовы пространства}.

\begin{definition}[Черновая запись начала темы]
Аффинное пространство можно понимать как множество точек, на котором соответствующее векторное пространство действует свободно и транзитивно.
\begin{knowledgefix}
Аффинно-евклидово пространство — аффинное пространство \(E\), связанное с евклидовым векторным пространством \(V\). Для точек \(A,B\in E\) расстояние задаётся формулой \(\rho(A,B)=\|\overrightarrow{AB}\|\).
\end{knowledgefix}
\end{definition}


% ============================================================
% Блок страниц 061--080: Аффинные евклидовы пространства, квадрики, тензоры
% ============================================================

\subsubsection{Аффинные евклидовы пространства}

\subsubsection*{Метрика в аффинном евклидовом пространстве}

В конце предыдущего блока была введена аффинная евклидова геометрия: к аффинному пространству присоединяется евклидово векторное пространство направлений. Если $A,B$ --- точки, то вектор от $A$ к $B$ обозначается
\[
    \overrightarrow{AB}\in V.
\]
Расстояние между точками определяется через длину этого вектора:
\[
    \rho(A,B)=\lVert \overrightarrow{AB}\rVert.
\]

\begin{proposition}[Свойства расстояния]
Для любых точек $A,B,C$ выполнены:
\begin{enumerate}[label=\arabic*)]
    \item $\rho(A,B)\ge 0$;
    \item $\rho(A,B)=0 \Longleftrightarrow A=B$;
    \item $\rho(A,B)=\rho(B,A)$;
    \item $\rho(A,C)\le \rho(A,B)+\rho(B,C)$.
\end{enumerate}
\end{proposition}

\begin{proof}
Первые три свойства следуют из свойств нормы и равенства
\[
    \overrightarrow{BA}=-\overrightarrow{AB}.
\]
Для неравенства треугольника используем тождество
\[
    \overrightarrow{AC}=\overrightarrow{AB}+\overrightarrow{BC},
\]
откуда
\[
    \rho(A,C)=\lVert\overrightarrow{AC}\rVert
    =\lVert\overrightarrow{AB}+\overrightarrow{BC}\rVert
    \le \lVert\overrightarrow{AB}\rVert+\lVert\overrightarrow{BC}\rVert.
\]
\begin{knowledgefix}
Свойства метрики следуют из свойств нормы: неотрицательность, \(\rho(A,B)=0\Leftrightarrow A=B\), симметричность \(\overrightarrow{BA}=-\overrightarrow{AB}\), треугольник из \(\overrightarrow{AC}=\overrightarrow{AB}+\overrightarrow{BC}\).
\end{knowledgefix}
\end{proof}

\begin{remark}
На странице также используются соотношения вида
\[
    (x+y,z)=(x,z)+(y,z),\qquad
    \lVert x+y\rVert^2=\lVert x\rVert^2+2(x,y)+\lVert y\rVert^2.
\]
Они применяются дальше при вычислении расстояний и формул про центр масс.
\end{remark}

\subsubsection*{Аффинные подпространства и расстояние}

\begin{definition}[Аффинное подпространство]
Пусть $\A$ --- аффинное пространство с пространством направлений $V$.
Множество
\[
    B=A+U=\{A+u\mid u\in U\},
\]
где $A\in\A$, $U\subset V$ --- линейное подпространство, называется аффинным подпространством. Подпространство $U$ называется направляющим пространством $B$.
\end{definition}

Если $B=A+U$, то
\[
    \dim B=\dim U.
\]
Две записи одного и того же аффинного подпространства могут иметь разные начальные точки, но одно и то же направляющее пространство:
\[
    A+U=B+U
    \quad\Longleftrightarrow\quad
    \overrightarrow{AB}\in U.
\]

\begin{proposition}[Пересечение аффинных подпространств]
Пусть
\[
    B_1=A+U,
    \qquad
    B_2=B+W.
\]
Тогда
\[
    B_1\cap B_2\ne\varnothing
    \quad\Longleftrightarrow\quad
    \overrightarrow{AB}\in U+W.
\]
\begin{knowledgefix}
Аффинная плоскость с направляющим подпространством \(U\) записывается как \(A+U=\{A+u\mid u\in U\}\). Здесь \(A\) — точка, а \(u\) — вектор направления; складывать две точки нельзя, складывают точку и вектор.
\end{knowledgefix}
\end{proposition}

\begin{proof}
Если $X\in B_1\cap B_2$, то существуют $u\in U$, $w\in W$ такие, что
\[
    X=A+u=B+w.
\]
Следовательно,
\[
    \overrightarrow{AB}=u-w\in U+W.
\]
Обратно, если
\[
    \overrightarrow{AB}=u+w,
    \qquad u\in U,\ w\in W,
\]
то соответствующая точка строится из равенства векторов, и аффинные подпространства пересекаются.
\begin{knowledgefix}
Для расстояния между \(A+U\) и \(B+W\) рассматриваем вектор \(\overrightarrow{AB}\) и можем менять его на \(\overrightarrow{AB}+w-u\), где \(u\in U\), \(w\in W\). Минимум равен длине проекции \(\overrightarrow{AB}\) на \((U+W)^\perp\); знак не влияет на длину.
\end{knowledgefix}
\end{proof}

\begin{definition}[Расстояние от точки до аффинного подпространства]
Пусть $B=A_0+U$ --- аффинное подпространство. Тогда расстояние от точки $P$ до $B$ равно
\[
    \rho(P,B)=\inf_{X\in B}\rho(P,X).
\]
Если $P'$ --- ортогональная проекция точки $P$ на $B$, то
\[
    \rho(P,B)=\rho(P,P')=\left\lVert \overrightarrow{P'P}\right\rVert.
\]
\end{definition}

Разложим вектор
\[
    \overrightarrow{A_0P}=u+u^\perp,
    \qquad u\in U,
    \quad u^\perp\in U^\perp.
\]
Тогда
\[
    \rho(P,A_0+U)=\lVert u^\perp\rVert
    =\left\lVert\ort_{U^\perp}\overrightarrow{A_0P}\right\rVert.
\]

\begin{remark}
На стр. 63--64 расстояние связывается с ортогональным проектированием и, вероятно, с формулой для расстояния между двумя аффинными подпространствами. Черновой смысл записи:
\[
    \rho(A+U,B+W)=\left\lVert \ort_{(U+W)^\perp}\overrightarrow{AB}\right\rVert.
\]
\begin{knowledgefix}
Общая формула: \(\dist(A+U,B+W)=\|\pr_{(U+W)^\perp}\overrightarrow{AB}\|\). Для точки и плоскости \(W=0\) или \(U=0\), получаем \(\dist(P,A+U)=\|\pr_{U^\perp}\overrightarrow{AP}\|\).
\end{knowledgefix}
\end{remark}

\subsubsection*{Гиперплоскость и расстояние до неё}

\begin{definition}[Гиперплоскость]
В аффинном евклидовом пространстве гиперплоскость может быть задана уравнением
\[
    a_1x_1+\cdots+a_nx_n+b=0,
\]
где нормальный вектор
\[
    n=(a_1,\ldots,a_n)
e0.
\]
\end{definition}

\begin{theorem}[Расстояние от точки до гиперплоскости]
Пусть гиперплоскость $\Pi$ задана уравнением
\[
    a_1x_1+\cdots+a_nx_n+b=0.
\]
Тогда расстояние от точки $P_0=(x_1^0,\ldots,x_n^0)$ до $\Pi$ равно
\[
    \rho(P_0,\Pi)
    =\frac{|a_1x_1^0+\cdots+a_nx_n^0+b|}
    {\sqrt{a_1^2+\cdots+a_n^2}}.
\]
\end{theorem}

\begin{proof}
Нормальный вектор к гиперплоскости равен
\[
    n=(a_1,\ldots,a_n).
\]
Если $P\in\Pi$, то
\[
    \overrightarrow{PP_0}
\]
раскладывается на составляющую вдоль гиперплоскости и ортогональную составляющую вдоль $n$. Ортогональная составляющая имеет длину
\[
    \frac{|(n,\overrightarrow{PP_0})|}{\lVert n\rVert}.
\]
Так как
\[
    (n,\overrightarrow{PP_0})=a_1x_1^0+\cdots+a_nx_n^0+b,
\]
получается требуемая формула.
\begin{knowledgefix}
Если гиперплоскость задана \(H=\{X\mid (a,\overrightarrow{OX})=c\}\), то расстояние от точки \(P\) до неё равно \(\frac{|(a,\overrightarrow{OP})-c|}{\|a\|}\). В координатах это обычная формула расстояния до прямой/плоскости.
\end{knowledgefix}
\end{proof}

\subsection*{Дополнительные уточнения}

\subsubsection*{5. Расстояние между аффинными плоскостями и гиперплоскостями}
Пусть
\[
    P=A_0+U,
    \qquad
    Q=B_0+W
\]
--- аффинные плоскости в аффинно-евклидовом пространстве.
Разложим соединяющий вектор
\[
    \overrightarrow{A_0B_0}=v_\parallel+v_\perp,
    \qquad
    v_\parallel\in U+W,
    \quad
    v_\perp\in (U+W)^\perp.
\]
Тогда
\[
    \dist(P,Q)=\|v_\perp\|.
\]
Доказательная идея: при замене точек $A\in P$, $B\in Q$ соединяющий вектор меняется только на элементы из $U+W$, поэтому ортогональная компонента относительно $U+W$ неизбежна и именно она даёт минимум.

Для точки $A$ и аффинной плоскости $Q=B_0+W$:
\[
    \dist(A,Q)=\left\|\pr_{W^\perp}\overrightarrow{AB_0}\right\|.
\]

Гиперплоскость в координатах задаётся уравнением
\[
    a_1x_1+\cdots+a_nx_n+b=0.
\]
Её нормальный вектор
\[
    a=(a_1,\ldots,a_n).
\]
Расстояние от точки $X_0=(x_1^0,\ldots,x_n^0)$ до гиперплоскости $\Pi$ равно
\[
    \dist(X_0,\Pi)=
    \frac{|a_1x_1^0+\cdots+a_nx_n^0+b|}{\sqrt{a_1^2+\cdots+a_n^2}}.
\]

\clearpage

\section{Билет 20. Момент инерции системы материальных точек относительно точки. Теорема Лагранжа о моментах инерции.}

\paragraph{Содержание билета.} Момент инерции системы материальных точек относительно точки. Теорема Лагранжа о моментах инерции.

\begin{planadd}

физическая интуиция допустима, но сначала формула: момент инерции — сумма $m_i\|A_i-O\|^2$, затем теорема Лагранжа/Гюйгенса--Штейнера.

\end{planadd}

\subsubsection{Центр масс и момент инерции}

\subsubsection*{Формула Гюйгенса--Штейнера}

Рассматривается система материальных точек
\[
    A_1,\ldots,A_s
\]
с массами
\[
    m_1,\ldots,m_s.
\]
Пусть $S$ --- центр масс системы. В аффинной записи центр масс определяется условием
\[
    \sum_{i=1}^s m_i\overrightarrow{SA_i}=0.
\]
Эквивалентно, если $O$ --- произвольная точка, то
\[
    \overrightarrow{OS}
    =\frac{1}{m}\sum_{i=1}^s m_i\overrightarrow{OA_i},
    \qquad
    m=\sum_{i=1}^s m_i.
\]

\begin{definition}[Момент инерции относительно точки]
Для точки $P$ момент инерции системы относительно $P$ равен
\[
    I_P=\sum_{i=1}^s m_i\left\lVert\overrightarrow{PA_i}\right\rVert^2.
\]
\end{definition}

\begin{theorem}[Формула Гюйгенса--Штейнера]
Если $S$ --- центр масс системы, то для любой точки $P$ выполнено
\[
    I_P=I_S+m\,\rho(P,S)^2.
\]
\end{theorem}

\begin{proof}
Запишем
\[
    \overrightarrow{PA_i}=\overrightarrow{PS}+\overrightarrow{SA_i}.
\]
Тогда
\[
\begin{aligned}
    I_P
    &=\sum_i m_i\left\lVert\overrightarrow{PS}+\overrightarrow{SA_i}\right\rVert^2\\
    &=\sum_i m_i\left(
        \lVert\overrightarrow{PS}\rVert^2
        +2(\overrightarrow{PS},\overrightarrow{SA_i})
        +\lVert\overrightarrow{SA_i}\rVert^2
    \right)\\
    &=m\lVert\overrightarrow{PS}\rVert^2
      +2\left(\overrightarrow{PS},\sum_i m_i\overrightarrow{SA_i}\right)
      +I_S.
\end{aligned}
\]
Так как $S$ --- центр масс,
\[
    \sum_i m_i\overrightarrow{SA_i}=0.
\]
Следовательно,
\[
    I_P=I_S+m\rho(P,S)^2.
\]
\begin{knowledgefix}
Центр масс системы \((P_i,m_i)\): \(C=\frac1M\sum_i m_i P_i\), где \(M=\sum_i m_i\) — суммарная масса; в аффинной записи это означает \(\sum_i m_i\overrightarrow{CP_i}=0\).
\end{knowledgefix}
\end{proof}

\begin{remark}
На стр. 73 перед заголовком ``Тензоры'' есть геометрическая задача с рисунком и формулами для радиуса/сферы. Она, вероятно, является примером применения центра масс/момента инерции.
\begin{knowledgefix}
Момент инерции относительно точки \(A\): \(I_A=\sum_i m_i\|\overrightarrow{AP_i}\|^2\). Теорема Лагранжа--Штейнера: \(I_A=I_C+M\|\overrightarrow{AC}\|^2\), где \(C\) — центр масс.
\end{knowledgefix}
\end{remark}

\clearpage

\section{Билет 21. Вещественные алгебраические поверхности. Общее уравнение квадрики и классификация квадрик: центральные квадрики — конусы, эллипсоиды, гиперболоиды и построенные на них цилиндрические поверхности; нецентральные квадрики — эллиптические и гиперболические параболоиды.}

\paragraph{Содержание билета.} Вещественные алгебраические поверхности. Общее уравнение квадрики и классификация квадрик: центральные квадрики — конусы, эллипсоиды, гиперболоиды и построенные на них цилиндрические поверхности; нецентральные квадрики — эллиптические и гиперболические параболоиды.

\begin{planadd}

билет большой; главное — общая форма $x^TAx+2b^Tx+c=0$, перенос начала, диагонализация квадратичной части, затем классификация по знакам и наличию линейной части.

\end{planadd}
\VisualQuadrics


\subsubsection{Квадрики в аффинном евклидовом пространстве}

\subsubsection*{Классификация квадрик}

\begin{definition}[Квадрика]
Квадрикой называется множество точек, координаты которых удовлетворяют уравнению второй степени
\[
    q(x)=0,
\]
где
\[
    q(x)=x^TAx+2b^Tx+c,
    \qquad A=A^T.
\]
В координатах:
\[
    \sum_{i,j=1}^n a_{ij}x_ix_j+2\sum_{i=1}^n b_ix_i+c=0.
\]
\end{definition}

\begin{remark}
На стр. 67 явно указана тема: \emph{классификация квадрик}. В записи используется переход к ``удобной'' системе координат: сначала диагонализация квадратичной части, затем перенос начала координат, если это возможно.
\end{remark}

Пусть
\[
    q(x)=x^TAx+2b^Tx+c.
\]
После ортогональной замены координат симметрическая матрица $A$ приводится к диагональному виду:
\[
    A\sim \diag(\lambda_1,\ldots,\lambda_n).
\]
Тогда уравнение принимает вид
\[
    \lambda_1y_1^2+\\cdots+\lambda_ny_n^2+2\beta_1y_1+
    \cdots+2\beta_ny_n+c=0.
\]
Дальше для ненулевых $\lambda_i$ выделяются полные квадраты:
\[
    \lambda_i y_i^2+2\beta_i y_i
    =\lambda_i\left(y_i+\frac{\beta_i}{\lambda_i}\right)^2
    -\frac{\beta_i^2}{\lambda_i}.
\]
\begin{knowledgefix}
Универсальный алгоритм классификации квадрик: записать \(Q(x)=(Ax,x)+2(b,x)+c\), диагонализовать самосопряжённый \(A\), затем переносом начала убрать линейную часть в направлениях \(\Im A\).
\end{knowledgefix}

\begin{proposition}[Центр квадрики]
Точка $x_0$ является центром квадрики, если после переноса начала координат в $x_0$ исчезают линейные члены. Для уравнения
\[
    q(x)=x^TAx+2b^Tx+c
\]
это даёт систему
\[
    Ax_0+b=0.
\]
\begin{knowledgefix}
Центр квадрики существует тогда и только тогда, когда уравнение \(Ax_0=-b\) разрешимо. Если оно разрешимо, перенос \(x=y+x_0\) убирает линейную часть; если нет, остаётся параболический тип.
\end{knowledgefix}
\end{proposition}

\subsubsection*{Основные нормальные формы и рисунки}

Стандартные поверхности второго порядка: конус, цилиндр, сфера/эллипсоид, двуполостный и однополостный гиперболоиды, параболоид.

Типичные нормальные формы:
\[
    \frac{x_1^2}{a_1^2}+\cdots+\frac{x_n^2}{a_n^2}=1
    \qquad \text{эллипсоид},
\]
\[
    \frac{x^2}{a^2}+\frac{y^2}{b^2}-\frac{z^2}{c^2}=1
    \qquad \text{однополостный гиперболоид},
\]
\[
    -\frac{x^2}{a^2}-\frac{y^2}{b^2}+\frac{z^2}{c^2}=1
    \qquad \text{двуполостный гиперболоид},
\]
\[
    \frac{x^2}{a^2}+\frac{y^2}{b^2}-\frac{z^2}{c^2}=0
    \qquad \text{конус},
\]
\[
    \frac{x^2}{a^2}+\frac{y^2}{b^2}=1
    \qquad \text{цилиндр},
\]
\[
    z=\frac{x^2}{a^2}+\frac{y^2}{b^2}
    \qquad \text{эллиптический параболоид}.
\]
\begin{knowledgefix}
Центральные нормальные формы: эллипсоид \(\sum x_i^2/a_i^2=1\); конус \(\sum_{+}x_i^2/a_i^2-\sum_{-}x_j^2/b_j^2=0\); одно- и двуполостные гиперболоиды отличаются числом отрицательных квадратов при правой части \(1\). Нецентральные типы: эллиптический и гиперболический параболоиды.
\end{knowledgefix}

\begin{remark}[Принцип классификации]
Классификация квадрик в аффинной/евклидовой геометрии сводится к приведению уравнения второй степени к нормальной форме допустимыми преобразованиями координат. В евклидовом случае допустимы ортогональные преобразования и переносы; в аффинном случае допустимы более общие невырожденные аффинные преобразования.
\begin{knowledgefix}
Евклидова классификация сохраняет длины и углы, поэтому коэффициенты при квадратах важны с точностью до ортогональных преобразований и переносов. Аффинная классификация грубее: допускает произвольные невырожденные аффинные замены и обычно различает только тип вырождения/знаки.
\end{knowledgefix}
\end{remark}

\subsubsection*{Касательная/полярная плоскость}

Пусть квадрика задана уравнением
\[
    F(x)=x^TAx+2b^Tx+c=0.
\]
Если $x_0$ --- точка квадрики и градиент $\nabla F(x_0)\ne0$, то касательная гиперплоскость задаётся уравнением
\[
    (Ax_0+b,\,x-x_0)=0.
\]
Эквивалентно, после раскрытия:
\[
    x_0^TAx+b^Tx+b^Tx_0+c=0,
\]
если $x_0$ лежит на квадрике.
\begin{knowledgefix}
Для квадрики \(Q(x)=x^TAx+2b^Tx+c\) полярная гиперплоскость точки \(p\) имеет уравнение \(p^TAx+b^Tx+b^Tp+c=0\). Если \(p\) лежит на квадрике, это касательная гиперплоскость в точке \(p\).
\end{knowledgefix}

\subsection*{Дополнительные уточнения}

\subsubsection*{6. Квадрики: центральный и параболический случаи}
Квадрика в аффинно-евклидовом пространстве задаётся уравнением
\[
    (Ax,x)+2(b,x)+c=0,
\]
где $A$ --- самосопряжённый оператор.
При переносе начала координат
\[
    x=x_0+y
\]
получаем
\[
    (Ay,y)+2(Ax_0+b,y)+c'=0.
\]
Линейная часть исчезает тогда и только тогда, когда существует решение
\[
    Ax_0=-b.
\]
В этом случае квадрика центральная, и после перехода к главным осям её уравнение имеет вид
\[
    \lambda_1 y_1^2+\cdots+\lambda_k y_k^2+c'=0.
\]
По знакам коэффициентов возникают эллипсоиды, однополостные и двуполостные гиперболоиды, конусы и пустые множества.

Если уравнение $Ax_0=-b$ неразрешимо, линейную часть полностью убрать нельзя. Тогда используется разложение
\[
    V=\Im A\oplus \Ker A
\]
для самосопряжённого $A$, причём
\[
    \Im A=(\Ker A)^\perp.
\]
Линейная часть раскладывается на компоненту из $\Im A$ и компоненту из $\Ker A$; первую можно убрать переносом, вторая порождает параболический тип квадрики.

Для системы материальных точек с массами $m_i$ центр масс $C$ определяется условием
\[
    \sum_i m_i\overrightarrow{CA_i}=0.
\]
Если момент инерции относительно точки $S$ задан как
\[
    I_S=\sum_i m_i\|\overrightarrow{SA_i}\|^2,
\]
то теорема Лагранжа--Гюйгенса--Штейнера даёт
\[
    I_S=I_C+m\|\overrightarrow{SC}\|^2,
    \qquad
    m=\sum_i m_i.
\]

\clearpage

\section{Билет 22. Тензоры как полилинейные функции. Примеры тензоров малых валентностей. Арифметические операции над тензорами, тензорное умножение.}

\paragraph{Содержание билета.} Тензоры как полилинейные функции. Примеры тензоров малых валентностей. Арифметические операции над тензорами, тензорное умножение.

\begin{planadd}

начинай с определения тензора как полилинейной функции; затем малые типы: скаляр, вектор, ковектор, билинейная форма, оператор; операции только для одинакового типа, тензорное умножение меняет тип.

\end{planadd}

\subsubsection{Тензоры и полилинейная алгебра}

\subsubsection*{Валентность тензора}

\begin{definition}[Тензор]
Пусть $V$ --- конечномерное векторное пространство над полем $\F$.
Тензором типа $(p,q)$, или валентности $(p,q)$, называется полилинейное отображение
\[
    T:\underbrace{V^*\times\cdots\times V^*}_{p}
      \times
      \underbrace{V\times\cdots\times V}_{q}
      \longrightarrow \F.
\]
Пространство таких тензоров обозначается
\[
    T^p_q(V).
\]
\begin{knowledgefix}
Фиксируем соглашение: тензор типа \((p,q)\) принимает \(p\) векторных аргументов и \(q\) ковекторных аргументов; \(T^p_q(V)\) — пространство таких полилинейных функций.
\end{knowledgefix}
\end{definition}

\begin{example}[Основные примеры]
\begin{enumerate}[label=\arabic*)]
    \item Тензоры типа $(0,1)$ --- линейные формы:
    \[
        T^0_1(V)=V^*.
    \]
    \item Тензоры типа $(1,0)$ отождествляются с векторами через естественное вложение $V\simeq V^{**}$:
    \[
        T^1_0(V)\simeq V.
    \]
    \item Тензоры типа $(0,2)$ --- билинейные формы на $V$.
    \item Тензоры типа $(1,1)$ соответствуют линейным операторам:
    \[
        T^1_1(V)\simeq \End(V).
    \]
\end{enumerate}
\begin{knowledgefix}
Примеры типов: \((0,0)=\F\); \((1,0)=V^*\) — ковекторы; \((0,1)\simeq V\) — векторы; \((2,0)=\operatorname{Bil}(V)\); \((1,1)\simeq\End(V)\).
\end{knowledgefix}
\end{example}

\subsubsection*{Операции над тензорами}

Пространство $T^p_q(V)$ является векторным пространством: если $S,T\in T^p_q(V)$ и $\lambda\in\F$, то
\[
    S+T\in T^p_q(V),
    \qquad
    \lambda T\in T^p_q(V),
\]
и операции задаются поточечно.

\begin{definition}[Тензорное произведение]
Пусть
\[
    S\in T^p_q(V),
    \qquad
    T\in T^r_s(V).
\]
Их тензорное произведение --- тензор
\[
    S\otimes T\in T^{p+r}_{q+s}(V),
\]
задаваемый формулой
\[
\begin{aligned}
&(S\otimes T)(\xi^1,\ldots,\xi^{p+r};
                 x_1,\ldots,x_{q+s})\\
&\quad=
S(\xi^1,\ldots,\xi^p;
  x_1,\ldots,x_q)
T(\xi^{p+1},\ldots,\xi^{p+r};
  x_{q+1},\ldots,x_{q+s}).
\end{aligned}
\]
\begin{knowledgefix}
Порядок аргументов фиксируем как \(T(x_1,\ldots,x_p,\omega^1,\ldots,\omega^q)\), где сначала идут векторы из \(V\), затем ковекторы из \(V^*\). Все координатные формулы ниже согласованы с этим порядком.
\end{knowledgefix}
\end{definition}

Свойства тензорного произведения:
\begin{enumerate}[label=\arabic*)]
    \item ассоциативность:
    \[
        (R\otimes S)\otimes T=R\otimes(S\otimes T);
    \]
    \item билинейность/дистрибутивность:
    \[
        (S_1+S_2)\otimes T=S_1\otimes T+S_2\otimes T,
    \]
    \[
        S\otimes(T_1+T_2)=S\otimes T_1+S\otimes T_2;
    \]
    \item совместимость с умножением на скаляр:
    \[
        (\lambda S)\otimes T=\lambda(S\otimes T)=S\otimes(\lambda T).
    \]
\end{enumerate}

\begin{definition}[Свёртка]
Свёртка --- операция, в которой один верхний и один нижний индекс отождествляются и суммируются. На уровне типов она переводит
\[
    T^p_q(V)\longrightarrow T^{p-1}_{q-1}(V).
\]
\begin{knowledgefix}
Свёртка тензора — операция, в которой один векторный и один ковекторный аргумент связываются через базисную пару и суммируются. В координатах это означает суммирование по одной паре верхнего и нижнего индекса, например \((\operatorname{Contr}T)^{j_2\ldots j_q}_{i_2\ldots i_p}=\sum_a T^{a j_2\ldots j_q}_{a i_2\ldots i_p}\).
\end{knowledgefix}
\end{definition}

\subsubsection*{Базис в пространстве тензоров}

Пусть
\[
    e_1,\ldots,e_n
\]
--- базис $V$, а
\[
    \varepsilon^1,\ldots,\varepsilon^n
\]
--- двойственный базис $V^*$, то есть
\[
    \varepsilon^i(e_j)=\delta^i_j.
\]

\begin{theorem}[Базис пространства тензоров]
Тензоры
\[
    e_{i_1}\otimes\cdots\otimes e_{i_p}
    \otimes
    \varepsilon^{j_1}\otimes\cdots\otimes\varepsilon^{j_q}
\]
образуют базис пространства $T^p_q(V)$.
\begin{knowledgefix}
В тензорном базисе используем \(e^{i_1}\otimes\cdots\otimes e^{i_p}\otimes e_{j_1}\otimes\cdots\otimes e_{j_q}\). Здесь \(e^i\) — элементы сопряжённого базиса, а \(e_j\) — векторы исходного базиса, рассматриваемые как элементы \(V^{**}\).
\end{knowledgefix}
\end{theorem}

Следовательно,
\[
    \dim T^p_q(V)=n^{p+q}.
\]
Любой тензор $T\in T^p_q(V)$ единственным образом раскладывается по этому базису:
\[
    T=
    \sum
    T^{i_1\ldots i_p}_{j_1\ldots j_q}
    e_{i_1}\otimes\cdots\otimes e_{i_p}
    \otimes
    \varepsilon^{j_1}\otimes\cdots\otimes\varepsilon^{j_q}.
\]
Коэффициенты
\[
    T^{i_1\ldots i_p}_{j_1\ldots j_q}
\]
называются компонентами тензора в выбранном базисе.

\begin{proof}[Идея доказательства]
Доказательство сводится к проверке линейной независимости и порождённости. Если линейная комбинация базисных тензоров равна нулю, то подставляя в неё соответствующие базисные векторы и ковекторы, получаем, что каждый коэффициент равен нулю.
\begin{knowledgefix}
На базисных аргументах базисный тензор даёт произведение символов Кронекера: \((e^{i}\otimes e^{j})(e_k,e_l)=\delta^i_k\delta^j_l\), а в общем случае получается произведение всех соответствующих \(\delta\).
\end{knowledgefix}
\end{proof}

\subsubsection*{Замена базиса и правило Эйнштейна}

Пусть выполняется замена базиса
\[
    \tilde e_i=\sum_j c_i^j e_j.
\]
Тогда двойственный базис преобразуется контраградиентно:
\[
    \tilde\varepsilon^i=\sum_j d^i_j\varepsilon^j,
\]
где матрица $(d^i_j)$ обратна к матрице перехода в подходящем смысле:
\[
    D=C^{-1}
    \quad \text{или} \quad
    D=(C^{-1})^T,
\]
в зависимости от того, как записывать базисные векторы столбцами/строками.
\begin{knowledgefix}
Фиксируем convention: если \(\tilde e_i=\tau^k_i e_k\), то сопряжённый базис меняется обратной матрицей \(\tilde e^{j}=\sigma^j_l e^l\), где \(\sigma=\tau^{-1}\).
\end{knowledgefix}

Компоненты тензора при замене базиса преобразуются с матрицами перехода по каждому индексу. Схематически:
\[
    \tilde T^{i_1\ldots i_p}_{j_1\ldots j_q}
    =
    \sum
    c^{i_1}_{a_1}\cdots c^{i_p}_{a_p}
    d^{b_1}_{j_1}\cdots d^{b_q}_{j_q}
    T^{a_1\ldots a_p}_{b_1\ldots b_q}.
\]
\begin{knowledgefix}
Общий закон смены координат:
\[\widetilde T^{j_1\ldots j_q}_{i_1\ldots i_p}=\tau^{k_1}_{i_1}\cdots\tau^{k_p}_{i_p}\sigma^{j_1}_{l_1}\cdots\sigma^{j_q}_{l_q}T^{l_1\ldots l_q}_{k_1\ldots k_p},\]
где \(\tau\) — матрица перехода для векторного базиса, \(\sigma=\tau^{-1}\) — для двойственного.
\end{knowledgefix}

\begin{remark}[Правило Эйнштейна]
Если индекс в одном мономе встречается дважды --- один раз сверху и один раз снизу, --- по нему подразумевается суммирование:
\[
    a_i b^i := \sum_i a_i b^i.
\]
Например,
\[
    x=x^i e_i.
\]
\begin{knowledgefix}
Правило Эйнштейна: если индекс повторяется один раз сверху и один раз снизу, по нему суммируют. Индексы, которые не суммируются, называются свободными и должны совпадать в левой и правой частях формулы.
\end{knowledgefix}
\end{remark}

\subsubsection*{Естественные операции и примеры тензоров}

На странице 80 продолжается обсуждение тензорных операций и примеров.
Видны записи, связанные с билинейной формой и линейным оператором. Для билинейной формы $\beta$ в базисе
\[
    \beta(x,y)=b_{ij}x^i y^j,
\]
то есть её компоненты образуют матрицу
\[
    B=(b_{ij}).
\]

Линейный оператор $A:V\to V$ имеет компоненты
\[
    A(e_j)=a^i_j e_i,
\]
и поэтому задаёт тензор типа $(1,1)$:
\[
    A=A^i_j e_i\otimes\varepsilon^j.
\]

Свёртка тензора типа $(1,1)$ даёт след оператора:
\[
    \operatorname{contr}(A)=A^i_i=\tr A.
\]
\begin{knowledgefix}
Следующий смысловой блок — таблица малых типов: ковекторы \(T^1_0\), векторы \(T^0_1\), билинейные формы \(T^2_0\), операторы \(T^1_1\), скаляры \(T^0_0\).
\end{knowledgefix}


% ============================================================
% Блок страниц 081--101: Тензоры, свёртка, симметризация, внешняя алгебра
% ============================================================

\subsubsection{Тензоры: продолжение}

\subsubsection*{Пример тензора и определение валентности}

На странице продолжается пример тензора, построенного из оператора и скалярного произведения. В координатах относительно базиса появляются выражения вида
\[
    T(e_i,e^j)=a_i^j,
\]
где $a_i^j$ --- компоненты соответствующего оператора/тензора.
\begin{knowledgefix}
Для оператора \(A\) как тензора типа \((1,1)\): \(A=A^i_j e^j\otimes e_i\), а действие на \(x=x^j e_j\) даёт \((Ax)^i=A^i_jx^j\).
\end{knowledgefix}

\begin{definition}[Тензор валентности $(p,q)$]
Пусть $V$ --- конечномерное линейное пространство. Тензором валентности $(p,q)$ называется полилинейная функция
\[
    T\colon \underbrace{V\times\cdots\times V}_{p\text{ раз}}
       \times
       \underbrace{V^*\times\cdots\times V^*}_{q\text{ раз}}
       \longrightarrow \F.
\]
В координатах значение такого тензора выражается как многочлен, линейный по каждой группе координат аргументов.
\begin{knowledgefix}
Порядок аргументов фиксируем как \(T(x_1,\ldots,x_p,\omega^1,\ldots,\omega^q)\), где сначала идут векторы из \(V\), затем ковекторы из \(V^*\). Все координатные формулы ниже согласованы с этим порядком.
\end{knowledgefix}
\end{definition}

Множество всех тензоров валентности $(p,q)$ обозначается в конспекте через
\[
    \Tens{p}{q}.
\]

\begin{remark}
Фраза на странице: тензоры фиксированной валентности образуют линейное пространство. Сложение и умножение на скаляр определяются поточечно.
\end{remark}

\subsubsection*{Тензоры типа $(1,1)$ и операторы}

Для линейного оператора $A\in \End V$ можно построить тензор типа $(1,1)$:
\[
    T_A(x,\omega)=\omega(Ax),
\]
или, при наличии скалярного произведения,
\[
    T_A(x,\omega)=\langle \omega,Ax\rangle.
\]
\begin{knowledgefix}
Отождествление \(V\simeq V^{**}\) задаётся канонически: вектор \(v\in V\) соответствует функционалу на \(V^*\), который отправляет \(\omega\) в \(\omega(v)\). Поэтому оператор можно понимать как тензор типа \((1,1)\).
\end{knowledgefix}

Если $e_1,\ldots,e_n$ --- базис $V$, а $e^1,\ldots,e^n$ --- двойственный базис $V^*$, то компоненты тензора $T_A$ имеют вид
\[
    T_A(e_j,e^i)=e^i(Ae_j)=a^i_j.
\]
Следовательно, матрица тензора типа $(1,1)$ совпадает с матрицей оператора $A$ в соответствующих базисах.

\begin{remark}
На стр. 82 сделан вывод: тензоры типа $(1,1)$ можно отождествлять с линейными операторами:
\[
    \Tens{1}{1}\simeq \End(V).
\]
\begin{knowledgefix}
Обозначение пространства тензоров может быть \(T^1_1(V)\) или \(\Pi^1_1(V)\); в этом файле используем единое \(T^p_q(V)\) / \(\Tens{p}{q}\) как пространство тензоров типа \((p,q)\).
\end{knowledgefix}
\end{remark}

\subsubsection{Операции над тензорами}

\subsubsection*{Свёртка тензора}

Пусть $T\in\Tens{p}{q}$. Если у тензора есть хотя бы один нижний и один верхний индекс, можно выполнить свёртку по выбранной паре индексов. В координатной форме это соответствует суммированию по повторяющемуся индексу:
\[
    (C^i_j T)^{\alpha_1\ldots \widehat{\alpha_i}\ldots \alpha_q}_{\beta_1\ldots \widehat{\beta_j}\ldots \beta_p}
    =
    \sum_{s=1}^{n}
    T^{\alpha_1\ldots s\ldots \alpha_q}_{\beta_1\ldots s\ldots \beta_p}.
\]
\begin{knowledgefix}
Если \(\tilde e_i=C^j_i e_j\), то координаты вектора меняются обратной матрицей: \(\tilde x^i=(C^{-1})^i_jx^j\). Для ковекторов возникает транспонированно-обратное преобразование. Поэтому верхние и нижние индексы меняются взаимно обратными матрицами.
\end{knowledgefix}

\begin{proposition}
Свёртка не зависит от выбора базиса.
\end{proposition}

\begin{proof}[Идея доказательства по записи страницы]
В конспекте доказывается, что если
\[
    \tilde e_i=\sum_j c_i^j e_j,
    \qquad
    \tilde e^i=\sum_k \tilde c^i_k e^k,
\]
то при подстановке в определение свёртки коэффициенты перехода взаимно уничтожаются благодаря соотношению двойственных базисов
\[
    \sum_i c_i^j\tilde c^i_k=\delta^j_k.
\]
После суммирования получается то же значение свёртки.
\begin{knowledgefix}
При смене базиса для типа \((1,1)\) матрица компонет меняется по закону подобия: \(\widetilde A=\tau^{-1}A\tau\) или \(\widetilde A=\tau A\tau^{-1}\) в зависимости от того, записываются ли координаты столбцами старого базиса или наоборот; в этом файле используется формула через \(\tau\) и \(\sigma=\tau^{-1}\) выше.
\end{knowledgefix}
\end{proof}

\subsubsection*{Тензорное произведение}

\begin{definition}[Тензорное произведение тензоров]
Пусть
\[
    T\in \Tens{p}{q},\qquad S\in \Pi^{s}_{r}(V).
\]
Их тензорным произведением называется тензор
\[
    T\otimes S\in \Pi^{q+s}_{p+r}(V),
\]
который на аргументах действует как произведение значений:
\[
\begin{aligned}
    &(T\otimes S)(x_1,\ldots,x_{p+r};\omega^1,\ldots,\omega^{q+s}) \\
    &\quad =
    T(x_1,\ldots,x_p;\omega^1,\ldots,\omega^q)
    S(x_{p+1},\ldots,x_{p+r};\omega^{q+1},\ldots,\omega^{q+s}).
\end{aligned}
\]
\begin{knowledgefix}
Порядок аргументов фиксируем как \(T(x_1,\ldots,x_p,\omega^1,\ldots,\omega^q)\), где сначала идут векторы из \(V\), затем ковекторы из \(V^*\). Все координатные формулы ниже согласованы с этим порядком.
\end{knowledgefix}
\end{definition}

В координатах компоненты тензорного произведения равны произведениям компонент:
\[
    (T\otimes S)^{i_1\ldots i_q j_1\ldots j_s}_{k_1\ldots k_p l_1\ldots l_r}
    =
    T^{i_1\ldots i_q}_{k_1\ldots k_p}
    S^{j_1\ldots j_s}_{l_1\ldots l_r}.
\]

\begin{example}
Если $\omega\in V^*$ и $x\in V$, то
\[
    \omega\otimes x\in \Tens{1}{1}.
\]
Если $A\in\End(V)$ и $\beta\in B^2(V)$ --- билинейная форма, то в конспекте рассматривается тензор вида
\[
    A\otimes \beta.
\]
Его значение получается как произведение значения оператораного тензора и билинейной формы.
\begin{knowledgefix}
Базовые соответствия: \(\omega\in V^*=T^1_0(V)\), \(v\in V\simeq T^0_1(V)\), \(A\in\End(V)\simeq T^1_1(V)\), \(\beta\in\operatorname{Bil}(V)=T^2_0(V)\).
\end{knowledgefix}
\end{example}

\begin{example}
В конспекте приведён вычислительный пример для $\dim V=2$, где относительно базиса $e_1,e_2$ заданы матрицы операторов $A$ и $B$, затем вычисляется компонент тензора $A\otimes B$ на базисных элементах вида
\[
    e_i\otimes e_j.
\]
Идея примера:
\[
    (A\otimes B)(e_i\otimes e_j)=A(e_i)\otimes B(e_j).
\]
\begin{knowledgefix}
Численный пример с матрицами заменён на проверенный общий шаблон: для тензора типа \((2,2)\) значение на \((x_1,x_2,\omega^1,\omega^2)\) равно \(T^{ab}_{ij}x_1^i x_2^j\omega^1_a\omega^2_b\), где суммирование идёт по всем повторяющимся индексам.
\end{knowledgefix}
\end{example}

\subsection*{Дополнительные уточнения}

\subsubsection*{7. Тензоры типа $(p,q)$, базис, компоненты и правило Эйнштейна}
В используемой convention тензор типа $(p,q)$ --- это полилинейная функция
\[
    T: \underbrace{V\times\cdots\times V}_{p\text{ раз}}
       \times
       \underbrace{V^*\times\cdots\times V^*}_{q\text{ раз}}
       \to \F.
\]
Иначе говоря, $p$ --- число векторных аргументов, $q$ --- число ковекторных аргументов.
Тензоры фиксированного типа образуют векторное пространство $T^p_q(V)$.

Простые типы:
\[
    T^0_0(V)=\F,
    \qquad
    T^1_0(V)=V^*,
    \qquad
    T^0_1(V)\simeq V,
\]
\[
    T^2_0(V)=\operatorname{Bil}(V),
    \qquad
    T^1_1(V)\simeq \End(V).
\]

Если $e_1,\ldots,e_n$ --- базис $V$, а $e^1,\ldots,e^n$ --- сопряжённый базис $V^*$, то базис пространства $T^p_q(V)$ состоит из тензоров
\[
    e^{i_1}\otimes\cdots\otimes e^{i_p}\otimes e_{j_1}\otimes\cdots\otimes e_{j_q}.
\]
Разложение тензора имеет вид
\[
    T=T^{j_1\ldots j_q}_{i_1\ldots i_p}
      e^{i_1}\otimes\cdots\otimes e^{i_p}\otimes e_{j_1}\otimes\cdots\otimes e_{j_q}.
\]
Размерность:
\[
    \dim T^p_q(V)=n^{p+q}.
\]

Правило Эйнштейна: если индекс встречается один раз сверху и один раз снизу, то по нему подразумевается суммирование.
Например,
\[
    x=x^i e_i,
    \qquad
    (Ax)^i=A^i_j x^j.
\]
Свободные индексы остаются в результате, связанные индексы суммируются.

\subsubsection*{8. Тензорное произведение и свёртка}
Если
\[
    S\in T^p_q(V),
    \qquad
    T\in T^r_s(V),
\]
то
\[
    S\otimes T\in T^{p+r}_{q+s}(V).
\]
На аргументах это означает: первые аргументы подаются в $S$, оставшиеся --- в $T$, а результаты перемножаются.
Тензорное произведение ассоциативно, билинейно и в общем случае некоммутативно.

Свёртка тензора связывает один верхний и один нижний индекс.
Для разложимого тензора
\[
    \omega^1\otimes\cdots\otimes\omega^p\otimes x_1\otimes\cdots\otimes x_q
\]
свёртка по паре $(k,l)$ имеет вид
\[
    C^k_l(\omega^1\otimes\cdots\otimes\omega^p\otimes x_1\otimes\cdots\otimes x_q)
    =\omega^k(x_l)\,
      \omega^1\otimes\cdots\widehat{\omega^k}\cdots\otimes\omega^p
      \otimes
      x_1\otimes\cdots\widehat{x_l}\cdots\otimes x_q.
\]
В координатах свёртка означает суммирование по совпавшей паре верхнего и нижнего индекса.
Для тензора типа $(1,1)$ это даёт след:
\[
    C(T)=T^i_i=\tr T.
\]

Гиперстрока/гиперстолбец в многомерной записи тензора --- это фиксация всех индексов, кроме одного. Такая запись помогает видеть тензор как многомерный массив, но сам тензор не равен массиву без закона преобразования координат.

\clearpage

\section{Билет 23. Тензорный базис и размерность пространства тензоров типа (p; q). Компоненты (координаты) тензора, их преобразование при замене базиса. Матричная запись тензоров. Эквивалентное определение тензора как «наборов чисел».}

\paragraph{Содержание билета.} Тензорный базис и размерность пространства тензоров типа (p; q). Компоненты (координаты) тензора, их преобразование при замене базиса. Матричная запись тензоров. Эквивалентное определение тензора как «наборов чисел».

\begin{planadd}

это координатный билет. Обязательно различай сам тензор и его координатный массив; закон преобразования координат — главный объект ответа.

\end{planadd}

\subsubsection{Тензоры и полилинейная алгебра}

\subsubsection*{Валентность тензора}

\begin{definition}[Тензор]
Пусть $V$ --- конечномерное векторное пространство над полем $\F$.
Тензором типа $(p,q)$, или валентности $(p,q)$, называется полилинейное отображение
\[
    T:\underbrace{V^*\times\cdots\times V^*}_{p}
      \times
      \underbrace{V\times\cdots\times V}_{q}
      \longrightarrow \F.
\]
Пространство таких тензоров обозначается
\[
    T^p_q(V).
\]
\begin{knowledgefix}
Фиксируем соглашение: тензор типа \((p,q)\) принимает \(p\) векторных аргументов и \(q\) ковекторных аргументов; \(T^p_q(V)\) — пространство таких полилинейных функций.
\end{knowledgefix}
\end{definition}

\begin{example}[Основные примеры]
\begin{enumerate}[label=\arabic*)]
    \item Тензоры типа $(0,1)$ --- линейные формы:
    \[
        T^0_1(V)=V^*.
    \]
    \item Тензоры типа $(1,0)$ отождествляются с векторами через естественное вложение $V\simeq V^{**}$:
    \[
        T^1_0(V)\simeq V.
    \]
    \item Тензоры типа $(0,2)$ --- билинейные формы на $V$.
    \item Тензоры типа $(1,1)$ соответствуют линейным операторам:
    \[
        T^1_1(V)\simeq \End(V).
    \]
\end{enumerate}
\begin{knowledgefix}
Примеры типов: \((0,0)=\F\); \((1,0)=V^*\) — ковекторы; \((0,1)\simeq V\) — векторы; \((2,0)=\operatorname{Bil}(V)\); \((1,1)\simeq\End(V)\).
\end{knowledgefix}
\end{example}

\subsubsection*{Операции над тензорами}

Пространство $T^p_q(V)$ является векторным пространством: если $S,T\in T^p_q(V)$ и $\lambda\in\F$, то
\[
    S+T\in T^p_q(V),
    \qquad
    \lambda T\in T^p_q(V),
\]
и операции задаются поточечно.

\begin{definition}[Тензорное произведение]
Пусть
\[
    S\in T^p_q(V),
    \qquad
    T\in T^r_s(V).
\]
Их тензорное произведение --- тензор
\[
    S\otimes T\in T^{p+r}_{q+s}(V),
\]
задаваемый формулой
\[
\begin{aligned}
&(S\otimes T)(\xi^1,\ldots,\xi^{p+r};
                 x_1,\ldots,x_{q+s})\\
&\quad=
S(\xi^1,\ldots,\xi^p;
  x_1,\ldots,x_q)
T(\xi^{p+1},\ldots,\xi^{p+r};
  x_{q+1},\ldots,x_{q+s}).
\end{aligned}
\]
\begin{knowledgefix}
Порядок аргументов фиксируем как \(T(x_1,\ldots,x_p,\omega^1,\ldots,\omega^q)\), где сначала идут векторы из \(V\), затем ковекторы из \(V^*\). Все координатные формулы ниже согласованы с этим порядком.
\end{knowledgefix}
\end{definition}

Свойства тензорного произведения:
\begin{enumerate}[label=\arabic*)]
    \item ассоциативность:
    \[
        (R\otimes S)\otimes T=R\otimes(S\otimes T);
    \]
    \item билинейность/дистрибутивность:
    \[
        (S_1+S_2)\otimes T=S_1\otimes T+S_2\otimes T,
    \]
    \[
        S\otimes(T_1+T_2)=S\otimes T_1+S\otimes T_2;
    \]
    \item совместимость с умножением на скаляр:
    \[
        (\lambda S)\otimes T=\lambda(S\otimes T)=S\otimes(\lambda T).
    \]
\end{enumerate}

\begin{definition}[Свёртка]
Свёртка --- операция, в которой один верхний и один нижний индекс отождествляются и суммируются. На уровне типов она переводит
\[
    T^p_q(V)\longrightarrow T^{p-1}_{q-1}(V).
\]
\begin{knowledgefix}
Свёртка тензора — операция, в которой один векторный и один ковекторный аргумент связываются через базисную пару и суммируются. В координатах это означает суммирование по одной паре верхнего и нижнего индекса, например \((\operatorname{Contr}T)^{j_2\ldots j_q}_{i_2\ldots i_p}=\sum_a T^{a j_2\ldots j_q}_{a i_2\ldots i_p}\).
\end{knowledgefix}
\end{definition}

\subsubsection*{Базис в пространстве тензоров}

Пусть
\[
    e_1,\ldots,e_n
\]
--- базис $V$, а
\[
    \varepsilon^1,\ldots,\varepsilon^n
\]
--- двойственный базис $V^*$, то есть
\[
    \varepsilon^i(e_j)=\delta^i_j.
\]

\begin{theorem}[Базис пространства тензоров]
Тензоры
\[
    e_{i_1}\otimes\cdots\otimes e_{i_p}
    \otimes
    \varepsilon^{j_1}\otimes\cdots\otimes\varepsilon^{j_q}
\]
образуют базис пространства $T^p_q(V)$.
\begin{knowledgefix}
В тензорном базисе используем \(e^{i_1}\otimes\cdots\otimes e^{i_p}\otimes e_{j_1}\otimes\cdots\otimes e_{j_q}\). Здесь \(e^i\) — элементы сопряжённого базиса, а \(e_j\) — векторы исходного базиса, рассматриваемые как элементы \(V^{**}\).
\end{knowledgefix}
\end{theorem}

Следовательно,
\[
    \dim T^p_q(V)=n^{p+q}.
\]
Любой тензор $T\in T^p_q(V)$ единственным образом раскладывается по этому базису:
\[
    T=
    \sum
    T^{i_1\ldots i_p}_{j_1\ldots j_q}
    e_{i_1}\otimes\cdots\otimes e_{i_p}
    \otimes
    \varepsilon^{j_1}\otimes\cdots\otimes\varepsilon^{j_q}.
\]
Коэффициенты
\[
    T^{i_1\ldots i_p}_{j_1\ldots j_q}
\]
называются компонентами тензора в выбранном базисе.

\begin{proof}[Идея доказательства]
Доказательство сводится к проверке линейной независимости и порождённости. Если линейная комбинация базисных тензоров равна нулю, то подставляя в неё соответствующие базисные векторы и ковекторы, получаем, что каждый коэффициент равен нулю.
\begin{knowledgefix}
На базисных аргументах базисный тензор даёт произведение символов Кронекера: \((e^{i}\otimes e^{j})(e_k,e_l)=\delta^i_k\delta^j_l\), а в общем случае получается произведение всех соответствующих \(\delta\).
\end{knowledgefix}
\end{proof}

\subsubsection*{Замена базиса и правило Эйнштейна}

Пусть выполняется замена базиса
\[
    \tilde e_i=\sum_j c_i^j e_j.
\]
Тогда двойственный базис преобразуется контраградиентно:
\[
    \tilde\varepsilon^i=\sum_j d^i_j\varepsilon^j,
\]
где матрица $(d^i_j)$ обратна к матрице перехода в подходящем смысле:
\[
    D=C^{-1}
    \quad \text{или} \quad
    D=(C^{-1})^T,
\]
в зависимости от того, как записывать базисные векторы столбцами/строками.
\begin{knowledgefix}
Фиксируем convention: если \(\tilde e_i=\tau^k_i e_k\), то сопряжённый базис меняется обратной матрицей \(\tilde e^{j}=\sigma^j_l e^l\), где \(\sigma=\tau^{-1}\).
\end{knowledgefix}

Компоненты тензора при замене базиса преобразуются с матрицами перехода по каждому индексу. Схематически:
\[
    \tilde T^{i_1\ldots i_p}_{j_1\ldots j_q}
    =
    \sum
    c^{i_1}_{a_1}\cdots c^{i_p}_{a_p}
    d^{b_1}_{j_1}\cdots d^{b_q}_{j_q}
    T^{a_1\ldots a_p}_{b_1\ldots b_q}.
\]
\begin{knowledgefix}
Общий закон смены координат:
\[\widetilde T^{j_1\ldots j_q}_{i_1\ldots i_p}=\tau^{k_1}_{i_1}\cdots\tau^{k_p}_{i_p}\sigma^{j_1}_{l_1}\cdots\sigma^{j_q}_{l_q}T^{l_1\ldots l_q}_{k_1\ldots k_p},\]
где \(\tau\) — матрица перехода для векторного базиса, \(\sigma=\tau^{-1}\) — для двойственного.
\end{knowledgefix}

\begin{remark}[Правило Эйнштейна]
Если индекс в одном мономе встречается дважды --- один раз сверху и один раз снизу, --- по нему подразумевается суммирование:
\[
    a_i b^i := \sum_i a_i b^i.
\]
Например,
\[
    x=x^i e_i.
\]
\begin{knowledgefix}
Правило Эйнштейна: если индекс повторяется один раз сверху и один раз снизу, по нему суммируют. Индексы, которые не суммируются, называются свободными и должны совпадать в левой и правой частях формулы.
\end{knowledgefix}
\end{remark}

\subsubsection*{Естественные операции и примеры тензоров}

На странице 80 продолжается обсуждение тензорных операций и примеров.
Видны записи, связанные с билинейной формой и линейным оператором. Для билинейной формы $\beta$ в базисе
\[
    \beta(x,y)=b_{ij}x^i y^j,
\]
то есть её компоненты образуют матрицу
\[
    B=(b_{ij}).
\]

Линейный оператор $A:V\to V$ имеет компоненты
\[
    A(e_j)=a^i_j e_i,
\]
и поэтому задаёт тензор типа $(1,1)$:
\[
    A=A^i_j e_i\otimes\varepsilon^j.
\]

Свёртка тензора типа $(1,1)$ даёт след оператора:
\[
    \operatorname{contr}(A)=A^i_i=\tr A.
\]
\begin{knowledgefix}
Следующий смысловой блок — таблица малых типов: ковекторы \(T^1_0\), векторы \(T^0_1\), билинейные формы \(T^2_0\), операторы \(T^1_1\), скаляры \(T^0_0\).
\end{knowledgefix}


% ============================================================
% Блок страниц 081--101: Тензоры, свёртка, симметризация, внешняя алгебра
% ============================================================

\subsubsection{Тензоры: продолжение}

\subsubsection*{Пример тензора и определение валентности}

На странице продолжается пример тензора, построенного из оператора и скалярного произведения. В координатах относительно базиса появляются выражения вида
\[
    T(e_i,e^j)=a_i^j,
\]
где $a_i^j$ --- компоненты соответствующего оператора/тензора.
\begin{knowledgefix}
Для оператора \(A\) как тензора типа \((1,1)\): \(A=A^i_j e^j\otimes e_i\), а действие на \(x=x^j e_j\) даёт \((Ax)^i=A^i_jx^j\).
\end{knowledgefix}

\begin{definition}[Тензор валентности $(p,q)$]
Пусть $V$ --- конечномерное линейное пространство. Тензором валентности $(p,q)$ называется полилинейная функция
\[
    T\colon \underbrace{V\times\cdots\times V}_{p\text{ раз}}
       \times
       \underbrace{V^*\times\cdots\times V^*}_{q\text{ раз}}
       \longrightarrow \F.
\]
В координатах значение такого тензора выражается как многочлен, линейный по каждой группе координат аргументов.
\begin{knowledgefix}
Порядок аргументов фиксируем как \(T(x_1,\ldots,x_p,\omega^1,\ldots,\omega^q)\), где сначала идут векторы из \(V\), затем ковекторы из \(V^*\). Все координатные формулы ниже согласованы с этим порядком.
\end{knowledgefix}
\end{definition}

Множество всех тензоров валентности $(p,q)$ обозначается в конспекте через
\[
    \Tens{p}{q}.
\]

\begin{remark}
Фраза на странице: тензоры фиксированной валентности образуют линейное пространство. Сложение и умножение на скаляр определяются поточечно.
\end{remark}

\subsubsection*{Тензоры типа $(1,1)$ и операторы}

Для линейного оператора $A\in \End V$ можно построить тензор типа $(1,1)$:
\[
    T_A(x,\omega)=\omega(Ax),
\]
или, при наличии скалярного произведения,
\[
    T_A(x,\omega)=\langle \omega,Ax\rangle.
\]
\begin{knowledgefix}
Отождествление \(V\simeq V^{**}\) задаётся канонически: вектор \(v\in V\) соответствует функционалу на \(V^*\), который отправляет \(\omega\) в \(\omega(v)\). Поэтому оператор можно понимать как тензор типа \((1,1)\).
\end{knowledgefix}

Если $e_1,\ldots,e_n$ --- базис $V$, а $e^1,\ldots,e^n$ --- двойственный базис $V^*$, то компоненты тензора $T_A$ имеют вид
\[
    T_A(e_j,e^i)=e^i(Ae_j)=a^i_j.
\]
Следовательно, матрица тензора типа $(1,1)$ совпадает с матрицей оператора $A$ в соответствующих базисах.

\begin{remark}
На стр. 82 сделан вывод: тензоры типа $(1,1)$ можно отождествлять с линейными операторами:
\[
    \Tens{1}{1}\simeq \End(V).
\]
\begin{knowledgefix}
Обозначение пространства тензоров может быть \(T^1_1(V)\) или \(\Pi^1_1(V)\); в этом файле используем единое \(T^p_q(V)\) / \(\Tens{p}{q}\) как пространство тензоров типа \((p,q)\).
\end{knowledgefix}
\end{remark}

\subsubsection{Операции над тензорами}

\subsubsection*{Свёртка тензора}

Пусть $T\in\Tens{p}{q}$. Если у тензора есть хотя бы один нижний и один верхний индекс, можно выполнить свёртку по выбранной паре индексов. В координатной форме это соответствует суммированию по повторяющемуся индексу:
\[
    (C^i_j T)^{\alpha_1\ldots \widehat{\alpha_i}\ldots \alpha_q}_{\beta_1\ldots \widehat{\beta_j}\ldots \beta_p}
    =
    \sum_{s=1}^{n}
    T^{\alpha_1\ldots s\ldots \alpha_q}_{\beta_1\ldots s\ldots \beta_p}.
\]
\begin{knowledgefix}
Если \(\tilde e_i=C^j_i e_j\), то координаты вектора меняются обратной матрицей: \(\tilde x^i=(C^{-1})^i_jx^j\). Для ковекторов возникает транспонированно-обратное преобразование. Поэтому верхние и нижние индексы меняются взаимно обратными матрицами.
\end{knowledgefix}

\begin{proposition}
Свёртка не зависит от выбора базиса.
\end{proposition}

\begin{proof}[Идея доказательства по записи страницы]
В конспекте доказывается, что если
\[
    \tilde e_i=\sum_j c_i^j e_j,
    \qquad
    \tilde e^i=\sum_k \tilde c^i_k e^k,
\]
то при подстановке в определение свёртки коэффициенты перехода взаимно уничтожаются благодаря соотношению двойственных базисов
\[
    \sum_i c_i^j\tilde c^i_k=\delta^j_k.
\]
После суммирования получается то же значение свёртки.
\begin{knowledgefix}
При смене базиса для типа \((1,1)\) матрица компонет меняется по закону подобия: \(\widetilde A=\tau^{-1}A\tau\) или \(\widetilde A=\tau A\tau^{-1}\) в зависимости от того, записываются ли координаты столбцами старого базиса или наоборот; в этом файле используется формула через \(\tau\) и \(\sigma=\tau^{-1}\) выше.
\end{knowledgefix}
\end{proof}

\subsubsection*{Тензорное произведение}

\begin{definition}[Тензорное произведение тензоров]
Пусть
\[
    T\in \Tens{p}{q},\qquad S\in \Pi^{s}_{r}(V).
\]
Их тензорным произведением называется тензор
\[
    T\otimes S\in \Pi^{q+s}_{p+r}(V),
\]
который на аргументах действует как произведение значений:
\[
\begin{aligned}
    &(T\otimes S)(x_1,\ldots,x_{p+r};\omega^1,\ldots,\omega^{q+s}) \\
    &\quad =
    T(x_1,\ldots,x_p;\omega^1,\ldots,\omega^q)
    S(x_{p+1},\ldots,x_{p+r};\omega^{q+1},\ldots,\omega^{q+s}).
\end{aligned}
\]
\begin{knowledgefix}
Порядок аргументов фиксируем как \(T(x_1,\ldots,x_p,\omega^1,\ldots,\omega^q)\), где сначала идут векторы из \(V\), затем ковекторы из \(V^*\). Все координатные формулы ниже согласованы с этим порядком.
\end{knowledgefix}
\end{definition}

В координатах компоненты тензорного произведения равны произведениям компонент:
\[
    (T\otimes S)^{i_1\ldots i_q j_1\ldots j_s}_{k_1\ldots k_p l_1\ldots l_r}
    =
    T^{i_1\ldots i_q}_{k_1\ldots k_p}
    S^{j_1\ldots j_s}_{l_1\ldots l_r}.
\]

\begin{example}
Если $\omega\in V^*$ и $x\in V$, то
\[
    \omega\otimes x\in \Tens{1}{1}.
\]
Если $A\in\End(V)$ и $\beta\in B^2(V)$ --- билинейная форма, то в конспекте рассматривается тензор вида
\[
    A\otimes \beta.
\]
Его значение получается как произведение значения оператораного тензора и билинейной формы.
\begin{knowledgefix}
Базовые соответствия: \(\omega\in V^*=T^1_0(V)\), \(v\in V\simeq T^0_1(V)\), \(A\in\End(V)\simeq T^1_1(V)\), \(\beta\in\operatorname{Bil}(V)=T^2_0(V)\).
\end{knowledgefix}
\end{example}

\begin{example}
В конспекте приведён вычислительный пример для $\dim V=2$, где относительно базиса $e_1,e_2$ заданы матрицы операторов $A$ и $B$, затем вычисляется компонент тензора $A\otimes B$ на базисных элементах вида
\[
    e_i\otimes e_j.
\]
Идея примера:
\[
    (A\otimes B)(e_i\otimes e_j)=A(e_i)\otimes B(e_j).
\]
\begin{knowledgefix}
Численный пример с матрицами заменён на проверенный общий шаблон: для тензора типа \((2,2)\) значение на \((x_1,x_2,\omega^1,\omega^2)\) равно \(T^{ab}_{ij}x_1^i x_2^j\omega^1_a\omega^2_b\), где суммирование идёт по всем повторяющимся индексам.
\end{knowledgefix}
\end{example}

\subsection*{Дополнительные уточнения}

\subsubsection*{7. Тензоры типа $(p,q)$, базис, компоненты и правило Эйнштейна}
В используемой convention тензор типа $(p,q)$ --- это полилинейная функция
\[
    T: \underbrace{V\times\cdots\times V}_{p\text{ раз}}
       \times
       \underbrace{V^*\times\cdots\times V^*}_{q\text{ раз}}
       \to \F.
\]
Иначе говоря, $p$ --- число векторных аргументов, $q$ --- число ковекторных аргументов.
Тензоры фиксированного типа образуют векторное пространство $T^p_q(V)$.

Простые типы:
\[
    T^0_0(V)=\F,
    \qquad
    T^1_0(V)=V^*,
    \qquad
    T^0_1(V)\simeq V,
\]
\[
    T^2_0(V)=\operatorname{Bil}(V),
    \qquad
    T^1_1(V)\simeq \End(V).
\]

Если $e_1,\ldots,e_n$ --- базис $V$, а $e^1,\ldots,e^n$ --- сопряжённый базис $V^*$, то базис пространства $T^p_q(V)$ состоит из тензоров
\[
    e^{i_1}\otimes\cdots\otimes e^{i_p}\otimes e_{j_1}\otimes\cdots\otimes e_{j_q}.
\]
Разложение тензора имеет вид
\[
    T=T^{j_1\ldots j_q}_{i_1\ldots i_p}
      e^{i_1}\otimes\cdots\otimes e^{i_p}\otimes e_{j_1}\otimes\cdots\otimes e_{j_q}.
\]
Размерность:
\[
    \dim T^p_q(V)=n^{p+q}.
\]

Правило Эйнштейна: если индекс встречается один раз сверху и один раз снизу, то по нему подразумевается суммирование.
Например,
\[
    x=x^i e_i,
    \qquad
    (Ax)^i=A^i_j x^j.
\]
Свободные индексы остаются в результате, связанные индексы суммируются.

\subsubsection*{9. Смена базиса и закон преобразования компонент}
Пусть новый базис выражен через старый как
\[
    \tilde e_j=e_i C^i_j.
\]
Тогда координаты вектора преобразуются обратной матрицей:
\[
    x^i=C^i_j\tilde x^j,
    \qquad
    \tilde x^j=(C^{-1})^j_i x^i.
\]
Сопряжённый базис преобразуется контраградиентно:
\[
    \tilde e^i=(C^{-1})^i_j e^j.
\]

Для тензора типа $(1,1)$ закон преобразования имеет матричный вид
\[
    \tilde A=C^{-1}AC.
\]
Для билинейной формы, то есть тензора типа $(2,0)$,
\[
    \tilde B=C^TBC
\]
в матричной записи, если столбцы $C$ задают новый базис через старый.

Общий закон: верхние индексы преобразуются как координаты векторов, нижние --- как координаты ковекторов, то есть с обратной/прямой матрицей в зависимости от выбранной convention записи базиса. Поэтому при verified-проходе важно сохранить ровно convention лектора.

\clearpage

\section{Билет 24. Свёртка тензоров. Интерпретация операций линейной алгебры в терминах тензорного умножения и свёртки. Тензорное произведение операторов. Метрический тензор, подъём и опускание индексов.}

\paragraph{Содержание билета.} Свёртка тензоров. Интерпретация операций линейной алгебры в терминах тензорного умножения и свёртки. Тензорное произведение операторов. Метрический тензор, подъём и опускание индексов.

\begin{planadd}

свёртка уменьшает один верхний и один нижний индекс; след оператора и действие оператора на вектор — базовые примеры; метрический тензор позволяет поднимать и опускать индексы.

\end{planadd}
\VisualContraction


\subsubsection{Тензоры: продолжение}

\subsubsection*{Пример тензора и определение валентности}

На странице продолжается пример тензора, построенного из оператора и скалярного произведения. В координатах относительно базиса появляются выражения вида
\[
    T(e_i,e^j)=a_i^j,
\]
где $a_i^j$ --- компоненты соответствующего оператора/тензора.
\begin{knowledgefix}
Для оператора \(A\) как тензора типа \((1,1)\): \(A=A^i_j e^j\otimes e_i\), а действие на \(x=x^j e_j\) даёт \((Ax)^i=A^i_jx^j\).
\end{knowledgefix}

\begin{definition}[Тензор валентности $(p,q)$]
Пусть $V$ --- конечномерное линейное пространство. Тензором валентности $(p,q)$ называется полилинейная функция
\[
    T\colon \underbrace{V\times\cdots\times V}_{p\text{ раз}}
       \times
       \underbrace{V^*\times\cdots\times V^*}_{q\text{ раз}}
       \longrightarrow \F.
\]
В координатах значение такого тензора выражается как многочлен, линейный по каждой группе координат аргументов.
\begin{knowledgefix}
Порядок аргументов фиксируем как \(T(x_1,\ldots,x_p,\omega^1,\ldots,\omega^q)\), где сначала идут векторы из \(V\), затем ковекторы из \(V^*\). Все координатные формулы ниже согласованы с этим порядком.
\end{knowledgefix}
\end{definition}

Множество всех тензоров валентности $(p,q)$ обозначается в конспекте через
\[
    \Tens{p}{q}.
\]

\begin{remark}
Фраза на странице: тензоры фиксированной валентности образуют линейное пространство. Сложение и умножение на скаляр определяются поточечно.
\end{remark}

\subsubsection*{Тензоры типа $(1,1)$ и операторы}

Для линейного оператора $A\in \End V$ можно построить тензор типа $(1,1)$:
\[
    T_A(x,\omega)=\omega(Ax),
\]
или, при наличии скалярного произведения,
\[
    T_A(x,\omega)=\langle \omega,Ax\rangle.
\]
\begin{knowledgefix}
Отождествление \(V\simeq V^{**}\) задаётся канонически: вектор \(v\in V\) соответствует функционалу на \(V^*\), который отправляет \(\omega\) в \(\omega(v)\). Поэтому оператор можно понимать как тензор типа \((1,1)\).
\end{knowledgefix}

Если $e_1,\ldots,e_n$ --- базис $V$, а $e^1,\ldots,e^n$ --- двойственный базис $V^*$, то компоненты тензора $T_A$ имеют вид
\[
    T_A(e_j,e^i)=e^i(Ae_j)=a^i_j.
\]
Следовательно, матрица тензора типа $(1,1)$ совпадает с матрицей оператора $A$ в соответствующих базисах.

\begin{remark}
На стр. 82 сделан вывод: тензоры типа $(1,1)$ можно отождествлять с линейными операторами:
\[
    \Tens{1}{1}\simeq \End(V).
\]
\begin{knowledgefix}
Обозначение пространства тензоров может быть \(T^1_1(V)\) или \(\Pi^1_1(V)\); в этом файле используем единое \(T^p_q(V)\) / \(\Tens{p}{q}\) как пространство тензоров типа \((p,q)\).
\end{knowledgefix}
\end{remark}

\subsubsection{Операции над тензорами}

\subsubsection*{Свёртка тензора}

Пусть $T\in\Tens{p}{q}$. Если у тензора есть хотя бы один нижний и один верхний индекс, можно выполнить свёртку по выбранной паре индексов. В координатной форме это соответствует суммированию по повторяющемуся индексу:
\[
    (C^i_j T)^{\alpha_1\ldots \widehat{\alpha_i}\ldots \alpha_q}_{\beta_1\ldots \widehat{\beta_j}\ldots \beta_p}
    =
    \sum_{s=1}^{n}
    T^{\alpha_1\ldots s\ldots \alpha_q}_{\beta_1\ldots s\ldots \beta_p}.
\]
\begin{knowledgefix}
Если \(\tilde e_i=C^j_i e_j\), то координаты вектора меняются обратной матрицей: \(\tilde x^i=(C^{-1})^i_jx^j\). Для ковекторов возникает транспонированно-обратное преобразование. Поэтому верхние и нижние индексы меняются взаимно обратными матрицами.
\end{knowledgefix}

\begin{proposition}
Свёртка не зависит от выбора базиса.
\end{proposition}

\begin{proof}[Идея доказательства по записи страницы]
В конспекте доказывается, что если
\[
    \tilde e_i=\sum_j c_i^j e_j,
    \qquad
    \tilde e^i=\sum_k \tilde c^i_k e^k,
\]
то при подстановке в определение свёртки коэффициенты перехода взаимно уничтожаются благодаря соотношению двойственных базисов
\[
    \sum_i c_i^j\tilde c^i_k=\delta^j_k.
\]
После суммирования получается то же значение свёртки.
\begin{knowledgefix}
При смене базиса для типа \((1,1)\) матрица компонет меняется по закону подобия: \(\widetilde A=\tau^{-1}A\tau\) или \(\widetilde A=\tau A\tau^{-1}\) в зависимости от того, записываются ли координаты столбцами старого базиса или наоборот; в этом файле используется формула через \(\tau\) и \(\sigma=\tau^{-1}\) выше.
\end{knowledgefix}
\end{proof}

\subsubsection*{Тензорное произведение}

\begin{definition}[Тензорное произведение тензоров]
Пусть
\[
    T\in \Tens{p}{q},\qquad S\in \Pi^{s}_{r}(V).
\]
Их тензорным произведением называется тензор
\[
    T\otimes S\in \Pi^{q+s}_{p+r}(V),
\]
который на аргументах действует как произведение значений:
\[
\begin{aligned}
    &(T\otimes S)(x_1,\ldots,x_{p+r};\omega^1,\ldots,\omega^{q+s}) \\
    &\quad =
    T(x_1,\ldots,x_p;\omega^1,\ldots,\omega^q)
    S(x_{p+1},\ldots,x_{p+r};\omega^{q+1},\ldots,\omega^{q+s}).
\end{aligned}
\]
\begin{knowledgefix}
Порядок аргументов фиксируем как \(T(x_1,\ldots,x_p,\omega^1,\ldots,\omega^q)\), где сначала идут векторы из \(V\), затем ковекторы из \(V^*\). Все координатные формулы ниже согласованы с этим порядком.
\end{knowledgefix}
\end{definition}

В координатах компоненты тензорного произведения равны произведениям компонент:
\[
    (T\otimes S)^{i_1\ldots i_q j_1\ldots j_s}_{k_1\ldots k_p l_1\ldots l_r}
    =
    T^{i_1\ldots i_q}_{k_1\ldots k_p}
    S^{j_1\ldots j_s}_{l_1\ldots l_r}.
\]

\begin{example}
Если $\omega\in V^*$ и $x\in V$, то
\[
    \omega\otimes x\in \Tens{1}{1}.
\]
Если $A\in\End(V)$ и $\beta\in B^2(V)$ --- билинейная форма, то в конспекте рассматривается тензор вида
\[
    A\otimes \beta.
\]
Его значение получается как произведение значения оператораного тензора и билинейной формы.
\begin{knowledgefix}
Базовые соответствия: \(\omega\in V^*=T^1_0(V)\), \(v\in V\simeq T^0_1(V)\), \(A\in\End(V)\simeq T^1_1(V)\), \(\beta\in\operatorname{Bil}(V)=T^2_0(V)\).
\end{knowledgefix}
\end{example}

\begin{example}
В конспекте приведён вычислительный пример для $\dim V=2$, где относительно базиса $e_1,e_2$ заданы матрицы операторов $A$ и $B$, затем вычисляется компонент тензора $A\otimes B$ на базисных элементах вида
\[
    e_i\otimes e_j.
\]
Идея примера:
\[
    (A\otimes B)(e_i\otimes e_j)=A(e_i)\otimes B(e_j).
\]
\begin{knowledgefix}
Численный пример с матрицами заменён на проверенный общий шаблон: для тензора типа \((2,2)\) значение на \((x_1,x_2,\omega^1,\omega^2)\) равно \(T^{ab}_{ij}x_1^i x_2^j\omega^1_a\omega^2_b\), где суммирование идёт по всем повторяющимся индексам.
\end{knowledgefix}
\end{example}

\subsubsection{Опускание и поднятие индексов}

\subsubsection*{Невырожденная билинейная форма как изоморфизм}

Пусть $V$ --- конечномерное пространство и $g$ --- невырожденная билинейная форма, например скалярное произведение.
Тогда каждому $y\in V$ сопоставляется функционал
\[
    y^*\colon V\to\F,
    \qquad
    y^*(x)=g(x,y).
\]
Получаем линейный изоморфизм
\[
    V\longrightarrow V^*,
    \qquad
    y\longmapsto y^*.
\]

\begin{remark}[Опускание индекса]
В координатах, если
\[
    g=g_{ij}e^i\otimes e^j,
\]
то вектор $y=y^j e_j$ переходит в ковектор
\[
    y_i=g_{ij}y^j.
\]
Это называется опусканием индекса.
\end{remark}

\begin{remark}[Поднятие индекса]
Если матрица $(g_{ij})$ невырождена, то существует обратная матрица $(g^{ij})$, и можно выполнять обратную операцию:
\[
    y^i=g^{ij}y_j.
\]
Это называется поднятием индекса.
\begin{knowledgefix}
Метрический тензор имеет компоненты \(g_{ij}=(e_i,e_j)\). Он опускает индекс: \(v_i=g_{ij}v^j\). Обратный метрический тензор \(g^{ij}\) поднимает индекс: \(\alpha^i=g^{ij}\alpha_j\).
\end{knowledgefix}
\end{remark}

\subsubsection*{Скалярное произведение тензоров}

При наличии скалярного произведения на $V$ можно ввести скалярное произведение на пространствах тензоров. Для базисных тензоров оно задаётся через произведение соответствующих компонент, а затем продолжается билинейно.
В координатах появляются формулы вида
\[
    \langle T,S\rangle
    =
    \sum T^{i_1\ldots i_q}_{j_1\ldots j_p}
          S^{j_1\ldots j_p}_{i_1\ldots i_q},
\]
после поднятия/опускания индексов с помощью $g$.
\begin{knowledgefix}
Формулы поднятия и опускания индексов: \(v_i=g_{ij}v^j\), \(\alpha^i=g^{ij}\alpha_j\), причём \(g^{ik}g_{kj}=\delta^i_j\). Для общего тензора операция применяется к выбранному индексу свёрткой с \(g\) или \(g^{-1}\).
\end{knowledgefix}

\subsection*{Дополнительные уточнения}

\subsubsection*{8. Тензорное произведение и свёртка}
Если
\[
    S\in T^p_q(V),
    \qquad
    T\in T^r_s(V),
\]
то
\[
    S\otimes T\in T^{p+r}_{q+s}(V).
\]
На аргументах это означает: первые аргументы подаются в $S$, оставшиеся --- в $T$, а результаты перемножаются.
Тензорное произведение ассоциативно, билинейно и в общем случае некоммутативно.

Свёртка тензора связывает один верхний и один нижний индекс.
Для разложимого тензора
\[
    \omega^1\otimes\cdots\otimes\omega^p\otimes x_1\otimes\cdots\otimes x_q
\]
свёртка по паре $(k,l)$ имеет вид
\[
    C^k_l(\omega^1\otimes\cdots\otimes\omega^p\otimes x_1\otimes\cdots\otimes x_q)
    =\omega^k(x_l)\,
      \omega^1\otimes\cdots\widehat{\omega^k}\cdots\otimes\omega^p
      \otimes
      x_1\otimes\cdots\widehat{x_l}\cdots\otimes x_q.
\]
В координатах свёртка означает суммирование по совпавшей паре верхнего и нижнего индекса.
Для тензора типа $(1,1)$ это даёт след:
\[
    C(T)=T^i_i=\tr T.
\]

Гиперстрока/гиперстолбец в многомерной записи тензора --- это фиксация всех индексов, кроме одного. Такая запись помогает видеть тензор как многомерный массив, но сам тензор не равен массиву без закона преобразования координат.

\subsubsection*{10. Метрический тензор, подъём и опускание индексов}

Метрический тензор задаётся скалярным произведением:
\[
    g(x,y)=(x,y).
\]
В базисе его компоненты
\[
    g_{ij}=(e_i,e_j).
\]
Если $x=x^i e_i$, то соответствующий ковектор $x^\flat$ имеет координаты
\[
    x_i=g_{ij}x^j.
\]
Это называется опусканием индекса.
Обратный метрический тензор имеет компоненты $g^{ij}$, где
\[
    g^{ik}g_{kj}=\delta^i_j.
\]
Подъём индекса:
\[
    x^i=g^{ij}x_j.
\]
Для общего тензора опускание верхнего индекса и подъём нижнего индекса выполняются свёрткой с $g_{ij}$ или $g^{ij}$ соответственно.

\clearpage

\section{Билет 25. Симметрические и кососимметрические тензоры. Симметризация и альтернирование (антисимметризация), их свойства.}

\paragraph{Содержание билета.} Симметрические и кососимметрические тензоры. Симметризация и альтернирование (антисимметризация), их свойства.

\begin{planadd}

симметризация — усреднение по перестановкам, альтернирование — усреднение со знаком перестановки. Отметь, что для $p\ge3$ всё пространство не обязано распадаться только на симметрическую и кососимметрическую части.

\end{planadd}
\VisualSymAlt


\subsubsection{Симметрические тензоры}

\subsubsection*{Действие группы перестановок и симметризация}

Пусть $T\in\Pi^0_p(V)$ --- ковариантный тензор порядка $p$.
Для перестановки $\sigma\in S_p$ определяется оператор $f_\sigma$ перестановки аргументов:
\[
    (f_\sigma T)(x_1,\ldots,x_p)
    =
    T(x_{\sigma(1)},\ldots,x_{\sigma(p)}).
\]
\begin{knowledgefix}
Оператор перестановки аргументов можно определить как \((F_\sigma T)(x_1,
\ldots,x_p)=T(x_{\sigma(1)},\ldots,x_{\sigma(p)})\) или через \(\sigma^{-1}\). Оба соглашения дают одинаковые операторы \(\Sym\) и \(\Alt\) после суммирования по всей группе \(S_p\).
\end{knowledgefix}

\begin{definition}[Симметрический тензор]
Тензор $T\in\Pi^0_p(V)$ называется симметрическим, если
\[
    f_\sigma T=T
    \qquad
    \forall\sigma\in S_p.
\]
Пространство симметрических тензоров обозначается
\[
    \Sym_p(V)\subset \Pi^0_p(V).
\]
\end{definition}

\begin{definition}[Оператор симметризации]
Оператор симметризации задаётся формулой
\[
    \Sym(T)=\frac{1}{p!}\sum_{\sigma\in S_p} f_\sigma(T).
\]
\end{definition}

\begin{proposition}[Свойства симметризации]
Для оператора $\Sym$ выполнены:
\begin{enumerate}[label=\arabic*)]
    \item $\Imop(\Sym)=\Sym_p(V)$;
    \item $\Ker(\Sym)$ состоит из тензоров, симметризация которых равна нулю;
    \item $\Sym^2=\Sym$.
\end{enumerate}
\end{proposition}

\begin{proof}[Набросок доказательства по страницам 88--90]
Если $T$ уже симметричен, то для всех $\sigma$ имеем $f_\sigma T=T$, поэтому
\[
    \Sym(T)=\frac{1}{p!}\sum_{\sigma\in S_p}T=T.
\]
Отсюда $\Sym^2=\Sym$ и $\Imop(\Sym)=\Sym_p(V)$.
Обратное включение получается из того, что симметризация любой формы инвариантна относительно перестановок аргументов.
\begin{knowledgefix}
Симметризация действительно даёт симметрический тензор, потому что для любой \(\tau\in S_p\): \(F_\tau\Sym T=\frac1{p!}\sum_{\sigma\in S_p}F_{\tau\sigma}T=\Sym T\) после перенумерации \(\rho=\tau\sigma\).
\end{knowledgefix}
\end{proof}

\begin{example}[Симметризация монома, стр. 89--90]
Для базисных ковекторов $e^i$ рассматриваются тензоры вида
\[
    e^{i_1}\otimes\cdots\otimes e^{i_p}.
\]
Их симметризация:
\[
    \Sym(e^{i_1}\otimes\cdots\otimes e^{i_p})
    =
    \frac{1}{p!}\sum_{\sigma\in S_p}
    e^{i_{\sigma(1)}}\otimes\cdots\otimes e^{i_{\sigma(p)}}.
\]
Если среди индексов есть повторения, одинаковые слагаемые группируются.
\begin{knowledgefix}
Для \(n=2,p=3\) базис симметрических ковариантных тензоров соответствует мономам степеней \(3\): \((e^1)^3,(e^1)^2e^2,e^1(e^2)^2,(e^2)^3\). Поэтому координаты часто обозначают \(\lambda_{30},\lambda_{21},\lambda_{12},\lambda_{03}\).
\end{knowledgefix}
\end{example}

\begin{lemma}[Базис пространства симметрических тензоров]
Симметрические мономы
\[
    \Sym\bigl((e^1)^{\otimes a_1}\otimes\cdots\otimes(e^n)^{\otimes a_n}\bigr),
    \qquad
    a_1+\cdots+a_n=p,
\]
образуют базис пространства $\Sym_p(V)$.
\end{lemma}

\begin{corollary}[Размерность]
\[
    \dim \Sym_p(V)=\binom{n+p-1}{p}.
\]
\begin{knowledgefix}
Размерность пространства симметрических \(p\)-тензоров на \(n\)-мерном пространстве равна числу мономов степени \(p\) от \(n\) переменных: \(\binom{n+p-1}{p}\).
\end{knowledgefix}
\end{corollary}

\subsubsection{Альтернирование и внешние формы}

\subsubsection*{Оператор альтернирования}

\begin{definition}[Оператор альтернирования]
Для $T\in\Pi^0_p(V)$ оператор альтернирования задаётся формулой
\[
    \Alt(T)=\frac{1}{p!}\sum_{\sigma\in S_p}\sgn(\sigma)f_\sigma(T).
\]
\end{definition}

\begin{example}
Для $p=2$:
\[
    \Alt(e^1\otimes e^2)
    =
    \frac{1}{2}(e^1\otimes e^2-e^2\otimes e^1).
\]
В конспекте рядом записано сравнение с симметризацией:
\[
    \Sym(e^1\otimes e^2)
    =
    \frac{1}{2}(e^1\otimes e^2+e^2\otimes e^1).
\]
\end{example}

\begin{lemma}
Для оператора $\Alt$ выполнены:
\begin{enumerate}[label=\arabic*)]
    \item $\Imop(\Alt)=\Lambda^p(V^*)$;
    \item $\Ker(\Alt)$ состоит из тензоров, альтернирование которых равно нулю;
    \item $\Alt^2=\Alt$.
\end{enumerate}
\end{lemma}

\begin{proof}[Идея доказательства]
Если $T$ кососимметричен, то
\[
    f_\sigma(T)=\sgn(\sigma)T,
\]
поэтому
\[
    \Alt(T)=\frac{1}{p!}\sum_{\sigma\in S_p}\sgn(\sigma)^2T=T.
\]
С другой стороны, альтернирование произвольного тензора даёт кососимметрический тензор.
\begin{knowledgefix}
Базис симметрических тензоров получают симметризацией \(\Sym(e^{i_1}\otimes\cdots\otimes e^{i_p})\) с неубывающими индексами \(i_1\le\cdots\le i_p\). Эти элементы порождают все симметрические тензоры и линейно независимы.
\end{knowledgefix}
\end{proof}

\begin{corollary}
Если среди индексов $i_1,\ldots,i_p$ есть совпадающие, то
\[
    \Alt(e^{i_1}\otimes\cdots\otimes e^{i_p})=0.
\]
\end{corollary}

\begin{definition}[Внешнее произведение базисных ковекторов]
Для попарно различных индексов вводится обозначение
\[
    e^{i_1}\wedge\cdots\wedge e^{i_p}
    =
    \Alt(e^{i_1}\otimes\cdots\otimes e^{i_p})
\]
или, с возможным нормировочным множителем,
\[
    e^{i_1}\wedge\cdots\wedge e^{i_p}
    =
    \sum_{\sigma\in S_p}\sgn(\sigma)
    e^{i_{\sigma(1)}}\otimes\cdots\otimes e^{i_{\sigma(p)}}.
\]
\begin{knowledgefix}
Фиксируем нормировку: \(\Alt T=\frac1{p!}\sum_{\sigma\in S_p}\sgn(\sigma)F_\sigma T\), \(\Sym T=\frac1{p!}\sum_{\sigma\in S_p}F_\sigma T\). Если у лектора внешнее произведение нормируется иначе, все формулы отличаются только общим фактором, но свойства сохраняются.
\end{knowledgefix}
\end{definition}

\begin{corollary}[Базис и размерность внешних форм]
Элементы
\[
    e^{i_1}\wedge\cdots\wedge e^{i_p},
    \qquad
    1\le i_1<\cdots<i_p\le n,
\]
образуют базис пространства $\Lambda^p(V^*)$.
Следовательно,
\[
    \dim \Lambda^p(V^*)=\binom{n}{p}.
\]
Если $p>n$, то
\[
    \Lambda^p(V^*)=0.
\]
\begin{knowledgefix}
Для \(p=2\) любой ковариантный тензор раскладывается на симметрическую и кососимметрическую части. Для \(p\ge3\) одних полностью симметрической и полностью кососимметрической частей недостаточно: существуют смешанные типы симметрии.
\end{knowledgefix}
\end{corollary}

\subsection*{Дополнительные уточнения}

\subsubsection*{11. Симметризация, альтернирование и симметрическое произведение}
Для чисто ковариантных $p$-тензоров $T\in T^p_0(V)$ и перестановки $\sigma\in S_p$ задаётся оператор перестановки аргументов
\[
    (F_\sigma T)(x_1,\ldots,x_p)=T(x_{\sigma(1)},\ldots,x_{\sigma(p)}).
\]
Оператор симметризации:
\[
    \Sym(T)=\frac1{p!}\sum_{\sigma\in S_p}F_\sigma T.
\]
Его образ --- пространство симметрических тензоров $S^p(V^*)$, причём
\[
    \Sym^2=\Sym.
\]
Базис $S^p(V^*)$ задаётся мономиальными симметризованными произведениями
\[
    (e^1)^{a_1}\cdots (e^n)^{a_n},
    \qquad
    a_1+\cdots+a_n=p,
\]
и
\[
    \dim S^p(V^*)=\binom{n+p-1}{p}.
\]

Оператор альтернирования:
\[
    \Alt(T)=\frac1{p!}\sum_{\sigma\in S_p}\sgn(\sigma)F_\sigma T.
\]
Кососимметрический тензор удовлетворяет
\[
    F_\sigma T=\sgn(\sigma)T.
\]
Внешнее произведение базисных ковекторов задаётся как
\[
    e^{i_1}\wedge\cdots\wedge e^{i_p}
    =\Alt(e^{i_1}\otimes\cdots\otimes e^{i_p}).
\]
Если среди индексов есть повторение, результат равен нулю. Базис $\Lambda^p(V^*)$ имеет вид
\[
    e^{i_1}\wedge\cdots\wedge e^{i_p},
    \qquad
    1\le i_1<\cdots<i_p\le n,
\]
поэтому
\[
    \dim\Lambda^p(V^*)=\binom np.
\]

Симметрическое умножение определяется формулой
\[
    T\odot S=\Sym(T\otimes S).
\]
Если $T\in S^p(V^*)$, $S\in S^r(V^*)$, то
\[
    T\odot S\in S^{p+r}(V^*).
\]
На аргументах:
\[
    (T\odot S)(x_1,\ldots,x_{p+r})
    =\frac1{(p+r)!}\sum_{\sigma\in S_{p+r}}
    T(x_{\sigma(1)},\ldots,x_{\sigma(p)})
    S(x_{\sigma(p+1)},\ldots,x_{\sigma(p+r)}).
\]
Симметрическое умножение билинейно, коммутативно и ассоциативно.

\clearpage

\section{Билет 26. Симметрическое умножение и его свойства. Симметрическая алгебра, её базис и размерность. Изоморфизм алгебре многочленов от нескольких переменных.}

\paragraph{Содержание билета.} Симметрическое умножение и его свойства. Симметрическая алгебра, её базис и размерность. Изоморфизм алгебре многочленов от нескольких переменных.

\begin{planadd}

Материал про полную симметрическую алгебру и изоморфизм с многочленами удобно оформить как стандартное дополнение к билету.

\end{planadd}

\subsubsection{Симметрические тензоры}

\subsubsection*{Действие группы перестановок и симметризация}

Пусть $T\in\Pi^0_p(V)$ --- ковариантный тензор порядка $p$.
Для перестановки $\sigma\in S_p$ определяется оператор $f_\sigma$ перестановки аргументов:
\[
    (f_\sigma T)(x_1,\ldots,x_p)
    =
    T(x_{\sigma(1)},\ldots,x_{\sigma(p)}).
\]
\begin{knowledgefix}
Оператор перестановки аргументов можно определить как \((F_\sigma T)(x_1,
\ldots,x_p)=T(x_{\sigma(1)},\ldots,x_{\sigma(p)})\) или через \(\sigma^{-1}\). Оба соглашения дают одинаковые операторы \(\Sym\) и \(\Alt\) после суммирования по всей группе \(S_p\).
\end{knowledgefix}

\begin{definition}[Симметрический тензор]
Тензор $T\in\Pi^0_p(V)$ называется симметрическим, если
\[
    f_\sigma T=T
    \qquad
    \forall\sigma\in S_p.
\]
Пространство симметрических тензоров обозначается
\[
    \Sym_p(V)\subset \Pi^0_p(V).
\]
\end{definition}

\begin{definition}[Оператор симметризации]
Оператор симметризации задаётся формулой
\[
    \Sym(T)=\frac{1}{p!}\sum_{\sigma\in S_p} f_\sigma(T).
\]
\end{definition}

\begin{proposition}[Свойства симметризации]
Для оператора $\Sym$ выполнены:
\begin{enumerate}[label=\arabic*)]
    \item $\Imop(\Sym)=\Sym_p(V)$;
    \item $\Ker(\Sym)$ состоит из тензоров, симметризация которых равна нулю;
    \item $\Sym^2=\Sym$.
\end{enumerate}
\end{proposition}

\begin{proof}[Набросок доказательства по страницам 88--90]
Если $T$ уже симметричен, то для всех $\sigma$ имеем $f_\sigma T=T$, поэтому
\[
    \Sym(T)=\frac{1}{p!}\sum_{\sigma\in S_p}T=T.
\]
Отсюда $\Sym^2=\Sym$ и $\Imop(\Sym)=\Sym_p(V)$.
Обратное включение получается из того, что симметризация любой формы инвариантна относительно перестановок аргументов.
\begin{knowledgefix}
Симметризация действительно даёт симметрический тензор, потому что для любой \(\tau\in S_p\): \(F_\tau\Sym T=\frac1{p!}\sum_{\sigma\in S_p}F_{\tau\sigma}T=\Sym T\) после перенумерации \(\rho=\tau\sigma\).
\end{knowledgefix}
\end{proof}

\begin{example}[Симметризация монома, стр. 89--90]
Для базисных ковекторов $e^i$ рассматриваются тензоры вида
\[
    e^{i_1}\otimes\cdots\otimes e^{i_p}.
\]
Их симметризация:
\[
    \Sym(e^{i_1}\otimes\cdots\otimes e^{i_p})
    =
    \frac{1}{p!}\sum_{\sigma\in S_p}
    e^{i_{\sigma(1)}}\otimes\cdots\otimes e^{i_{\sigma(p)}}.
\]
Если среди индексов есть повторения, одинаковые слагаемые группируются.
\begin{knowledgefix}
Для \(n=2,p=3\) базис симметрических ковариантных тензоров соответствует мономам степеней \(3\): \((e^1)^3,(e^1)^2e^2,e^1(e^2)^2,(e^2)^3\). Поэтому координаты часто обозначают \(\lambda_{30},\lambda_{21},\lambda_{12},\lambda_{03}\).
\end{knowledgefix}
\end{example}

\begin{lemma}[Базис пространства симметрических тензоров]
Симметрические мономы
\[
    \Sym\bigl((e^1)^{\otimes a_1}\otimes\cdots\otimes(e^n)^{\otimes a_n}\bigr),
    \qquad
    a_1+\cdots+a_n=p,
\]
образуют базис пространства $\Sym_p(V)$.
\end{lemma}

\begin{corollary}[Размерность]
\[
    \dim \Sym_p(V)=\binom{n+p-1}{p}.
\]
\begin{knowledgefix}
Размерность пространства симметрических \(p\)-тензоров на \(n\)-мерном пространстве равна числу мономов степени \(p\) от \(n\) переменных: \(\binom{n+p-1}{p}\).
\end{knowledgefix}
\end{corollary}

\subsubsection{Симметрическое произведение}

\subsubsection*{Определение и свойства}

\begin{definition}[Симметрическое произведение]
Пусть $T\in\Sym_p(V)$, $S\in\Sym_r(V)$. Их симметрическим произведением называется
\[
    T\odot S=\Sym(T\otimes S).
\]
Иногда в конспекте перед симметризацией появляется числовой множитель, зависящий от $p$ и $r$:
\[
    T\odot S=\frac{(p+r)!}{p!r!}\Sym(T\otimes S).
\]
\begin{knowledgefix}
Симметрическое произведение можно определить как \(S\odot T=\frac{(p+q)!}{p!q!}\Sym(S\otimes T)\) или просто \(\Sym(S\otimes T)\). В этом файле фиксируем первую нормировку, чтобы произведение базисных мономов совпадало с обычным умножением многочленов.
\end{knowledgefix}
\end{definition}

\begin{proposition}[Свойства симметрического произведения]
Для симметрических тензоров выполнены:
\begin{enumerate}[label=\arabic*)]
    \item билинейность:
    \[
        (T_1+T_2)\odot S=T_1\odot S+T_2\odot S,
        \qquad
        T\odot(S_1+S_2)=T\odot S_1+T\odot S_2;
    \]
    \item однородность по скаляру:
    \[
        (\lambda T)\odot S=\lambda(T\odot S)=T\odot(\lambda S);
    \]
    \item коммутативность:
    \[
        T\odot S=S\odot T;
    \]
    \item ассоциативность:
    \[
        (T\odot S)\odot R=T\odot(S\odot R).
    \]
\end{enumerate}
\end{proposition}

\begin{proof}[Идея доказательства по страницам 94--95]
Билинейность следует из билинейности тензорного произведения и линейности оператора $\Sym$.
Коммутативность доказывается тем, что перестановка блока первых $p$ аргументов с блоком следующих $r$ аргументов после полной симметризации не меняет результат.
Ассоциативность доказывается аналогично: полная симметризация тензора $T\otimes S\otimes R$ не зависит от порядка, в котором сначала симметризуются блоки.
\begin{knowledgefix}
Доказательство ассоциативности симметрического произведения сводится к тому, что повторная симметризация не меняет уже симметризованный результат: \(\Sym\circ\Sym=\Sym\), а суммы по перестановкам можно перегруппировать.
\end{knowledgefix}
\end{proof}

\subsection*{Дополнительные уточнения}

\subsubsection*{11. Симметризация, альтернирование и симметрическое произведение}
Для чисто ковариантных $p$-тензоров $T\in T^p_0(V)$ и перестановки $\sigma\in S_p$ задаётся оператор перестановки аргументов
\[
    (F_\sigma T)(x_1,\ldots,x_p)=T(x_{\sigma(1)},\ldots,x_{\sigma(p)}).
\]
Оператор симметризации:
\[
    \Sym(T)=\frac1{p!}\sum_{\sigma\in S_p}F_\sigma T.
\]
Его образ --- пространство симметрических тензоров $S^p(V^*)$, причём
\[
    \Sym^2=\Sym.
\]
Базис $S^p(V^*)$ задаётся мономиальными симметризованными произведениями
\[
    (e^1)^{a_1}\cdots (e^n)^{a_n},
    \qquad
    a_1+\cdots+a_n=p,
\]
и
\[
    \dim S^p(V^*)=\binom{n+p-1}{p}.
\]

Оператор альтернирования:
\[
    \Alt(T)=\frac1{p!}\sum_{\sigma\in S_p}\sgn(\sigma)F_\sigma T.
\]
Кососимметрический тензор удовлетворяет
\[
    F_\sigma T=\sgn(\sigma)T.
\]
Внешнее произведение базисных ковекторов задаётся как
\[
    e^{i_1}\wedge\cdots\wedge e^{i_p}
    =\Alt(e^{i_1}\otimes\cdots\otimes e^{i_p}).
\]
Если среди индексов есть повторение, результат равен нулю. Базис $\Lambda^p(V^*)$ имеет вид
\[
    e^{i_1}\wedge\cdots\wedge e^{i_p},
    \qquad
    1\le i_1<\cdots<i_p\le n,
\]
поэтому
\[
    \dim\Lambda^p(V^*)=\binom np.
\]

Симметрическое умножение определяется формулой
\[
    T\odot S=\Sym(T\otimes S).
\]
Если $T\in S^p(V^*)$, $S\in S^r(V^*)$, то
\[
    T\odot S\in S^{p+r}(V^*).
\]
На аргументах:
\[
    (T\odot S)(x_1,\ldots,x_{p+r})
    =\frac1{(p+r)!}\sum_{\sigma\in S_{p+r}}
    T(x_{\sigma(1)},\ldots,x_{\sigma(p)})
    S(x_{\sigma(p+1)},\ldots,x_{\sigma(p+r)}).
\]
Симметрическое умножение билинейно, коммутативно и ассоциативно.
\begin{planadd}
\textbf{Стандартное продолжение билета: эта часть нужна для полного ответа.}
Симметрическая алгебра пространства $V^*$ определяется как
\[
    S(V^*)=\bigoplus_{p=0}^{\infty} S^p(V^*),
\]
где $S^p(V^*)$ — пространство симметрических ковариантных $p$-тензоров.
Если $e^1,\ldots,e^n$ — базис $V^*$, то базис $S^p(V^*)$ задаётся мономами
\[
    (e^1)^{\odot k_1}\odot\cdots\odot(e^n)^{\odot k_n},
    \qquad k_1+\cdots+k_n=p.
\]
Отсюда
\[
    \dim S^p(V^*)=\binom{n+p-1}{p}.
\]
Изоморфизм с алгеброй многочленов: $e^i$ соответствует переменной $x_i$, а
симметрическое произведение соответствует обычному умножению многочленов.
\end{planadd}

\clearpage

\section{Билет 27. Внешнее умножение и его свойства. Грассманова (внешняя) алгебра, её базис и размерность. Изоморфизм алгебре грассмановых многочленов.}

\paragraph{Содержание билета.} Внешнее умножение и его свойства. Грассманова (внешняя) алгебра, её базис и размерность. Изоморфизм алгебре грассмановых многочленов.

\begin{planadd}

внешнее умножение строится через альтернирование тензорного произведения; базис задаётся $e^{i_1}\wedge\cdots\wedge e^{i_p}$ при $i_1<\cdots<i_p$, размерность $\binom np$.

\end{planadd}

\subsubsection{Альтернирование и внешние формы}

\subsubsection*{Оператор альтернирования}

\begin{definition}[Оператор альтернирования]
Для $T\in\Pi^0_p(V)$ оператор альтернирования задаётся формулой
\[
    \Alt(T)=\frac{1}{p!}\sum_{\sigma\in S_p}\sgn(\sigma)f_\sigma(T).
\]
\end{definition}

\begin{example}
Для $p=2$:
\[
    \Alt(e^1\otimes e^2)
    =
    \frac{1}{2}(e^1\otimes e^2-e^2\otimes e^1).
\]
В конспекте рядом записано сравнение с симметризацией:
\[
    \Sym(e^1\otimes e^2)
    =
    \frac{1}{2}(e^1\otimes e^2+e^2\otimes e^1).
\]
\end{example}

\begin{lemma}
Для оператора $\Alt$ выполнены:
\begin{enumerate}[label=\arabic*)]
    \item $\Imop(\Alt)=\Lambda^p(V^*)$;
    \item $\Ker(\Alt)$ состоит из тензоров, альтернирование которых равно нулю;
    \item $\Alt^2=\Alt$.
\end{enumerate}
\end{lemma}

\begin{proof}[Идея доказательства]
Если $T$ кососимметричен, то
\[
    f_\sigma(T)=\sgn(\sigma)T,
\]
поэтому
\[
    \Alt(T)=\frac{1}{p!}\sum_{\sigma\in S_p}\sgn(\sigma)^2T=T.
\]
С другой стороны, альтернирование произвольного тензора даёт кососимметрический тензор.
\begin{knowledgefix}
Базис симметрических тензоров получают симметризацией \(\Sym(e^{i_1}\otimes\cdots\otimes e^{i_p})\) с неубывающими индексами \(i_1\le\cdots\le i_p\). Эти элементы порождают все симметрические тензоры и линейно независимы.
\end{knowledgefix}
\end{proof}

\begin{corollary}
Если среди индексов $i_1,\ldots,i_p$ есть совпадающие, то
\[
    \Alt(e^{i_1}\otimes\cdots\otimes e^{i_p})=0.
\]
\end{corollary}

\begin{definition}[Внешнее произведение базисных ковекторов]
Для попарно различных индексов вводится обозначение
\[
    e^{i_1}\wedge\cdots\wedge e^{i_p}
    =
    \Alt(e^{i_1}\otimes\cdots\otimes e^{i_p})
\]
или, с возможным нормировочным множителем,
\[
    e^{i_1}\wedge\cdots\wedge e^{i_p}
    =
    \sum_{\sigma\in S_p}\sgn(\sigma)
    e^{i_{\sigma(1)}}\otimes\cdots\otimes e^{i_{\sigma(p)}}.
\]
\begin{knowledgefix}
Фиксируем нормировку: \(\Alt T=\frac1{p!}\sum_{\sigma\in S_p}\sgn(\sigma)F_\sigma T\), \(\Sym T=\frac1{p!}\sum_{\sigma\in S_p}F_\sigma T\). Если у лектора внешнее произведение нормируется иначе, все формулы отличаются только общим фактором, но свойства сохраняются.
\end{knowledgefix}
\end{definition}

\begin{corollary}[Базис и размерность внешних форм]
Элементы
\[
    e^{i_1}\wedge\cdots\wedge e^{i_p},
    \qquad
    1\le i_1<\cdots<i_p\le n,
\]
образуют базис пространства $\Lambda^p(V^*)$.
Следовательно,
\[
    \dim \Lambda^p(V^*)=\binom{n}{p}.
\]
Если $p>n$, то
\[
    \Lambda^p(V^*)=0.
\]
\begin{knowledgefix}
Для \(p=2\) любой ковариантный тензор раскладывается на симметрическую и кососимметрическую части. Для \(p\ge3\) одних полностью симметрической и полностью кососимметрической частей недостаточно: существуют смешанные типы симметрии.
\end{knowledgefix}
\end{corollary}

\subsubsection{Внешнее произведение и внешняя алгебра}

\subsubsection*{Внешнее произведение}

\begin{definition}[Внешнее произведение]
Пусть $T\in\Lambda^p(V^*)$, $S\in\Lambda^r(V^*)$. Их внешним произведением называется
\[
    T\wedge S=\Alt(T\otimes S),
\]
или, с нормировкой,
\[
    T\wedge S=\frac{(p+r)!}{p!r!}\Alt(T\otimes S).
\]
\begin{knowledgefix}
Для внешнего или симметрического произведения множители с \(p!\) и \(r!\) возникают из числа перестановок внутри блоков аргументов. При нормировке через альтернирование/симметризацию стандартная формула содержит коэффициент \(\frac{(p+r)!}{p!r!}\) перед проектором на степень \(p+r\).
\end{knowledgefix}
\end{definition}

\begin{remark}
На странице 96 записано, что симметрическая алгебра $S(V^*)$ может быть бесконечномерной, а для внешней алгебры ситуация другая: из-за кососимметричности произведения с повторяющимися множителями обращаются в ноль.
\end{remark}

\begin{proposition}[Свойства внешнего произведения]
Для внешних форм выполнены:
\begin{enumerate}[label=\arabic*)]
    \item билинейность;
    \item ассоциативность:
    \[
        (T\wedge S)\wedge R=T\wedge(S\wedge R);
    \]
    \item градуированная коммутативность:
    \[
        S\wedge T=(-1)^{pr}T\wedge S,
    \]
    где $T\in\Lambda^p(V^*)$, $S\in\Lambda^r(V^*)$.
\end{enumerate}
\end{proposition}

\begin{proof}[По структуре страниц 97--99]
Билинейность следует из линейности $\Alt$ и билинейности $\otimes$.
Ассоциативность доказывается тем, что обе части равны альтернированию полного тензорного произведения $T\otimes S\otimes R$ с одинаковой нормировкой.
Для коммутативности блоков нужно перенести блок из $r$ аргументов через блок из $p$ аргументов; это даёт знак
\[
    (-1)^{pr}.
\]
\begin{knowledgefix}
Ассоциативность внешнего произведения доказывается разбиением перестановок на перестановки внутри блоков и shuffle-перестановки между блоками. Это даёт стандартную формулу \((\alpha\wedge\beta)\wedge\gamma=\alpha\wedge(\beta\wedge\gamma)\).
\end{knowledgefix}
\end{proof}

\begin{corollary}
Если $S\in\Lambda^p(V^*)$ и $p$ нечётно, то
\[
    S\wedge S=0,
\]
поскольку
\[
    S\wedge S=(-1)^{p^2}S\wedge S=-S\wedge S.
\]
\begin{knowledgefix}
Градуированная коммутативность: если \(\alpha\in\Lambda^p V^*\), \(\beta\in\Lambda^q V^*\), то \(\alpha\wedge\beta=(-1)^{pq}\beta\wedge\alpha\). Поэтому \(\alpha\wedge\alpha=0\) при нечётном \(p\) над полем характеристики \(\ne2\).
\end{knowledgefix}
\end{corollary}

\subsubsection{Внешняя алгебра}

\subsubsection*{Определение внешней алгебры}

\begin{definition}[Внешняя алгебра]
Внешней алгеброй пространства $V^*$ называется прямая сумма
\[
    \Lambda(V^*)=\bigoplus_{p=0}^{n}\Lambda^p(V^*),
\]
где
\[
    \Lambda^0(V^*)=\F,
    \qquad
    \Lambda^1(V^*)=V^*.
\]
Умножение в этой алгебре задаётся внешним произведением $\wedge$.
\end{definition}

Так как $\Lambda^p(V^*)=0$ при $p>n$, то внешняя алгебра конечномерна:
\[
    \dim\Lambda(V^*)=
    \sum_{p=0}^{n}\binom{n}{p}=2^n.
\]

\begin{remark}
В конспекте подчёркнуто: внешняя алгебра является ассоциативной, но не коммутативной в обычном смысле. Вместо этого выполнена градуированная коммутативность.
\end{remark}

\begin{example}[Произведение базисных ковекторов]
Для базисных ковекторов выполнено
\[
    e^i\wedge e^j=-e^j\wedge e^i,
    \qquad
    e^i\wedge e^i=0.
\]
Например, произведение может быть ненулевым:
\[
    e^1\wedge e^2\wedge e^3\ne 0,
\]
если $e^1,e^2,e^3$ различны и $\dim V\ge 3$.
\begin{knowledgefix}
Если во внешнем произведении повторяется один и тот же ковектор, результат равен нулю: \(e^i\wedge e^i=0\). Это следует из кососимметричности при перестановке двух одинаковых аргументов.
\end{knowledgefix}
\end{example}

\begin{proposition}[Определитель как внешняя форма]
Пусть $e_1,\ldots,e_n$ --- базис $V$, а $e^1,\ldots,e^n$ --- двойственный базис. Тогда
\[
    e^1\wedge\cdots\wedge e^n
\]
является $n$-линейной кососимметрической формой. Для векторов
\[
    x_j=\sum_{i=1}^{n}x^i_j e_i
\]
выполнено
\[
    (e^1\wedge\cdots\wedge e^n)(x_1,\ldots,x_n)
    =
    \det(x^i_j).
\]
\begin{knowledgefix}
Базис \(\Lambda^p V^*\): элементы \(e^{i_1}\wedge\cdots\wedge e^{i_p}\), где \(1\le i_1<\cdots<i_p\le n\). Размерность равна \(\binom np\). Нормировочный множитель зависит от соглашения для \(\wedge\), но базис и размерность не меняются.
\end{knowledgefix}
\end{proposition}

\begin{corollary}
Если $\tilde e_1,\ldots,\tilde e_n$ --- другой базис $V$, то соответствующая верхняя внешняя форма меняется с множителем, равным определителю матрицы перехода.
\begin{knowledgefix}
Если \(\tilde e_i=C^j_i e_j\), то верхняя внешняя форма меняется умножением на \(\det C\) или \((\det C)^{-1}\) в зависимости от того, переход записан для базиса или для координат. Главное инвариантное утверждение: \(\Lambda^n V^*\) одномерно, и любой базисный объёмный элемент меняется на ненулевой скалярный множитель.
\end{knowledgefix}
\end{corollary}

\begin{remark}[Грассманова алгебра]
В конце страницы 101 внешняя алгебра описывается как градуированная антикоммутативная алгебра с образующими $x_1,\ldots,x_n$ и соотношениями
\[
    x_i x_j=-x_j x_i\quad (i\ne j),
    \qquad
    x_i^2=0.
\]
Это соответствует правилам внешнего произведения базисных ковекторов.
\begin{knowledgefix}
Грассманова, или внешняя, алгебра: \(\Lambda(V^*)=\bigoplus_{p=0}^n\Lambda^p(V^*)\) с умножением \(\wedge\). Её базис образуют внешние мономы \(e^{i_1}\wedge\cdots\wedge e^{i_p}\), \(i_1<\cdots<i_p\).
\end{knowledgefix}
\end{remark}
\begin{planadd}
\textbf{Стандартное оформление последней части билета.}
Внешняя алгебра
\[
    \Lambda(V^*)=\bigoplus_{p=0}^{n}\Lambda^p(V^*)
\]
имеет базис
\[
    e^{i_1}\wedge\cdots\wedge e^{i_p}, \qquad 1\le i_1<\cdots<i_p\le n.
\]
Поэтому
\[
    \dim \Lambda^p(V^*)=\binom np.
\]
Изоморфизм с алгеброй грассмановых многочленов задаётся соответствием
$e^i\leftrightarrow \theta_i$, где переменные антикоммутируют:
$\theta_i\theta_j=-\theta_j\theta_i$ и $\theta_i^2=0$.
\end{planadd}

\clearpage

\end{document}
